\documentclass[11pt]{amsart}
\usepackage[margin=1.1in]{geometry}
\usepackage{amsmath,amssymb,amsthm,mathtools,bm}
\usepackage{enumitem}
\usepackage[colorlinks=true,linkcolor=blue,urlcolor=blue]{hyperref}

\numberwithin{equation}{section}
\theoremstyle{plain}
\newtheorem{theorem}{Theorem}[section]
\newtheorem{proposition}[theorem]{Proposition}
\newtheorem{lemma}[theorem]{Lemma}
\newtheorem{corollary}[theorem]{Corollary}
% \theoremstyle{definition} 
\newtheorem{definition}[theorem]{Definition}
\newtheorem{example}[theorem]{Example}
\newtheorem{assumption}[theorem]{Assumption}
\newtheorem{exercise}[theorem]{Exercise}
% \theoremstyle{remark}
\newtheorem{remark}[theorem]{Remark}

\newcommand{\R}{\mathbb{R}}
\newcommand{\E}{\mathbb{E}}
\newcommand{\Prob}{\mathbb{P}}
\newcommand{\cP}{\mathcal{P}}
\newcommand{\cK}{\mathcal{K}}
\newcommand{\cL}{\mathcal{L}}
\newcommand{\dd}{\,\mathrm{d}}
\newcommand{\abs}[1]{\left\lvert#1\right\rvert}
\newcommand{\norm}[1]{\left\lVert#1\right\rVert}
\newcommand{\ip}[2]{\left\langle#1,#2\right\rangle}
\newcommand{\one}{\mathbf{1}}
\DeclareMathOperator{\Law}{Law}
\DeclareMathOperator{\supp}{supp}
\DeclareMathOperator{\diam}{diam}
\DeclareMathOperator{\diver}{div}
\DeclareMathOperator{\Id}{Id}
\DeclareMathOperator{\Lip}{Lip}
\DeclareMathOperator*{\argmin}{arg\,min}

\newcommand{\blue}[1]{{\color{blue}#1}}
\newcommand{\cyan}[1]{{\color{cyan}#1}}
\newcommand{\red}[1]{{\color{red}#1}}


\title{Interacting Particle Systems and Transformers}
\author{Fei Lu}
\date{Fall 2026. Disclaimer: the notes are written with the help from chatGPT, and the author takes full responsibility for any errors. If you find any errors, please contact the author at \href{mailto:feilu@math.jhu.edu}{feilu@math.jhu.edu}.}

\begin{document}
\begin{abstract}
These notes develop the interacting-particle viewpoint on self-attention.
The mathematical spine is the passage from finite particle updates to
measure-valued dynamics, followed by continuous depth, collective behavior,
stability, and the inverse problem of learning an interaction law.  The notes
distinguish exact finite-token identities from mean-field and continuous-depth
approximations, and they state the assumptions under which those approximations
are justified.  A final discussion connects the common object---a learned
nonlocal operator on empirical measures---to generative modeling, in-context
learning, operator learning, optimal transport, inverse problems, and
architecture design.
\end{abstract}

\maketitle
\tableofcontents

\section*{Course thread, prerequisites, and notation}
For particles or tokens $X^1,\ldots,X^N\in\R^d$, define the empirical measure
\begin{equation}\label{eq:empirical-measure}
    \mu^N=\frac1N\sum_{i=1}^N\delta_{X^i}.
\end{equation}
It forgets labels but retains the distribution of states, making it the natural
state variable for permutation-equivariant dynamics.  This loss of labels is a
feature for exchangeable particle clouds and a limitation for ordered
sequences.  Positions, particle species, boundary labels, and causal
information must be retained as \emph{marks}; see
Section~\ref{sec:self-attention}.

We write $\cP_p(\R^d)$ for probability measures with finite $p$th moment.  If
$T:\R^d\to\R^m$ is measurable, its pushforward is
\[
    T_\#\mu(A)=\mu(T^{-1}(A)),\qquad
    \int f\,\dd(T_\#\mu)=\int f(T(x))\,\mu(\dd x).
\]
For $\mu,\nu\in\cP_1(\R^d)$, the $1$-Wasserstein distance is
\begin{equation}\label{eq:w1-definition}
    W_1(\mu,\nu)
    =\inf_{\pi\in\Gamma(\mu,\nu)}
      \int_{\R^d\times\R^d}\abs{x-y}\,\pi(\dd x,\dd y),
\end{equation}
where $\Gamma(\mu,\nu)$ is the set of couplings.  We use the same notation for a
measure and its density when no confusion can arise.  All gradients and
divergences act on the displayed spatial variable.

The sections follow one chain:
\[
\begin{split}
  \text{finite interacting particles}
  &\longrightarrow \text{Wasserstein transport and gradient flow}\\
  &\longrightarrow \text{attention as a nonlocal particle update}\\
  &\longrightarrow \text{continuous depth, clustering, and stability}\\
  &\longrightarrow \text{learning the interaction from data}.
\end{split}
\]
The prerequisites are It\^o's formula, basic ODE/SDE well-posedness, weak
convergence of measures, and Gronwall's inequality.  The notes recall the
specific facts needed from each.  Statements are intentionally finite-horizon
unless a dissipative hypothesis gives genuine long-time control.

\begin{remark}[Three levels that must not be conflated]\label{rem:three-levels}
The identity expressing a self-attention layer through its empirical measure is
exact.  The limit $N\to\infty$ is a probabilistic approximation requiring
exchangeability and regularity.  The limit of many small residual layers is a
numerical-analysis approximation requiring a small step size and controlled
variation of the layer parameters.  Neither asymptotic limit is automatic for
a trained language model.
\end{remark}


%%%%%%
\section{Mean-field interacting particle systems}\label{sec:ips}
\paragraph{\textbf{Key question}}
What equation describes the limit of the empirical measure of $N$
interacting particles as $N\to\infty$, and why does the limiting equation
close?

\subsection{The finite-particle system}
Consider $N$ particles in $\R^d$ with pairwise interaction
$b:\R^d\times\R^d\to\R^d$ and additive Brownian noise:
\begin{equation}\label{eq:particles}
    \dd X_t^i=\frac1N\sum_{j=1}^N b(X_t^i,X_t^j)\dd t
       +\sigma\dd B_t^i,
    \qquad i=1,\ldots,N.
\end{equation}
Here $B^1,\ldots,B^N$ are independent standard Brownian motions in $\R^d$,
$\sigma\geq0$, and the initial variables $X_0^1,\ldots,X_0^N$ are i.i.d.\ with
law $\mu_0\in\cP_2(\R^d)$, independently of the Brownian motions.  The case
$\sigma=0$ is deterministic and $\sigma>0$ is the stochastic system in the Ito sense.

The factor $1/N$ is the mean-field scaling: each pair interaction is weak, but
each particle feels $N$ interactions.  We include $j=i$ for convenience.  If
the self-interaction is omitted, the drift changes by
$N^{-1}b(X_t^i,X_t^i)$, which is lower order under the moment estimates below.

\begin{assumption}[Lipschitz interaction]\label{ass:regular-mf}
There is a constant $L_b\geq0$ such that
\[
  \abs{b(x,y)-b(x',y')}
  \leq L_b\bigl(\abs{x-x'}+\abs{y-y'}\bigr),
  \qquad x,x',y,y'\in\R^d.
\]
\end{assumption}
In particular, if $b$ is globally Lipschitz, then it has linear growth:
\begin{equation}\label{eq:b-linear-growth}
  \abs{b(x,y)}\leq \abs{b(0,0)}+L_b(\abs{x}+\abs{y}).
\end{equation}
One may replace the global Lipschitz assumption with a local Lipschitz assumption plus a linear growth condition, but the global Lipschitz assumption is stronger and easier to use.

\begin{proposition}[Finite-system well-posedness and moments]
\label{prop:finite-wp}
Under Assumption~\ref{ass:regular-mf}, \eqref{eq:particles} has a unique global
strong solution.  For every $T<\infty$,
\begin{equation}\label{eq:finite-moment}
  \sup_{N\geq1}\max_{1\leq i\leq N}
  \E\left[\sup_{0\leq t\leq T}\abs{X_t^i}^2\right]
  \leq C_T\left(1+\int_{\R^d}\abs{x}^2\mu_0(\dd x)\right),
\end{equation}
where $C_T$ may depend on $T,L_b,\abs{b(0,0)},\sigma$, and $d$, but not on
$N$.
\end{proposition}
\begin{proof}
For $\boldsymbol{x}=(x_1,\ldots,x_N)$, set
\[
 F_i^N(\boldsymbol{x})=\frac1N\sum_{j=1}^N b(x_i,x_j),
 \qquad
 \norm{\boldsymbol{x}}_N^2=\frac1N\sum_{i=1}^N\abs{x_i}^2.
\]
The Lipschitz assumption and Jensen's inequality give
\[
 \norm{F^N(\boldsymbol{x})-F^N(\boldsymbol{x}')}{}_N^2
 \leq 4L_b^2\norm{\boldsymbol{x}-\boldsymbol{x}'}_N^2.
\]
Thus the drift of the $Nd$-dimensional SDE is globally Lipschitz, so the
standard SDE theorem gives a unique global strong solution; see
\cite[Chapter~5]{KaratzasShreve1991}.  Moreover,
\eqref{eq:b-linear-growth} implies
\[
 \abs{F_i^N(\boldsymbol{x})}^2
 \leq C\left(1+\abs{x_i}^2+\frac1N\sum_{j=1}^N\abs{x_j}^2\right).
\]
Applying the elementary inequality
$\sup_{r\leq t}\abs{\int_0^r f_s\dd s}^2
\leq t\int_0^t\abs{f_s}^2\dd s$ to the integral form of
\eqref{eq:particles} gives, for $t\leq T$,
\begin{align*}
 \E\sup_{0\leq r\leq t}\abs{X_r^i}^2
 &\leq 3\E\abs{X_0^i}^2
 +3t\int_0^t\E\abs{F_i^N(\boldsymbol X_s)}^2\dd s
 +3\sigma^2\E\sup_{0\leq r\leq t}\abs{B_r^i}^2\\
 &\leq 3\E\abs{X_0^i}^2
 +3T\int_0^t\E\abs{F_i^N(\boldsymbol X_s)}^2\dd s
 +12d\sigma^2T, 
\end{align*}
where in the last step, Doob's $L^2$ maximal inequality (see Theorem \ref{thm:appendix-doob})
\footnote{TBV: Doob's $L^p$ maximal inequality states that for a continuous martingale $M$ with $M_0=0$ and $p>1$, $\E\sup_{0\leq t\leq T}\abs{M_t}^p\leq (p/(p-1))^p\E\abs{M_T}^p$.  The constant is sharp.  See \cite[Theorem~3.2]{KaratzasShreve1991}. Alternatively,  this is also the Brownian special case of the Burkholder--Davis--Gundy
inequality.  No stochastic-integral estimate is needed here because the
integral equation contains the additive term $\sigma B_t^i$ directly.  If one instead applies It\^o's formula to $\abs{X_t^i}^2$, then BDG is needed to control the martingale $2\sigma\int_0^t X_s^i\cdot\dd B_s^i$. Also, when the noise is multiplicative, then BDG can control the martingale term in the It\^o integral. }
gives
\[
 \E\sup_{0\leq r\leq t}\abs{B_r^i}^2
 \leq 4\E\abs{B_t^i}^2=4dt.
\] 

Set
\[
 u_i(t):=\E\sup_{0\leq r\leq t}\abs{X_r^i}^2,
 \qquad
 A_N(t):=\max_{1\leq i\leq N}u_i(t).
\]
For each $i$, the preceding estimate contains the integral of
$u_i(s)+N^{-1}\sum_j u_j(s)$, which is at most $2A_N(s)$.  Since the
initial variables have common law $\mu_0$, taking the maximum over $i$ yields
\[
 A_N(t)
 \leq C_T\left(1+\int\abs{x}^2\mu_0(\dd x)
 +\int_0^t A_N(s)\dd s\right).
\]
Gronwall's inequality now gives \eqref{eq:finite-moment}.  Notice that the maximum is outside the
expectation: the stronger quantity
$\E[\max_i\sup_{r\leq t}\abs{X_r^i}^2]$ generally grows with $N$ and would
not give the claimed uniform estimate.
\end{proof}

Singular forces such as Coulomb kernels are not covered by
Assumption~\ref{ass:regular-mf}; their well-posedness and mean-field limits
require estimates adapted to the singularity.

\begin{remark}[Singular interactions remain an active area] TBV. 
The Lipschitz coupling used in this section breaks down for Coulomb, Riesz,
Biot--Savart, and strongly attractive kernels, and there is no single
replacement covering all of them.  Several complementary methods have
developed rapidly.  Jabin and Wang proved quantitative propagation of chaos
in relative entropy for stochastic systems with
$\dot W^{-1,\infty}$ interaction kernels, a class that includes the
Biot--Savart kernel for the two-dimensional viscous vortex model
\cite{JabinWang2018}.  Bresch, Jabin, and Wang combined entropy and modulated
free energy to treat singular attractive interactions and obtained a
quantitative particle derivation of the Patlak--Keller--Segel equation in the
optimal subcritical regime \cite{BreschJabinWang2023}.  More recent results
include strong trajectory-level propagation of chaos for singular
$L^p$ interactions \cite{HaoRocknerZhang2024}, sharp uniform-in-time
mean-field convergence for periodic Riesz flows
\cite{ChodronRosenzweigSerfaty2025}, and sharp local chaos estimates for
$\dot W^{-1,\infty}$ interactions \cite{Wang2026}.  These theorems rely on
diffusion strength, sign, dimension, integrability, or modulated-energy
structure; none makes a generic singular force as routine as the globally
Lipschitz model above.
\end{remark}


\subsection{The weak equation for the empirical measure}

The empirical measure of the particle system is
\begin{equation}\label{eq:empirical-measure-time}
    \mu_t^N:=\frac1N\sum_{i=1}^N\delta_{X_t^i}.
\end{equation}
For a measure $\mu$, use the compact notation
\begin{equation}\label{eq:b-mu}
    b[\mu](x):=\int_{\R^d}b(x,y)\,\mu(\dd y).
\end{equation}
Then the drift of particle $i$ is exactly
$b[\mu_t^N](X_t^i)$.  This identity is the first closure step: the other
particles enter the equation for $X_t^i$ only through $\mu_t^N$.


For $\varphi\in C_b^2(\R^d)$, It\^o's formula for the $i$th particle gives
\[
 \dd\varphi(X_t^i)
 =\nabla\varphi(X_t^i)\cdot b[\mu_t^N](X_t^i)\dd t
 +\frac{\sigma^2}{2}\Delta\varphi(X_t^i)\dd t
 +\sigma\nabla\varphi(X_t^i)\cdot\dd B_t^i.
\]
Averaging over $i$ yields the semimartingale identity
\begin{align}
    \dd\int\varphi\,\dd\mu_t^N
    &=\int\left[b[\mu_t^N]\cdot\nabla\varphi
       +\frac{\sigma^2}{2}\Delta\varphi\right]\dd\mu_t^N\dd t
       +\dd M_t^{N,\varphi},\label{eq:empirical-weak}\\
    M_t^{N,\varphi}
    &=\frac{\sigma}{N}\sum_{i=1}^N
      \int_0^t\nabla\varphi(X_s^i)\cdot\dd B_s^i.
      \label{eq:empirical-martingale}
\end{align}
For a continuous square-integrable martingale $M$ with $M_0=0$, we write
$\langle M\rangle_t$ for its predictable quadratic variation: it is the unique increasing predictable process such that
$M_t^2-\langle M\rangle_t$ is a martingale; see \cite[Chapter~3]{KaratzasShreve1991}.  
In particular,
\[
 \E\abs{M_t}^2=\E\langle M\rangle_t.
\]
The independence of the Brownian motions and the stochastic-integral
quadratic-variation formula yield
\begin{equation}\label{eq:martingale-bracket}
  \left\langle M^{N,\varphi}\right\rangle_t
  =\frac{\sigma^2}{N^2}\sum_{i=1}^N
    \int_0^t\abs{\nabla\varphi(X_s^i)}^2\dd s
  \leq\frac{\sigma^2t}{N}\norm{\nabla\varphi}_\infty^2.
\end{equation}
Thus the stochastic fluctuation of any fixed smooth observable is of size
$N^{-1/2}$.  Notice that \eqref{eq:martingale-bracket} depends on the test function, so it is not a metric between the entire empirical measure and its limit.

For simplicity, we can write the weak-form equation as
\[
 \partial_t \mu_t^N = -\diver\bigl(\mu_t^N b[\mu_t^N]\bigr)  + \frac{\sigma^2}{2} \Delta \mu_t^N + \text{martingale noise}.
\]

%%%%%%%%%%%%%% ============ 
\subsection{The nonlinear limit and its well-posedness}
The candidate limiting process is the McKean--Vlasov SDE with additive noise:
\begin{equation}\label{eq:mv-sde}
  \dd\overline X_t=b[\mu_t](\overline X_t)\dd t+\sigma\dd B_t,
  \qquad \mu_t=\Law(\overline X_t),\qquad
  \Law(\overline X_0)=\mu_0, 
\end{equation}
where $ b[\mu](x)=\int_{\R^d}b(x,y)\,\mu(\dd y)$. 
It is nonlinear in law: the unknown law $\mu_t$ enters the drift that generates that same law. 

Next, we prove well-posedness of the McKean--Vlasov SDE.  The proof is based on a fixed-point construction and its propagation-of-chaos consequence are classical; see \cite{Sznitman1991} and \cite{CarmonaDelarue2018}. We start from a bound on the drift in terms of the Wasserstein distance between measures. 
For $\mu,\nu\in\cP_2(\R^d)$, let
\[
 W_2(\mu,\nu)^2
 :=\inf_{\pi\in\Gamma(\mu,\nu)}
   \int_{\R^d\times\R^d}\abs{x-y}^2\pi(\dd x,\dd y), 
\]
where $\Gamma(\mu,\nu):= \{ \pi \in \mathcal{P}(\R^d \times \R^d) : \pi(\cdot, \R^d) = \mu, \pi(\R^d, \cdot) = \nu \}$ is the set of all couplings of $\mu$ and $\nu$. 


\begin{lemma}[Lipschitz continuity of the mean-field drift]
\label{lem:measure-drift-lipschitz}
Under Assumption~\ref{ass:regular-mf}, for all
$x,x'\in\R^d$ and $\mu,\nu\in\cP_2(\R^d)$,
\begin{equation}\label{eq:measure-drift-lipschitz}
 \abs{b[\mu](x)-b[\nu](x')}
 \leq L_b\bigl(\abs{x-x'}+W_2(\mu,\nu)\bigr).
\end{equation}
\end{lemma}
\begin{proof}
Let $\pi\in\Gamma(\mu,\nu)$ be any coupling.  Since the first and second
marginals of $\pi$ are $\mu$ and $\nu$, respectively,
\[
 b[\mu](x)-b[\nu](x')
 =\int_{\R^d\times\R^d}
   \bigl(b(x,y)-b(x',y')\bigr)\,\pi(\dd y,\dd y').
\]
The triangle inequality and Assumption~\ref{ass:regular-mf} give
\begin{align*}
 \abs{b[\mu](x)-b[\nu](x')}
 &\leq L_b\int
   \bigl(\abs{x-x'}+\abs{y-y'}\bigr)\,\pi(\dd y,\dd y')\\
 &=L_b\abs{x-x'}
   +L_b\int\abs{y-y'}\,\pi(\dd y,\dd y')\\
 &\leq L_b\abs{x-x'}
   +L_b\left(\int\abs{y-y'}^2\,\pi(\dd y,\dd y')\right)^{1/2},
\end{align*}
where the last step is Cauchy--Schwarz and $\pi$ has total mass one.
Taking the infimum over $\pi\in\Gamma(\mu,\nu)$ proves
\eqref{eq:measure-drift-lipschitz}.
\end{proof}

\begin{theorem}[Well-posedness of the McKean--Vlasov SDE]
\label{thm:mv-wellposed}
Under Assumption~\ref{ass:regular-mf}, \eqref{eq:mv-sde} has a unique global
strong solution and $\mu\in C([0,\infty);\cP_2(\R^d))$ with the Wasserstein-2 metric.  For every $T<\infty$,
\begin{equation}\label{eq:mv-moment}
 \E\left[\sup_{0\leq t\leq T}\abs{\overline X_t}^2\right]
 \leq C_T\left(1+\int_{\R^d}\abs{x}^2\mu_0(\dd x)\right).
\end{equation}
\end{theorem}
\begin{proof}   The proof is standard, but we include it for completeness and to highlight the role of the Wasserstein distance in the contraction argument. 

\textbf{Part I. Strong well-posedness.}  We first prove the strong well-posedness by using a fixed-point argument on the space of continuous curves of measures. 
Fix $h\in (0,T]$ and let
\[
 \mathcal C_h=\left\{m\in C([0,h];\cP_2(\R^d)):m_0=\mu_0\right\},
 \qquad d_h(m,n)=\sup_{0\leq t\leq h}W_2(m_t,n_t).
\]
One can verify that $\mathcal C_h$ with $d_h$ is a complete metric space.  For a deterministic curve $m\in\mathcal
C_h$, the SDE
\begin{equation}\label{eq:frozen-law-sde}
 \dd X_t^m=b[m_t](X_t^m)\dd t+\sigma\dd B_t,
 \qquad X_0^m=X_0,\quad\Law(X_0)=\mu_0,
\end{equation}
has a unique strong solution because its drift is globally Lipschitz in $x$
and has linear growth.  Define
$$\Phi(m)_t=\Law(X_t^m).$$  

We will show that $\Phi$ is a contraction on $\mathcal C_h$ for sufficiently small $h$ in two steps. 

\paragraph{\textbf{Step 1. $\Phi$ maps $\mathcal C_h$ into itself.}}  We need to show that $\Phi(m)\in\mathcal C_h$, i.e., that $\Phi(m)_0=\mu_0$ and that $t\mapsto \Phi(m)_t$ is continuous in the Wasserstein-2 metric, which we obtain by showing mean-square continuity of $X^m$. 

Since $m\in C([0,h];\cP_2(\R^d))$, we have
\[
 K_m:=\sup_{0\leq t\leq h}
 \int_{\R^d}\abs{y}^2m_t(\dd y)<\infty.
\]
Indeed, 
$\bigl(\int\abs{y}^2m_t(\dd y)\bigr)^{1/2}
=W_2(m_t,\delta_0)$ is continuous in $t$.  By
\eqref{eq:b-linear-growth} and Cauchy--Schwarz,
\begin{equation}\label{eq:frozen-drift-growth}
 \abs{b[m_t](x)}^2 % = \abs{\int b(x,y)m_t(\dd y)}^2 
 \leq 2\abs{b(0,0)}^2+2L_b^2\left(\abs{x}+\int\abs{y}m_t(\dd y)\right)^2
 \leq C\left(1+\abs{x}^2+K_m\right).
\end{equation}
The integral form of \eqref{eq:frozen-law-sde}, the same maximal estimate used
in Proposition~\ref{prop:finite-wp}, and
\eqref{eq:frozen-drift-growth} give
\[
 Q_m(t):=\E\sup_{0\leq r\leq t}\abs{X_r^m}^2
 \leq C_{h,m}\left(1+\E\abs{X_0}^2\right)
 +C_h\int_0^tQ_m(s)\dd s.
\]
Gronwall's inequality therefore gives
\begin{equation}\label{eq:frozen-moment}
 \E\sup_{0\leq t\leq h}\abs{X_t^m}^2
 \leq C_{h,m}\left(1+\E\abs{X_0}^2\right)<\infty.
\end{equation}
For $0\leq s\leq t\leq h$, the increment equation gives
\begin{align}
 \E\abs{X_t^m-X_s^m}^2
 &\leq 2(t-s)\int_s^t\E\abs{b[m_r](X_r^m)}^2\dd r
      +2d\sigma^2(t-s)\notag\\
 &\leq C_{h,m}(t-s),\label{eq:frozen-ms-continuity}
\end{align}
where \eqref{eq:frozen-drift-growth} and \eqref{eq:frozen-moment} are used in
the last step.  Thus $X^m$ is mean-square continuous, and the coupling
$(X_t^m,X_s^m)$ shows
\[
 W_2(\Phi(m)_t,\Phi(m)_s)^2
 \leq\E\abs{X_t^m-X_s^m}^2\longrightarrow0
 \qquad\text{as }t\to s.
\]
Together with $\Phi(m)_0=\mu_0$ and \eqref{eq:frozen-moment}, this proves
$\Phi(m)\in\mathcal C_h$.

\paragraph{\textbf{Step 2. $\Phi:\mathcal C_h\to\mathcal C_h$ is a contraction.}}
To prove contraction, let $X^m$ and $X^n$ solve \eqref{eq:frozen-law-sde} for $m$ and $n$ with the same $X_0$ and Brownian motion.  By \eqref{eq:measure-drift-lipschitz}, with
\[
 D(t):=\E\left[\sup_{0\leq r\leq t}\abs{X_r^m-X_r^n}^2\right],
\]
Cauchy--Schwarz gives, for $t\leq h$,
\[
 D(t)
 \leq 2L_b^2h\int_0^t
       \left(D(s)+d_h(m,n)^2\right)\dd s.
\]
Indeed, if $\Delta_t=X_t^m-X_t^n$, then the common Brownian motion and common
initial variable cancel, so
\[
 \Delta_r=\int_0^r
 \bigl(b[m_s](X_s^m)-b[n_s](X_s^n)\bigr)\dd s.
\]
For every $r\leq t\leq h$, Cauchy--Schwarz in time and
Lemma~\ref{lem:measure-drift-lipschitz} give
\begin{align*}
 \abs{\Delta_r}^2
 &\leq r\int_0^r
 \abs{b[m_s](X_s^m)-b[n_s](X_s^n)}^2\dd s\\
 &\leq 2L_b^2h\int_0^t
 \left(\abs{\Delta_s}^2+W_2(m_s,n_s)^2\right)\dd s\\
 &\leq 2L_b^2h\int_0^t
 \left(\sup_{0\leq u\leq s}\abs{\Delta_u}^2
       +d_h(m,n)^2\right)\dd s.
\end{align*}
Taking the supremum over $r\leq t$ and then expectation proves the stated
bound for $D(t)$.
Gronwall's inequality therefore gives
\begin{equation}\label{eq:phi-contraction}
 d_h(\Phi(m),\Phi(n))^2
 = \sup_{0\leq t\leq h}W_2(\Phi(m)_t,\Phi(n)_t)^2
 \leq \sup_{0\leq t\leq h}D(t)  
 \leq 2L_b^2h^2e^{2L_b^2h^2}d_h(m,n)^2.
\end{equation}
Choose $h$ so that the coefficient on the right is less than one. 

 Banach's fixed-point theorem gives a unique curve $\mu=\Phi(\mu)$ on $[0,h]$; the corresponding $X^\mu$ solves \eqref{eq:mv-sde}.  Repeating the argument on successive intervals of length $h$, with the terminal law from one interval as the initial law for the next, constructs the solution on every finite time interval. The resulting $\mu\in C([0,\infty);\cP_2(\R^d))$ is unique in the sense that any other solution must coincide with it on each interval of length $h$.  

If two solutions are driven by the same $(X_0,B)$, uniqueness of the fixed
point first gives equality of their law curves.  They then solve the same
ordinary Lipschitz SDE, so pathwise uniqueness gives equality of the
processes.

\textbf{Part II. The moment bound.}
To prove \eqref{eq:mv-moment}, note that the Lipschitz continuity of $b$ and the linear growth condition imply, 
\[
 Q(t): =\E[\sup_{0\leq r\leq t}\abs{\overline X_r}^2] \leq C_T\left(1+\E\abs{\overline X_0}^2\right)
       +C_T\int_0^tQ(s)\dd s,\qquad t\leq T.
\]
To derive it, first use the linear growth bound \eqref{eq:b-linear-growth}, Jensen's inequality, and
$\mu_s=\Law(\overline X_s)$:
\begin{align}
 \E\abs{b[\mu_s](\overline X_s)}^2
 &\leq C\E\left[
 1+\abs{\overline X_s}^2
 +\left(\int\abs{y}\,\mu_s(\dd y)\right)^2\right]\notag\\
 &\leq C\left(1+\E\abs{\overline X_s}^2
 +\int\abs{y}^2\mu_s(\dd y)\right)\notag\\
 &=C\left(1+2\E\abs{\overline X_s}^2\right)
 \leq C\left(1+Q(s)\right).
 \label{eq:mv-drift-moment}
\end{align}
Next, the integral form of \eqref{eq:mv-sde} $ 
\overline X_t=\overline X_0+\int_0^t b[\mu_s](\overline X_s)\dd s+\sigma B_t 
$ gives, for $t\leq T$, 
\begin{align*}
 Q(t)
 &\leq 3\E\abs{\overline X_0}^2
 +3t\int_0^t\E\abs{b[\mu_s](\overline X_s)}^2\dd s
 +3\sigma^2\E\sup_{0\leq r\leq t}\abs{B_r}^2\\
 &\leq 3\E\abs{\overline X_0}^2
 +3TC\int_0^t(1+Q(s))\dd s+12d\sigma^2T,
\end{align*}
where Doob's inequality was used in the last term.  Absorbing the constant
terms into $C_T$ yields the claimed inequality. Gronwall's inequality gives \eqref{eq:mv-moment}.
\end{proof}

Taking expectations in It\^o's formula for \eqref{eq:mv-sde}, we have that the law $\mu_t$ solves the weak Fokker--Planck equation, i.e., $(\mu_t)_{0\leq t\leq T}\subset\cP_2(\R^d)$ is continuous in time and satisfies, 
\begin{equation}\label{eq:mv-weak} 
\begin{aligned}
  \partial_t \mu_t & = -\diver\bigl(\mu_t b[\mu_t]\bigr)  + \frac{\sigma^2}{2} \Delta \mu_t, \quad   \Leftrightarrow  \\
   \int\varphi\,\dd\mu_t-\int\varphi\,\dd\mu_0
   & =\int_0^t\int\left[
       b[\mu_s]\cdot\nabla\varphi
       +\frac{\sigma^2}{2}\Delta\varphi\right]\dd\mu_s\dd s, \quad \forall \varphi\in C_c^\infty(\R^d). 
\end{aligned}
\end{equation} 
If $\mu_t$ has a smooth density $\rho_t$,
we have the strong form of the nonlinear Fokker--Planck equation
\begin{equation}\label{eq:mv-fp}
    \partial_t\rho_t
    +\diver\bigl(\rho_t b[\rho_t]\bigr)
    =\frac{\sigma^2}{2}\Delta\rho_t.
\end{equation}
The weak formulation remains meaningful when no smooth density exists and is
therefore the primary equation.

\subsection{Mean-field limit: propagation of chaos}
The disappearance of the martingale in \eqref{eq:empirical-weak} suggests a
deterministic limit, but closure also requires replacing
$b[\mu_t^N]$ by $b[\mu_t]$.  A synchronous coupling makes this precise.

\begin{theorem}[Finite-time propagation of chaos]\label{thm:poc} Let $N\geq 2$. 
Under Assumption~\ref{ass:regular-mf}, let
$\overline X^1,\ldots,\overline X^N$ be independent copies of
\eqref{eq:mv-sde}, with $\overline X_0^i=X_0^i$ and driven by the same $B^i$ as
$X^i$ in \eqref{eq:particles}.   For every $T<\infty$ there is $C_T$, independent of $N$, such that
\begin{equation}\label{eq:poc-rate}
  \max_{1\leq i\leq N}
  \E\left[\sup_{0\leq t\leq T}
    \abs{X_t^i-\overline X_t^i}^2\right]\leq\frac{C_T}{N}.
\end{equation}
Consequently, for every fixed $k$ and $t\leq T$,
\begin{equation}\label{eq:k-chaos-rate}
 W_2^2\left(\Law(X_t^1,\ldots,X_t^k),\mu_t^{\otimes k}\right)
 \leq \frac{kC_T}{N},
\end{equation}
where $W_2$ on the left is computed with the Euclidean norm on $\R^{dk}$.
\end{theorem}
\begin{proof}
Set $\Delta_t^i=X_t^i-\overline X_t^i$ and let
\[
 D(t)=\max_{1\leq i\leq N}
 \E\left[\sup_{0\leq r\leq t}\abs{\Delta_r^i}^2\right].
\]

The Brownian terms cancel in the
synchronous coupling.  Add and subtract
$N^{-1}\sum_j b(\overline X_t^i,\overline X_t^j)$ to write
\begin{equation}\label{eq:coupling-error-decomposition}
 \Delta_t^i=\int_0^t\left(A_s^i+R_s^i\right)\dd s,
\end{equation}
where
\begin{align*}
 A_s^i=\frac1N\sum_{j=1}^N
 \left[b(X_s^i,X_s^j)-b(\overline X_s^i,\overline X_s^j)\right],\quad 
 R_s^i=\frac1N\sum_{j=1}^N b(\overline X_s^i,\overline X_s^j)
       -b[\mu_s](\overline X_s^i).
\end{align*}
The Lipschitz property and Jensen's inequality imply
\begin{equation}\label{eq:A-bound}
 \E \abs{A_s^i}^2
 \leq 2L_b^2\E \left(\abs{\Delta_s^i}^2
       +\frac1N\sum_{j=1}^N\abs{\Delta_s^j}^2\right) \leq4L_b^2D(s).
\end{equation}



To estimate the empirical sampling error $R_s^i$, define
$g_s(x,y)=b(x,y)-b[\mu_s](x)$.  Conditional on
$\overline X_s^i=x$, the variables
$g_s(x,\overline X_s^j)$, $j\ne i$, are independent and centered.  Hence all
their cross terms vanish, while the self term is kept separate:
\begin{align}
 \E\abs{R_s^i}^2
 &=\frac{N-1}{N^2}
   \E\abs{g_s(\overline X_s^i,\overline X_s^j)}^2
   +\frac1{N^2}
   \E\abs{g_s(\overline X_s^i,\overline X_s^i)}^2
 \leq \frac{C_T}{N}, \label{eq:R-bound}
\end{align}
where $j\ne i$ in the first term and $C_T$ is independent of $N$. 
To justify the last inequality, observe that $ \abs{g_s(x,y)} = \abs{b(x,y)-b[\mu_s](x)} 
 \leq L_b\int\abs{y-z}\,\mu_s(\dd z)$ 
and use the uniform second-moment bound \eqref{eq:mv-moment}.  More precisely,
set $ m_2(s):=\int_{\R^d}\abs{z}^2\mu_s(\dd z)
       =\E\abs{\overline X_s}^2$. 
Jensen's inequality with respect to $\mu_s$ gives, for every $x,y\in\R^d$,
\begin{align*}
 \abs{g_s(x,y)}^2
 &\leq L_b^2\int_{\R^d}\abs{y-z}^2\mu_s(\dd z) \leq 2L_b^2\left(\abs{y}^2+m_2(s)\right).
\end{align*}
Since each $\overline X_s^j$ has law $\mu_s$, this estimate applies both to
the off-diagonal term $j\ne i$ and to the self term $j=i$:
\begin{align*}
 \E\abs{g_s(\overline X_s^i,\overline X_s^j)}^2  \leq4L_b^2m_2(s),\quad 
 \E\abs{g_s(\overline X_s^i,\overline X_s^i)}^2  \leq4L_b^2m_2(s).
\end{align*}
Finally, 
\[
 \sup_{0\leq s\leq T}m_2(s) = \sup_{0\leq s\leq T}\E\abs{\overline X_s}^2
 \leq C_T
\]
by \eqref{eq:mv-moment}, which proves the uniform bound in
\eqref{eq:R-bound}.



Lastly, for every $i$ and $t\leq T$, first take the supremum in
\eqref{eq:coupling-error-decomposition} and apply Cauchy--Schwarz in time:
\begin{align*}
 \sup_{0\leq r\leq t}\abs{\Delta_r^i}^2
 &\leq t\int_0^t\abs{A_s^i+R_s^i}^2\dd s \leq2T\int_0^t\left(\abs{A_s^i}^2+
       \abs{R_s^i}^2\right)\dd s.
\end{align*}
Taking expectations, using \eqref{eq:A-bound} and \eqref{eq:R-bound}, and then maximizing over $i$, we have 
\begin{align*}
 D(t)
 &\leq8TL_b^2\int_0^tD(s)\dd s
   +2T\int_0^t\frac{C_T}{N}\dd s \leq C_T\int_0^tD(s)\dd s+\frac{C_T}{N}.
\end{align*}
Gronwall's inequality proves \eqref{eq:poc-rate}.  Also,
\[
 \Law(\overline X_t^1,\ldots,\overline X_t^k)=\mu_t^{\otimes k},
\]
and $(\overline X_t^1,\ldots,\overline X_t^k)$ is coupled to
$(X_t^1,\ldots,X_t^k)$.  Therefore
\[
 W_2^2\left(\Law(X_t^1,\ldots,X_t^k),\mu_t^{\otimes k}\right)
 \leq\sum_{i=1}^k\E\abs{X_t^i-\overline X_t^i}^2,
\]
which gives \eqref{eq:k-chaos-rate} by applying \eqref{eq:poc-rate} to each term in the sum.
\end{proof}

\begin{remark}[Chaos does not mean finite particles do not interact]
At finite $N$, the $X^i$ are correlated.  ``Propagation of chaos'' says that
every fixed collection becomes asymptotically independent as the surrounding
population grows.  It does not assert independence of a number of particles
that grows proportionally with $N$.
\end{remark}

\begin{corollary}[Convergence of empirical measures]
\label{cor:empirical-measure-convergence}
Under the assumptions of Theorem~\ref{thm:poc}, for every $T<\infty$, the empirical measure $\mu_t^N$ in \eqref{eq:empirical-measure-time} converges to the deterministic limit $\mu_t=\Law(\overline X_t)$ in the sense that
\begin{equation}\label{eq:uniform-empirical-convergence}
 \sup_{0\leq t\leq T}W_2(\mu_t^N,\mu_t)
 \longrightarrow0
 \qquad\text{in probability}.
\end{equation}
Also, $\mu_\cdot\in ([0,T];\cP_2(\R^d))$ and satisfies the mean-field equation
\eqref{eq:mv-weak}.
\end{corollary}
\begin{proof} Let $\overline X_t^i$ be the independent copies of the McKean--Vlasov SDE \eqref{eq:mv-sde}. 
Define the empirical measure of the independent nonlinear processes by
\[
 \overline\mu_t^N:=\frac1N\sum_{i=1}^N\delta_{\overline X_t^i}.
\]
Pairing $X_t^i$ with $\overline X_t^i$ gives a coupling of
$\mu_t^N$ and $\overline\mu_t^N$: 
\[
 \pi_t^N:=\frac1N\sum_{i=1}^N
 \delta_{(X_t^i,\overline X_t^i)}
 \in\Gamma(\mu_t^N,\overline\mu_t^N).
\]
Using this admissible coupling in the definition of $W_2$ gives, pathwise,
\[
 W_2^2(\mu_t^N,\overline\mu_t^N)
 \leq\frac1N\sum_{i=1}^N\abs{X_t^i-\overline X_t^i}^2.
\]
 Taking the supremum in $t$ and using
$\sup_t\sum_i f_i(t)\leq\sum_i\sup_t f_i(t)$ yields
\[
 \sup_{0\leq t\leq T}W_2^2(\mu_t^N,\overline\mu_t^N)
 \leq\frac1N\sum_{i=1}^N
 \sup_{0\leq t\leq T}\abs{X_t^i-\overline X_t^i}^2.
\]
Taking expectations and using \eqref{eq:poc-rate} gives
\begin{equation}\label{eq:coupled-empirical-error}
 \E\sup_{0\leq t\leq T}
 W_2^2(\mu_t^N,\overline\mu_t^N)\leq\frac{C_T}{N}.
\end{equation}
Hence this first distance converges to zero in probability.

It remains to compare $\overline\mu_t^N$ with $\mu_t$.  For each $t$, $\overline\mu_t^N$ is the empirical measure of $N$ i.i.d.\ samples from $\mu_t$.  The law of large numbers gives $W_2^2(\mu_t,\overline\mu_t^N)\to 0$ in probability (in fact, Lemma~\ref{lem:empirical-wasserstein-slln} proves almost-sure convergence). But we need a uniform-in-time version.  To this end, we lift the problem to path space.
  Let
\[
 \mathsf X=C([0,T];\R^d),\qquad
 \norm{\omega}_\infty=\sup_{0\leq t\leq T}\abs{\omega(t)},
\]
and let $\mathsf P=\Law(\overline X_{\cdot})$ on $\mathsf X$.  The moment bound
\eqref{eq:mv-moment} implies $\mathsf P\in\cP_2(\mathsf X)$.  The paths
$\overline X_{\cdot}^1,\overline X_{\cdot}^2,\ldots$ are i.i.d.\ with law
$\mathsf P$, so their empirical path measure
\[
 \mathsf P^N=\frac1N\sum_{i=1}^N
 \delta_{\overline X_{\cdot}^i}
\]
converges almost surely to $\mathsf P$ in $W_2$ on $\mathsf X$.  Here are the
two ingredients behind this statement.  First, $\mathsf X$ is a separable
Polish space, so it admits a countable convergence-determining family
$\{f_m\}_{m\geq1}\subset C_b(\mathsf X)$.  For each $m$, the scalar strong law
gives
\[
 \int_{\mathsf X}f_m\,\dd\mathsf P^N
 =\frac1N\sum_{i=1}^Nf_m(\overline X_{\cdot}^i)
 \longrightarrow\E f_m(\overline X_{\cdot})
 =\int_{\mathsf X}f_m\,\dd\mathsf P
 \qquad\text{almost surely}.
\]
Intersecting these probability-one events over the countable family proves
$\mathsf P^N\Rightarrow\mathsf P$ almost surely.  This is only weak
convergence; by itself it does not control quadratic transportation cost.

Second, \eqref{eq:mv-moment} makes
$\norm{\overline X_{\cdot}}_\infty^2$ integrable, so another application of
the scalar strong law gives convergence of the second moments:
\[
 \frac1N\sum_{i=1}^N\norm{\overline X_{\cdot}^i}_\infty^2
 \longrightarrow
 \E\norm{\overline X_{\cdot}}_\infty^2.
\]
On a separable Banach space, weak convergence together with convergence of
these second moments is equivalent to $W_2$ convergence
\cite[Theorem~6.9]{Villani2009}, so $W_{2,\mathsf X}(\mathsf P^N,\mathsf P)\to 0$ as $N\to\infty$.  We need this stronger conclusion because all
evaluation maps $e_t$ are simultaneously $1$-Lipschitz for the supremum norm;
the resulting Wasserstein contraction below controls every $t\in[0,T]$ with
one path-space bound.  Weak convergence alone would not yield this uniform
$W_2$ estimate.

For the evaluation map $e_t(\omega)=\omega(t)$,
$(e_t)_\#\mathsf P^N=\overline\mu_t^N$ and
$(e_t)_\#\mathsf P=\mu_t$.  Since
$\abs{e_t(\omega)-e_t(\eta)}\leq\norm{\omega-\eta}_\infty$, pushing forward
any path-space coupling gives
\[
 \sup_{0\leq t\leq T}W_2(\overline\mu_t^N,\mu_t)
 \leq W_{2,\mathsf X}(\mathsf P^N,\mathsf P)
 \longrightarrow0
 \qquad\text{almost surely}.
\]
The triangle inequality and \eqref{eq:coupled-empirical-error} now prove
\eqref{eq:uniform-empirical-convergence}.

Finally, apply It\^o's formula to $\varphi(\overline X_t)$ for
$\varphi\in C_c^\infty(\R^d)$:
\[
 \varphi(\overline X_t)-\varphi(\overline X_0)
 =\int_0^t\left[
 b[\mu_s](\overline X_s)\cdot\nabla\varphi(\overline X_s)
 +\frac{\sigma^2}{2}\Delta\varphi(\overline X_s)\right]\dd s
 +\sigma\int_0^t\nabla\varphi(\overline X_s)\cdot\dd B_s.
\]
The stochastic integral has zero expectation.  Taking expectations and using
$\mu_s=\Law(\overline X_s)$ gives exactly \eqref{eq:mv-weak}.  The
mean-square continuity established in \eqref{eq:frozen-ms-continuity}, applied
at the fixed point $m=\mu$, gives $\mu\in C([0,T];\cP_2(\R^d))$.
\end{proof}

\begin{remark}[The empirical-measure rate is dimension dependent]
\label{rem:empirical-measure-rate}
The bound \eqref{eq:coupled-empirical-error} is not by itself a rate for
$W_2(\mu_t^N,\mu_t)$, because one must also estimate the sampling error
$W_2(\overline\mu_t^N,\mu_t)$.  At a fixed time $t$, suppose that $\mu_t$ has
a finite $q$th moment, with $q>4$ when $d\leq4$ and
$q>2d/(d-2)$ when $d>4$.  Theorem~1 of
\cite{FournierGuillin2015}, applied with $p=2$, gives
\[
 \E W_2^2(\overline\mu_t^N,\mu_t)\leq C_t r_d(N),
 \qquad
 r_d(N)=
 \begin{cases}
  N^{-1/2},&d<4,\\
  N^{-1/2}\log(1+N),&d=4,\\
  N^{-2/d},&d>4.
 \end{cases}
\]
Here $C_t$ depends on $d,q$, and the $q$th moment of $\mu_t$.  Since
\[
 W_2^2(\mu_t^N,\mu_t)
 \leq2W_2^2(\mu_t^N,\overline\mu_t^N)
     +2W_2^2(\overline\mu_t^N,\mu_t),
\]
\eqref{eq:coupled-empirical-error} yields the fixed-time estimate
\[
 \E W_2^2(\mu_t^N,\mu_t)
 \leq C_TN^{-1}+2C_t r_d(N).
\]
Thus, under generic moment assumptions, the i.i.d.\ sampling error rather
than the propagation-of-chaos coupling is the slower term.  The path-space
argument above proves uniform-in-time convergence but does not provide one of
these finite-dimensional rates: $C([0,T];\R^d)$ is infinite dimensional, and
a quantitative path-space estimate requires additional control of temporal
regularity or metric entropy.
\end{remark}

When the system is deterministic, the empirical measure convergence can be proved directly without the path-space argument.  The following proposition is a
deterministic analogue of Corollary~\ref{cor:empirical-measure-convergence} and is useful for the analysis of the mean-field limit of deterministic particle systems. 
Note that this result is not a consequence of Theorem~\ref{thm:poc} because the initial positions are not assumed to be i.i.d., but rather an arbitrary deterministic point cloud. It is a deterministic statement about the stability of the characteristic flow, and it is also a stronger statement than Corollary~\ref{cor:empirical-measure-convergence} because it gives a uniform-in-time bound (on the Wasserstein distance between the empirical measure and the mean-field limit), whereas Corollary~\ref{cor:empirical-measure-convergence} only gives a convergence. 
The proof is based on the stability of the characteristic flow, and it is simpler than the proof of Theorem~\ref{thm:poc} because there is no stochastic term to control.  

\begin{proposition}[Deterministic empirical approximation]
\label{prop:deterministic-empirical-convergence}
 Assume $\sigma=0$ and let $b$ satisfy Assumption~\ref{ass:regular-mf}.  Let
$\mu_0^N=N^{-1}\sum_{i=1}^N\delta_{x_0^{i,N}}\in\cP_2(\R^d)$ be deterministic,
let $\mu_t^N$ be the empirical measure of the particle ODE
\eqref{eq:particles}, and let $\mu_t$ solve the mean-field continuity equation \eqref{eq:mv-weak} 
with initial law $\mu_0\in\cP_2(\R^d)$.  Then
\begin{equation}\label{eq:deterministic-empirical-stability}
 \sup_{0\leq t\leq T}W_2(\mu_t^N,\mu_t)
 \leq e^{2L_bT}W_2(\mu_0^N,\mu_0).
\end{equation}
Consequently, $W_2(\mu_0^N,\mu_0)\to0$ implies uniform-in-time convergence of
the deterministic empirical measures.  If the initial points are instead
i.i.d.\ with law $\mu_0$, the same conclusion holds almost surely.
\end{proposition}
\begin{proof}
Let $\pi_0\in\Gamma(\mu_0^N,\mu_0)$ be an optimal coupling that achieves the minimum in the definition of the Wasserstein distance. Here we
use \emph{one-particle characteristic maps}, not the flow on the full
configuration space $(\R^d)^N$ (For a prescribed measure curve $(m_t)_{t\geq0}$, its one-particle characteristic map $f_t^m$ sends an initial position $x$ to the endpoint at
time $t$ of the nonautonomous ODE $  \dot z_s=b[m_s](z_s)$ with $z_0=x$.
Thus the measure curve determines the velocity field, and the characteristic
map transports each initial point along that field.)
Once the two measure curves $\mu_t^N$ and
$\mu_t$ are fixed, define characteristic maps $f_t^N,f_t:\R^d\to\R^d$ by
\begin{align*}
 \frac{\dd}{\dd t}f_t^N(x) &=b[\mu_t^N](f_t^N(x)),\qquad f_0^N(x)=x, \\
 \frac{\dd}{\dd t}f_t(y)   &=b[\mu_t](f_t(y)),\qquad f_0(y)=y.
\end{align*}
In particular, for every initial atom $x_0^{i}$, the particle equation reads
\[
 \dot X_t^{i}
 =\frac1N\sum_{j=1}^N b(X_t^{i},X_t^{j})
 =b[\mu_t^N](X_t^{i}),\qquad X_0^{i}=x_0^{i}.
\]
This is exactly the characteristic ODE defining $f_t^N(x_0^{i})$.
Uniqueness of that ODE therefore gives $X_t^{i}=f_t^N(x_0^{i})$.
Hence
$(f_t^N)_\#\mu_0^N=\mu_t^N$.  Similarly,
$(f_t)_\#\mu_0=\mu_t$.  Therefore,
\[
\pi_t=(f_t^N,f_t)_\#\pi_0\in\Gamma(\mu_t^N,\mu_t).
\]

Let $(X_0,Y_0)$ have law $\pi_0$ and set $X_t=f_t^N(X_0)$ and $Y_t=f_t(Y_0)$.  Then $(X_t,Y_t)$ has law $\pi_t$ and
\[ D(t):=\int_{\R^d\times\R^d}\abs{x-y}^2\pi_t(\dd x,\dd y)= \E|X_t-Y_t|^2.
\]
The chain rule gives
\[
 D'(t)=2\int (x-y)\cdot
 \bigl(b[\mu_t^N](x)-b[\mu_t](y)\bigr)\,\pi_t(\dd x,\dd y).
\]
Because the first and second marginals of $\pi_t$ are $\mu_t^N$ and
$\mu_t$, respectively,
\[
 b[\mu_t^N](x)-b[\mu_t](y)
 =\int\bigl[b(x,x')-b(y,y')\bigr]  \pi_t(\dd x',\dd y').
\]
Substitution and Fubini's theorem give the first equality:
\begin{align*}
 D'(t)
 &=2\iint (x-y)\cdot
   \bigl[b(x,x')-b(y,y')\bigr]\,
   \pi_t(\dd x,\dd y)\pi_t(\dd x',\dd y')\\
 &\leq2L_b\iint\abs{x-y}
   \bigl(\abs{x-y}+\abs{x'-y'}\bigr)\,
   \pi_t(\dd x,\dd y)\pi_t(\dd x',\dd y')\\
 &\leq4L_bD(t),
\end{align*}
where the last step uses
$(\int\abs{x-y}\dd\pi_t)^2\leq D(t)$.  Gronwall's inequality yields
$D(t)\leq e^{4L_bt}D(0)$.  Since $\pi_t$ is an admissible coupling and
$D(0)=W_2^2(\mu_0^N,\mu_0)$, taking square roots proves
\eqref{eq:deterministic-empirical-stability}.  The last assertion follows
from Lemma~\ref{lem:empirical-wasserstein-slln}.
\end{proof}

%%%%%%%%%%%===============
\subsection{Examples}

\begin{example}[linear consensus]
Let $b(x,y)=y-x$. Then, the system is 
$$
dX_t^i= \frac 1 N \sum_{j=1}^N (X_j^i - X_t^i) \dd t+\sigma\dd B_t = [ - X_t^i + \frac 1 N  \sum_{j=1}^N  X_t^j] \dd t+\sigma\dd B_t.
$$
  For the deterministic system, define
\[
  m_t^N=\frac1N\sum_iX_t^i,
  \qquad
  \mathcal V_t^N=\frac1N\sum_i\abs{X_t^i-m_t^N}^2.
\]
Because the pair forces cancel after summation, $\dot m_t^N=0$.  Moreover,
$\dot X_t^i=m_0^N-X_t^i$, so
\[
  X_t^i=m_0^N+e^{-t}(X_0^i-m_0^N),
  \qquad \mathcal V_t^N=e^{-2t}\mathcal V_0^N.
\]

At the mean-field level the drift is $m_t-x$, where
$m_t=\int x\,\mu_t(\dd x)$ is constant. The McKean--Vlasov SDE is
\[
  \dd\overline X_t=(m_0-\overline X_t)\dd t+\sigma\dd B_t, 
\] 
which is an Ornstein--Uhlenbeck process.  The law $\mu_t$ is Gaussian with mean $m_0$ and variance $\sigma^2(1-e^{-2t})/2$. 
The centered second moment
$S_t=\int\abs{x-m_t}^2\mu_t(\dd x)$ obeys
\[
  \dot S_t=-2S_t+d\sigma^2,
  \qquad
  S_t=e^{-2t}S_0+\frac{d\sigma^2}{2}(1-e^{-2t}).
\]
Attraction tries to collapse the law, while diffusion maintains a nonzero
width.  This balance reappears in noisy attention and generative samplers.

\end{example}


\begin{exercise}[Mean-field scaling and its alternatives] TBV. 
Replace $1/N$ in \eqref{eq:particles} by $N^{-\gamma}$ for the consensus
kernel.  Determine the relaxation time $\tau_N$, defined by
$e^{-t/\tau_N}$ being the contraction factor of a centered particle
deviation in the deterministic system, as a function of $N$ and $\gamma$.
Which value produces order-one interaction dynamics (i.e., on
fixed times $t=O(1)$, this contraction factor has a nontrivial limit in
$(0,1)$, or equivalently that $\tau_N$ stays bounded away from both zero and
infinity) as $N\to\infty$?
\emph{Hint:} compute the equation for $X^i-m^N$.
\end{exercise}
\begin{proof}[Solution]
For $b(x,y)=y-x$, the rescaled system is
\begin{equation}\label{eq:scaled-consensus}
 \dd X_t^i=N^{1-\gamma}(m_t^N-X_t^i)\dd t+\sigma\dd B_t^i.
\end{equation}
First set $\sigma=0$.  The empirical mean is conserved because
\[
 \frac1N\sum_{i=1}^N N^{-\gamma}
       \sum_{j=1}^N(X_t^j-X_t^i)=0.
\]
Thus $Y_t^i:=X_t^i-m_t^N$ solves
\[
 \dot Y_t^i=-N^{1-\gamma}Y_t^i,
 \qquad
 Y_t^i=e^{-t N^{1-\gamma}}Y_0^i.
\]
The $e$-folding relaxation time for particle deviations is therefore
\begin{equation}\label{eq:relaxation-time}
 \tau_N=N^{\gamma-1}.
\end{equation}
The empirical variance decays as
$\mathcal V_t^N=e^{-2N^{1-\gamma}t}\mathcal V_0^N$, so its $e$-folding time is
$\tau_N/2$; both notions have the same $N$-scaling.

If $\sigma>0$, then
\[
 \dd m_t^N=\frac{\sigma}{N}\sum_{j=1}^N\dd B_t^j,
 \qquad
 \dd Y_t^i=-N^{1-\gamma}Y_t^i\dd t
 +\sigma\left(\dd B_t^i-\frac1N\sum_{j=1}^N\dd B_t^j\right).
\]
Thus noise adds fluctuations but does not change the relaxation rate of the
initial deviation. 

The precise statement is obtained by writing
$\lambda_N=N^{1-\gamma}=\tau_N^{-1}$.  When $\sigma=0$,
\[
 \max_{1\leq i\leq N}\abs{Y_t^i}
 =e^{-\lambda_Nt}\max_{1\leq i\leq N}\abs{Y_0^i}.
\]
Hence, if $\gamma<1$, then for every fixed $t>0$ the right side converges to
zero, provided the initial deviations are uniformly bounded.  The nontrivial
decay profile is visible on the boundary-layer time $t=\tau_Ns$, since
$Y_{\tau_Ns}^i=e^{-s}Y_0^i$.  If $\gamma>1$, then for every fixed $T<\infty$,
\[
 \sup_{0\leq t\leq T}\max_i\abs{Y_t^i-Y_0^i}
 \leq\bigl(1-e^{-\lambda_NT}\bigr)\max_i\abs{Y_0^i}
 \longrightarrow0,
\]
so the interaction has no leading-order effect on that time interval.  For
$\gamma=1$, $Y_t^i=e^{-t}Y_0^i$, which is a nontrivial order-one contraction.
Thus $\boxed{\gamma=1}$ is the unique critical scaling.

For fixed $\sigma>0$, the term \emph{frozen} refers only to the interaction,
not to the full noisy dynamics.  Indeed, conditional on $Y_0^i$,
\[
 \E\!\left[\abs{Y_t^i}^2\mid Y_0^i\right]
 =e^{-2\lambda_Nt}\abs{Y_0^i}^2+{}
 \frac{d\sigma^2(1-N^{-1})}{2\lambda_N}
   \bigl(1-e^{-2\lambda_Nt}\bigr).
\]
Thus $\gamma<1$ still gives mean-square consensus at each fixed $t>0$ under
uniformly bounded initial second moments, whereas for $\gamma>1$ the
interaction vanishes but the centered Brownian fluctuations remain.
\end{proof}

\begin{exercise}[Weak equation with state-dependent diffusion]
Let the noise be $\Sigma(X_t^i,\mu_t^N)\dd B_t^i$, where
\[
 \Sigma:\R^d\times\cP_2(\R^d)\longrightarrow\R^{d\times d}.
\]
Derive the analogue of \eqref{eq:empirical-weak}.
\end{exercise}
\begin{proof}[Solution]
Set
\[
 a(x,\mu):=\Sigma(x,\mu)\Sigma(x,\mu)^\top.
\]
It\^o's formula, applied before averaging, gives
\begin{align*}
 \dd\varphi(X_t^i)
 &=\nabla\varphi(X_t^i)\cdot b[\mu_t^N](X_t^i)\dd t \quad+\frac12\operatorname{tr}\!\left(
 a(X_t^i,\mu_t^N)D^2\varphi(X_t^i)\right)\dd t\\
 &\quad+\nabla\varphi(X_t^i)^\top
 \Sigma(X_t^i,\mu_t^N)\dd B_t^i.
\end{align*}
Therefore
\begin{align}
 \dd\int\varphi\,\dd\mu_t^N
 &=\int\left[b[\mu_t^N]\cdot\nabla\varphi
 +\frac12\operatorname{tr}\!\left(a(\,\cdot\,,\mu_t^N)
 D^2\varphi\right)\right]\dd\mu_t^N\dd t
 +\dd M_t^{N,\varphi},\label{eq:state-dependent-weak}\\
 M_t^{N,\varphi}
 &=\frac1N\sum_{i=1}^N\int_0^t
 \nabla\varphi(X_s^i)^\top\Sigma(X_s^i,\mu_s^N)\dd B_s^i.
 \label{eq:state-dependent-martingale}
\end{align}
Since the Brownian motions are independent, the quadratic variation is
\begin{equation}\label{eq:state-dependent-bracket}
 \left\langle M^{N,\varphi}\right\rangle_t
 =\frac1{N^2}\sum_{i=1}^N\int_0^t
 \abs{\Sigma(X_s^i,\mu_s^N)^\top\nabla\varphi(X_s^i)}^2\dd s.
\end{equation}
In particular, if
$\sup_{x,\mu}\abs{\Sigma(x,\mu)^\top\nabla\varphi(x)}<\infty$, the martingale
is again of order $N^{-1/2}$ in $L^2$.  Formally passing to a deterministic
limit gives
\[
 \partial_t\rho
 =-\diver\bigl(\rho b[\rho]\bigr)
 +\frac12\sum_{k,\ell=1}^d
 \partial_{k\ell}\bigl(a_{k\ell}(x,\rho)\rho\bigr).
\]
There are no derivatives of $\Sigma$ in the weak generator; in the density
equation the derivatives act on the full product $a_{k\ell}\rho$.
\end{proof}





\paragraph{\textbf{Literature and storyline.}}
The modern mean-field viewpoint begins with Kac's stochastic model for a
dilute gas \cite{Kac1956}.  His decisive simplification was to retain the
many-particle symmetry and collision mechanism while discarding enough
microscopic detail to make the large-$N$ question mathematically visible.
In Kac's terminology, ``chaos'' is the asymptotic product structure of fixed
particle marginals, not disorder in the trajectories.  His autobiography
recalls that, even as a student, he resisted a formal derivation until he
understood its motivation; the same instinct later informed his use of simple
probabilistic models \cite{Kac1985Enigmas}.  McKean
then connected nonlinear parabolic equations with law-dependent Markov
processes \cite{McKean1966}; this is the conceptual origin of the
McKean--Vlasov SDE used above.  Sznitman's Saint-Flour notes made coupling and
propagation of chaos into a reusable proof framework \cite{Sznitman1991}.
The mathematical storyline of this section follows that history: first write
one particle's drift through the empirical law, then identify the nonlinear
limit process, and finally prove that finitely many tagged particles become
asymptotically independent.  The point is not merely to derive a PDE; it is to
justify when a population can be replaced by one statistically representative
nonlinear particle.


\subsection{Remarks and connections forward}
\begin{itemize}[leftmargin=*]
  \item We have three scales: microscopic (particles), mesoscopic (empirical measure), and macroscopic (mean-field PDE):
 \begin{align*}
   \text{Micro-scale }  \{X_t^i\}_{i=1}^N  \longrightarrow \text{Meso-scale }   \mu_t^N \longrightarrow \text{Macro-scale } \mu_t .
 \end{align*}
  % \begin{tikzpicture}
  %     \node (micro) at (0,0) {Microscopic};
  %     \node (meso) at (2,0) {Mesoscopic};
  %     \node (macro) at (4,0) {Macroscopic};
  %   \end{tikzpicture}

  \item Equation~\eqref{eq:mv-fp} is already a generative-dynamics template:
  a learned drift transports probability while Brownian noise smooths it.
  \item The map $\mu_0\mapsto\mu_t$ is a nonlinear solution operator.  Learning
  it from many initial conditions is an operator-learning problem.
  \item The empirical weak identity \eqref{eq:empirical-weak} will become a
  derivative-free regression equation for the inverse problem in
  Section~\ref{sec:learning-ips}.
  \item In a Transformer, tokens are not i.i.d.\ at later layers; the point of
  the theorem is that weak mean-field interaction can nevertheless create an
  asymptotically independent nonlinear description under the stated model.
\end{itemize}



%%%%%%%%%%%%%===========
\section{Gradient flow and optimal transport}\label{sec:gradient-flow}

\paragraph{\textbf{Key question}}
When does an IPS have a decreasing energy (that is, when is it a gradient
flow),
and why is the relevant geometry on probability measures supplied by optimal
transport?

The passage from particles to measures in Section~\ref{sec:ips} does more than reduce a large system to one equation.  It reveals a geometry.  A velocity
field moves mass through the continuity equation; optimal transport measures
the least kinetic action required by such a motion; and, when the velocity is
generated by an energy, the evolution becomes steepest descent in this
transport geometry.  Thus a particle ODE, a nonlinear PDE, and an optimization
principle can be three descriptions of the same dynamics.

\subsection{The continuity equation as a path in probability space}

Consider a sufficiently regular curve $(\mu_t)_{0\leq t\leq T}$ in
$\cP_2(\R^d)$ and a velocity field $v_t:\R^d\to\R^d$ satisfying
\begin{equation}\label{eq:gf-continuity}
 \partial_t\mu_t+\diver(\mu_t v_t)=0.
\end{equation}
The equation says that mass is neither created nor destroyed: it is transported
with instantaneous velocity $v_t$.
Recall that for $0\leq s\leq t\leq T$, the \emph{characteristic flow}
$F_{s,t}:\R^d\to\R^d$ is defined pointwise by
\[
 \frac{\dd}{\dd r}F_{s,r}(x)=v_r(F_{s,r}(x)),\qquad F_{s,s}(x)=x.
\]
Thus $F_{s,t}(x)$ is the position at time $t$ of a particle placed at $x$ at
time $s$.  Uniqueness gives the flow property
$F_{r,t}\circ F_{s,r}=F_{s,t}$; we abbreviate $F_t=F_{0,t}$.
If $v_t$ is Lipschitz in space, Proposition~\ref{prop:appendix-characteristics}
gives $\mu_t=(F_t)_\#\mu_0$.  Conversely, differentiating this pushforward
identity gives the weak form of \eqref{eq:gf-continuity}.  This is precisely
the deterministic part of the mean-field equation in
Section~\ref{sec:ips}, now viewed as
a curve in the space of probability measures.

The instantaneous kinetic energy of this motion is
\begin{equation}\label{eq:kinetic-action-density}
 K(t)= \norm{v_t}_{L^2(\mu_t)}^2
 :=\int_{\R^d}\abs{v_t(x)}^2\,\mu_t(\dd x),
\end{equation}
and its action over $[0,T]$ is the time integral of
\eqref{eq:kinetic-action-density}.  The measure $\mu_t$ specifies where mass
is present, while $v_t$ specifies how that mass moves.  The pair
$(\mu_t,v_t)$, rather than either component alone, is the natural dynamical
object.

\begin{theorem}[Benamou--Brenier dynamic formulation]
\label{thm:benamou-brenier}
For $\mu_0,\mu_1\in\cP_2(\R^d)$,
\begin{equation}\label{eq:benamou-brenier}
 W_2^2(\mu_0,\mu_1)
 =\inf_{(\mu,v)}
   \int_0^1\int_{\R^d}\abs{v_t(x)}^2\,\mu_t(\dd x)\dd t,
\end{equation}
where the infimum is over narrowly continuous curves
$\mu:[0,1]\to\cP_2(\R^d)$ and Borel velocity fields $v$ with finite action
that satisfy \eqref{eq:gf-continuity} in the distributional sense and have
the prescribed endpoints.
\end{theorem}

See \cite{BenamouBrenier2000} and
\cite[Chapter~7]{Villani2009}.  The static definition of $W_2$ minimizes the
quadratic displacement over couplings of two endpoint measures.  Theorem
\ref{thm:benamou-brenier} gives the complementary dynamic picture: the same quantity is the least kinetic action among all mass-preserving paths joining the endpoints (hence it is a geodesic). 
 In particular, $W_2$ is not merely a convenient metric for
weak convergence; it measures the cost of motion allowed by the continuity
equation.

Furthermore, let $X_0\sim\mu_0$, $X_1\sim\mu_1$ be jointly distributed random variables with law $\pi$, and $(\mu_t,v_t)$ be the corresponding curve in the space of probability measures and velocity fields satisying the continuity equation. Then, $\dot X_t = v_t(X_t)$ and 
\[
 W_2^2(\mu_0,\mu_1)
 \leq\E\abs{X_0-X_1}^2 = \E |\int_0^1 \dot X_tdt|^2  
 = \E |\int_0^1 v_t(X_t)dt|^2
 \leq \E \int_0^1 |v_t(X_t)|^2 dt.  
\]
There can be many curves and velocity fields connecting the endpoints. One curve with constant velocity $v_t \equiv X_1 - X_0$ is $X_t = (1-t)X_0 + tX_1$. The next subsection shows that when $\pi$ is the optimal coupling, this curve is the constant-speed geodesic. 


\vspace{2mm}
\paragraph{\textbf{Interpolation and geodesics.}}
If $\pi\in \Gamma(\mu_0,\mu_1)$ is an optimal coupling (i.e., it attains the infimum in the definition of $W_2$: $\int_{\R^d\times\R^d}|x-y|^2\,\pi(\dd x,\dd y)
 =W_2^2(\mu_0,\mu_1)$), then the displacement interpolation
$T_t:\R^d\times\R^d\to\R^d$ defined by $T_t(x,y)=(1-t)x+ty$ gives a constant-speed $W_2$ geodesic in $\cP_2(\R^d)$. Specifically, let
\begin{equation}\label{eq:displacement-interpolation}
 \mu_t=(T_t)_\#\pi,
 \qquad 0\leq t\leq1,
\end{equation}
Then, the curve $(\mu_t)_{0\leq t\leq1}$ satisfies
\[
 W_2(\mu_s,\mu_t)=|t-s|W_2(\mu_0,\mu_1),
 \qquad 0\leq s,t\leq1.
\]

To prove it, take a \emph{jointly distributed}
pair $(X_0,X_1)$ with law $\pi$. Thus, $X_0\sim\mu_0$ and
$X_1\sim\mu_1$, and they are coupled optimally (rather than chosen
independently) so that $W_2(\mu_0,\mu_1)=\bigl(\E\abs{X_0-X_1}^2\bigr)^{1/2}$.  Set
\[
 X_t=(1-t)X_0+tX_1.
\]
The pushforward--law correspondence in
Lemma~\ref{lem:pushforward-random-variable} gives
$\Law(X_t)=(T_t)_\#\pi=\mu_t$.  Hence $(X_s,X_t)$ is an admissible coupling
of $(\mu_s,\mu_t)$, and
\begin{equation}\label{eq:geodesic-upper-bound}
 W_2(\mu_s,\mu_t)
 \leq\bigl(\E\abs{X_s-X_t}^2\bigr)^{1/2} = |t-s| \bigl(\E\abs{X_0-X_1}^2\bigr)^{1/2}
 =\abs{t-s}W_2(\mu_0,\mu_1).
\end{equation}
To obtain the reverse inequality, suppose $0\leq s<t\leq1$.  Apply
\eqref{eq:geodesic-upper-bound} to the two end intervals and use the triangle
inequality:
\begin{align*}
 W_2(\mu_0,\mu_1)
 &\leq W_2(\mu_0,\mu_s)+W_2(\mu_s,\mu_t)
       +W_2(\mu_t,\mu_1)\\
 &\leq sW_2(\mu_0,\mu_1)+W_2(\mu_s,\mu_t)
       +(1-t)W_2(\mu_0,\mu_1).
\end{align*}
Rearranging gives
$W_2(\mu_s,\mu_t)\geq(t-s)W_2(\mu_0,\mu_1)$, which, together with the
upper bound, proves equality.

When the optimal coupling is induced by
a map $T$, meaning $\pi=(\Id,T)_\#\mu_0$, one may equivalently take
$X_0\sim\mu_0$ and $X_1=T(X_0)$.  The pair $(X_0,X_1)$ then has law
$(\Id,T)_\#\mu_0$, whose first marginal is $\mu_0$ and whose second marginal
is $T_\#\mu_0=\mu_1$.  By the composition rule for pushforwards in
Lemma~\ref{lem:pushforward-random-variable},
\[
 \mu_t=(T_t)_\#(\Id,T)_\#\mu_0
 =\bigl(T_t\circ(\Id,T)\bigr)_\#\mu_0
 =((1-t)\Id+tT)_\#\mu_0.
\]
Thus particles travel on straight lines, but the
resulting path of densities is generally nonlinear.  This interpolation of
mass, rather than linear interpolation of densities, is the basic geometry
behind Wasserstein gradient flow.




\vspace{3mm}

\paragraph{\textbf{Kinetic action and gradient flow.}}  Gradient flow is steepest descent of an energy function in a metric space.  In Euclidean space, the steepest descent direction is the negative gradient, and the steepest descent curve is the solution of the ODE $\dot z_t=-\nabla\mathcal E(z_t)$.  The steepest descent property is equivalent to the energy dissipation identity
\[
 \frac{\dd}{\dd t}\mathcal E(z_t)
 =D\mathcal E(z_t)[\dot z_t]
 =-\norm{\nabla\mathcal E(z_t)}^2
 =-\norm{\dot z_t}^2.
\]  
The energy dissipation identity is a fundamental property of gradient flows, indicating that the rate of change of the energy is equal to the negative square of the norm of the tangent vector to the flow.

Note that a free energy ranks states, but it does not by itself define a gradient
flow: the word \emph{steepest} has meaning only after a metric has specified
the size of a tangent displacement.  On a Riemannian manifold
$(\mathcal M,g)$, the gradient is characterized by
\[
 D\mathcal E(z)[\xi]
 =\langle \operatorname{grad}\mathcal E(z),\xi\rangle_z,
 \qquad
 \dot z_t=-\operatorname{grad}\mathcal E(z_t).
\]
Our aim is to formulate the same construction for a free energy
$\mathcal E:\cP_2(\R^d)\to(-\infty,+\infty]$, using $W_2$ to measure
motion.

\emph{The tangent metric supplied by $W_2$.}
The space $\cP_2(\R^d)$ consists of probability measures with finite second
moment. Rigorously, $(\cP_2(\R^d),W_2)$ is a geodesic metric space.
Otto's calculus interprets it, at sufficiently regular measures, as a formal
infinite-dimensional Riemannian manifold with tangent space
\[
 T_\mu\cP_2(\R^d)
 :=\overline{\left\{\nabla\zeta:\zeta\in C_c^\infty(\R^d)\right\}}
 ^{L^2(\mu;\R^d)}
\]
and inner product
\[
 \langle v,w\rangle_\mu
 :=\int_{\R^d}v(x)\cdot w(x)\,\mu(\dd x).
\]

This tangent norm is intrinsic to $W_2$.  If $(\mu_t)$ is an absolutely
continuous curve in $(\cP_2(\R^d),W_2)$, then for almost every $t$ there is
a unique (exercise: verify it) minimal-norm tangent velocity $v_t$ satisfying
\[
 \partial_t\mu_t+\diver(\mu_t v_t)=0.
\]
Its norm equals the metric derivative:
\begin{equation}\label{eq:w2-metric-speed}
 \abs{\dot\mu_t}_{W_2}
 :=\lim_{h\to0}\frac{W_2(\mu_{t+h},\mu_t)}{\abs h}
 =\norm{v_t}_{L^2(\mu_t)}.
\end{equation}
Other velocity fields can represent the same curve, but have no smaller
$L^2(\mu_t)$ norm.  Consequently,
\[
 K(t)=\int_{\R^d}\abs{v_t}^2\dd\mu_t
 =\abs{\dot\mu_t}_{W_2}^2
\]
for the minimal velocity.  The time integral of $K(t)$ is therefore the
squared-speed action of the curve, and the Benamou--Brenier formula
\eqref{eq:benamou-brenier} says that minimizing this action over curves with
fixed endpoints gives $W_2^2$.  This is the geometric reason that kinetic
action, rather than an arbitrary penalty, enters Wasserstein gradient flow.
The rigorous metric theory behind this formal manifold picture is developed
in \cite[Chapters~8, 10--11]{AmbrosioGigliSavare2008}.

\emph{Why fastest decrease is ill posed.}
Suppose that $\mathcal E$ has a sufficiently regular first variation
\[
 \phi_\mu(x):=\frac{\delta\mathcal E}{\delta\mu}[\mu](x),
 \qquad
 \nabla\phi_\mu\in T_\mu\cP_2(\R^d).
\]
For a smooth curve with $\mu_0=\mu$ and initial tangent velocity $v$, the
continuity equation and integration by parts give
\begin{equation}\label{eq:wasserstein-directional-energy}
 D\mathcal E(\mu)[v]
 :=\left.\frac{\dd}{\dd t}\mathcal E(\mu_t)\right|_{t=0}
 =\int_{\R^d}\nabla\phi_\mu(x)\cdot v(x)\,\mu(\dd x).
\end{equation}
Thus the formal $W_2$ gradient of $\mathcal E$ is
$\operatorname{grad}_{W_2}\mathcal E(\mu)=\nabla\phi_\mu$.

Maximizing \eqref{eq:wasserstein-directional-energy} would select ascent, not
descent.  One might instead minimize it, or maximize its negative, but this
problem is unbounded whenever $\nabla\phi_\mu\neq0$.  Indeed, for
$v=-c\nabla\phi_\mu$,
\[
 D\mathcal E(\mu)[v]
 =-c\int_{\R^d}\abs{\nabla\phi_\mu}^2\dd\mu
 \longrightarrow-\infty
 \qquad\text{as }c\to\infty.
\]
The energy derivative alone therefore demands an arbitrarily large velocity
and supplies no time scale.  Constraining
$\norm{v}_{L^2(\mu)}\leq1$ selects the normalized descent direction, but the
unit-speed constraint imposes an arbitrary clock.

\emph{The Onsager principle: decrease balanced by transport cost.}
At a fixed state $\mu$, define the dissipation potential
\[
 \Psi_\mu(v):=\frac12\norm{v}_{L^2(\mu)}^2
 =\frac12\int_{\R^d}\abs v^2\dd\mu
\]
and the \emph{Onsager functional}, also called the \emph{Rayleighian}, by
\begin{align}
 \mathcal R_\mu(v)
 &:=\underbrace{\Psi_\mu(v)}_{\text{potential dissipation}}
   +\underbrace{D\mathcal E(\mu)[v]}
      _{\text{directional energy change}} \notag\\
 &=\frac12\int_{\R^d}\abs{v+\nabla\phi_\mu}^2\dd\mu
   -\frac12\int_{\R^d}\abs{\nabla\phi_\mu}^2\dd\mu.
 \label{eq:onsager-functional}
\end{align}
Thus the potential dissipation term is one half of the instantaneous kinetic-action density.  By
\eqref{eq:w2-metric-speed}, it is precisely one half of the squared
$W_2$ speed.  The minimizer 
\[
 v_\mu
 =\argmin_{v\in T_\mu\cP_2(\R^d)}\mathcal R_\mu(v)
\]
is therefore the penalized version of fastest descent, and \eqref{eq:onsager-functional} gives
\[
 v_\mu=-\nabla\phi_\mu
 =-\nabla\frac{\delta\mathcal E}{\delta\mu}[\mu].
\]
The quadratic penalty selects both the direction and its magnitude.  The
coefficient $1/2$ corresponds to unit mobility; changing it rescales the
velocity, and hence the time variable.

\emph{The Wasserstein gradient-flow equation.}
Substituting the minimizing velocity into the continuity equation gives
\begin{equation}\label{eq:kinetic-role-gradient-flow}
 \partial_t\mu_t
 =\diver\left(
   \mu_t\nabla\frac{\delta\mathcal E}{\delta\mu}[\mu_t]\right).
\end{equation}
Along a smooth solution,
\begin{equation}\label{eq:wasserstein-energy-dissipation}
 \frac{\dd}{\dd t}\mathcal E(\mu_t)
 =-\int_{\R^d}
   \left|\nabla\frac{\delta\mathcal E}{\delta\mu}[\mu_t]\right|^2
   \dd\mu_t
 =-\norm{v_t}_{L^2(\mu_t)}^2
 =-\abs{\dot\mu_t}_{W_2}^2\leq0.
\end{equation}
The energy supplies the force $-\nabla\phi_{\mu_t}$, while the $W_2$ tangent
metric converts this force into a velocity. 

\emph{What is guaranteed to decrease?}
Equation~\eqref{eq:wasserstein-energy-dissipation} shows that the free energy
$\mathcal E(\mu_t)$ decreases.  The kinetic quantity
$K(t)=\abs{\dot\mu_t}_{W_2}^2$ need not.  Already for the Euclidean gradient
flow $\dot x_t=-E'(x_t)$,
\[
 \frac12K'(t)=-E''(x_t)K(t).
\]
For example, if $E(x)=1-\cos x$ and $x_t\in(\pi/2,\pi)$, then
$\frac{\dd}{\dd t}E(x_t)=-\sin^2x_t<0$ but
$K'(t)=-2\sin^2x_t\cos x_t>0$.  Formally, the Wasserstein analogue is
\[
 \frac12K'(t)
 =-\operatorname{Hess}_{W_2}\mathcal E(\mu_t)[v_t,v_t].
\]
Thus displacement convexity makes $K(t)$ nonincreasing, but gradient-flow
structure alone does not.  By contrast, the minimal kinetic quantity along
a constant-speed $W_2$ geodesic is constant by definition.

\emph{Connections forward.}
This abstract construction is the template for the examples below.  For the
interaction energy \eqref{eq:interaction-energy}, it identifies a symmetric
potential IPS as both a finite-dimensional gradient descent and the
continuum $W_2$ gradient flow \eqref{eq:wasserstein-interaction-gf}.  Adding
Boltzmann entropy gives the diffusive free energy \eqref{eq:free-energy} and
the McKean--Vlasov/Fokker--Planck equation
\eqref{eq:free-energy-gradient-flow}.  The JKO scheme
\eqref{eq:jko-scheme} is the finite-step version: it replaces the
infinitesimal squared-speed penalty by $W_2^2/(2h)$.

Later, Section~\ref{sec:continuous-depth} applies the same test to attention.
An attention field always defines a continuity equation in the
continuous-depth limit, but it is a standard $W_2$ gradient flow only when
the field has the variational form
\[
 \Phi[\mu](x)
 =-\nabla_x\frac{\delta\mathcal E}{\delta\mu}[\mu](x)
\]
for a single energy $\mathcal E$; generic row-normalized attention need not
have this structure.



%%% 
\subsection{Potential interactions: one gradient flow at two scales}

We now identify the additional structure that turns a general IPS into
a gradient flow.  Let $U,W\in C^2(\R^d)$ and suppose that $W$ is even:
$W(z)=W(-z)$.  Consider the deterministic particle system
\begin{equation}\label{eq:potential-particle-system}
 \dot X_t^i=-\nabla U(X_t^i)
 -\frac1N\sum_{j=1}^N\nabla W(X_t^i-X_t^j),
 \qquad i=1,\ldots,N.
\end{equation}
It is a special case of Section~\ref{sec:ips} with
$b(x,y)=-\nabla U(x)-\nabla W(x-y)$.  Define the energy
\begin{equation}\label{eq:finite-particle-energy}
 \mathcal F_N(X^1,\ldots,X^N)
 =\frac1N\sum_{i=1}^NU(X^i)
 +\frac1{2N^2}\sum_{i,j=1}^NW(X^i-X^j).
\end{equation}

\begin{proposition}[Particle descent and Wasserstein descent]
\label{prop:ips-wasserstein-gradient-flow}
For the potential system \eqref{eq:potential-particle-system},
\begin{equation}\label{eq:particle-gradient-flow}
 \dot X_t^i=-N\nabla_{X^i}\mathcal F_N(X_t^1,\ldots,X_t^N),
\end{equation}
and therefore
\begin{equation}\label{eq:particle-energy-dissipation}
 \frac{\dd}{\dd t}\mathcal F_N(X_t^1,\ldots,X_t^N)
 =-\frac1N\sum_{i=1}^N\abs{\dot X_t^i}^2\leq0.
\end{equation}
At the level of measures, define
\begin{equation}\label{eq:interaction-energy}
 \mathcal F(\mu)
 =\int_{\R^d}U(x)\,\mu(\dd x)
 +\frac12\iint_{\R^d\times\R^d}
 W(x-y)\,\mu(\dd x)\mu(\dd y).
\end{equation}
Its formal first variation is
\begin{equation}\label{eq:interaction-first-variation}
 \frac{\delta\mathcal F}{\delta\mu}[\mu](x)
 =U(x)+(W*\mu)(x),
\end{equation}
and the mean-field equation associated with
\eqref{eq:potential-particle-system} is
\begin{equation}\label{eq:wasserstein-interaction-gf}
 \partial_t\mu_t
 =\diver\left(\mu_t\nabla
     \frac{\delta\mathcal F}{\delta\mu}[\mu_t]\right)
 =\diver\bigl(\mu_t\nabla(U+W*\mu_t)\bigr).
\end{equation}
Along a smooth solution,
\begin{equation}\label{eq:energy-dissipation}
 \frac{\dd}{\dd t}\mathcal F(\mu_t)
 =-\int_{\R^d}
 \left|\nabla\frac{\delta\mathcal F}{\delta\mu}[\mu_t](x)\right|^2
 \mu_t(\dd x)\leq0.
\end{equation}
\end{proposition}

\begin{proof}
The following calculation is formal unless the density, potential, and decay justify the differentiations and integration by parts; the metric theory of gradient flows supplies weak formulations under substantially milder hypotheses.\footnote{One convenient sufficient set of classical assumptions is
$U,W\in C^2(\R^d)$, $W$ even, with $\nabla U$ and $\nabla W$ of at most
linear growth, together with a positive solution
$\rho\in C^1([0,T];C^0(\R^d))\cap C^0([0,T];C^2(\R^d))$ having finite
second moment.  Assume in addition either compact spatial support or enough
decay that all displayed integrands are integrable, differentiation may pass
under the integral sign, and the boundary flux
$\rho_t\nabla(U+W*\rho_t)$ has vanishing contribution at infinity.  These
conditions are deliberately stronger than necessary.  Weak metric solutions
under lower semicontinuity, coercivity, and slope assumptions are treated in
\cite[Chapters~10--11]{AmbrosioGigliSavare2008} and
\cite[Chapters~23--24]{Villani2009}.}

Since $W$ is even, $\nabla W$ is odd.  Differentiating
\eqref{eq:finite-particle-energy} with respect to $X^i$, the two appearances
of $X^i$ in the double sum contribute equally, and hence
\[
 \nabla_{X^i}\mathcal F_N
 =\frac1N\nabla U(X^i)
 +\frac1{N^2}\sum_{j=1}^N\nabla W(X^i-X^j).
\]
This proves \eqref{eq:particle-gradient-flow}.  The ordinary Euclidean chain
rule then gives
\[
 \frac{\dd}{\dd t}\mathcal F_N(X_t)
 =\sum_i\nabla_{X^i}\mathcal F_N(X_t)\cdot\dot X_t^i
 =-\frac1N\sum_i\abs{\dot X_t^i}^2.
\]

To calculate the first variation, let $\mu_\varepsilon=\mu+\varepsilon\eta$ for a finite signed measure $\eta$ with $\eta(\R^d)=0$ and sufficient moments, such that $\mu_\varepsilon$ remains nonnegative for the one-sided or two-sided variations under consideration.  The zero-mass condition is exactly the tangent constraint imposed by $\mu_\varepsilon(\R^d)=1$.  Expanding the two terms of the energy and using
the symmetry of $W$ gives
\begin{align*}
\mathcal F(\mu_\varepsilon)-\mathcal F(\mu)
 &=\varepsilon\int_{\R^d}(U+W*\mu)\,\eta(\dd y)  +\frac{\varepsilon^2}{2}
   \iint_{\R^d\times\R^d}W(x-y)\,\eta(\dd x)\eta(\dd y).
\end{align*}
Hence,
$\langle \frac{\delta \mathcal F}{\delta\mu}(\mu), \eta\rangle
 =\lim_{\varepsilon\to 0}\frac{\mathcal F(\mu_\varepsilon)-\mathcal F(\mu)}{\varepsilon} = \int(U+W*\mu)\,\dd\eta$, which yields
\eqref{eq:interaction-first-variation}, up to an additive constant that is
invisible to zero-mass variations.

For the subsequent classical PDE
calculation we impose the density and regularity assumptions stated above. The velocity in the continuity equation is
\[
 v_t=-\nabla\frac{\delta\mathcal F}{\delta\mu}[\mu_t]
     =-\nabla(U+W*\mu_t),
\]
and \eqref{eq:gf-continuity} becomes
\eqref{eq:wasserstein-interaction-gf}.  Finally, using the continuity equation
and integrating by parts, we obtain the energy dissipation law \eqref{eq:energy-dissipation}.
% \begin{align*}
%  \frac{\dd}{\dd t}\mathcal F(\mu_t)
%  &=\int\frac{\delta\mathcal F}{\delta\mu}[\mu_t]
%        \partial_t\mu_t\dd x =\int\frac{\delta\mathcal F}{\delta\mu}[\mu_t]
%        \diver\left(\mu_t\nabla
%        \frac{\delta\mathcal F}{\delta\mu}[\mu_t]\right)\dd x =-\int\left|\nabla
%        \frac{\delta\mathcal F}{\delta\mu}[\mu_t]\right|^2\dd\mu_t.
% \end{align*}
\end{proof}

The equality
$\mathcal F_N(X^1,\ldots,X^N)=\mathcal F(\mu^N)$ is exact under the convention
that the $i=j$ terms are included.  Thus the finite and continuum energy laws
are not merely analogous: they are the same energy restricted to empirical
measures.  The factor $N$ in \eqref{eq:particle-gradient-flow} compensates for
the energy-per-particle normalization, and the normalized Euclidean speed
$N^{-1}\sum_i|\dot X^i|^2$ is exactly
$\int|v|^2\dd\mu^N$.  This is the clearest point of contact between IPS and
Wasserstein geometry.

The logical direction is important.  A potential IPS with the symmetry
above produces a Wasserstein gradient flow.  A general interaction
$b(x,y)$ need not arise from a symmetric potential, but the continuity equation
and its characteristics remain valid even when no scalar energy exists.

\subsection{Diffusion as entropy descent}

Brownian noise adds a second fundamental energy.  For
$\mu\in\cP_2(\R^d)$, define the Boltzmann entropy by
\begin{equation}\label{eq:boltzmann-entropy}
 \operatorname{Ent}(\mu)
 =\begin{cases}
   \displaystyle\int_{\R^d}\rho(x)\log\rho(x)\dd x,
      &\mu(\dd x)=\rho(x)\dd x,\\
   +\infty,&\text{otherwise}.
  \end{cases}
\end{equation}
The integral is understood in the extended sense; the finite second moment
prevents its negative part from being infinite.  Its first variation at a
positive density is $1+\log\rho$.  Indeed, for a bounded variation
$\eta$ with $\int\eta\dd x=0$ and $\rho+\varepsilon\eta>0$, the pointwise
identity $(s\log s)'=1+\log s$ and differentiation under the integral give
\[
 \left.\frac{\dd}{\dd\varepsilon}
 \operatorname{Ent}(\rho+\varepsilon\eta)\right|_{\varepsilon=0}
 =\int_{\R^d}(1+\log\rho)\eta\dd x.
\]
Therefore
\[
 \diver\bigl(\rho\nabla \frac{\delta \operatorname{Ent}}{\delta\rho}\bigr)= \diver\bigl(\rho\nabla(1+\log\rho)\bigr)=\Delta\rho.
\]
Consequently, the heat equation is the Wasserstein gradient flow of entropy.

For $\beta=\sigma^2/2$, the potential McKean--Vlasov equation from Section~\ref{sec:ips} is formally the gradient flow of the free energy
\begin{equation}\label{eq:free-energy}
 \mathcal F_\beta(\rho)
 =\mathcal F(\rho)+\beta\operatorname{Ent}(\rho):
\end{equation}
\begin{equation}\label{eq:free-energy-gradient-flow}
 \partial_t\rho
 =\diver\bigl(\rho\nabla(U+W*\rho)\bigr)+\beta\Delta\rho
 =\diver\left(\rho\nabla
       \frac{\delta\mathcal F_\beta}{\delta\rho}\right).
\end{equation}

Along a smooth positive solution,
\begin{equation}\label{eq:free-energy-dissipation}
 \frac{\dd}{\dd t}\mathcal F_\beta(\rho_t)
 =-\int_{\R^d}
 \abs{\nabla\bigl(U+W*\rho_t+\beta\log\rho_t\bigr)}^2
 \rho_t\dd x\leq0.
\end{equation}
At the particle level, diffusion is produced by random Brownian increments;
at the law level, the same effect is deterministic steepest descent of
entropy.  This equivalence is one of the most striking links among
probability, PDEs, and geometry: microscopic randomness becomes macroscopic
dissipation.

\subsection{The JKO scheme: implicit Euler in Wasserstein space}

The variational structure suggests a way to construct the flow without
first writing its PDE.  Given a time step $h>0$ and a measure $\mu^k$, define
\begin{equation}\label{eq:jko-scheme}
 \mu^{k+1}\in\argmin_{\mu\in\cP_2(\R^d)}
 \left\{\mathcal F_\beta(\mu)
       +\frac1{2h}W_2^2(\mu,\mu^k)\right\}.
\end{equation}
This is the Jordan--Kinderlehrer--Otto, or JKO, minimizing-movement scheme
\cite{JordanKinderlehrerOtto1998}; see also
\cite{AmbrosioGigliSavare2008}.  It is the metric analogue of the implicit
Euler step
\[
 x^{k+1}\in\argmin_x
 \left\{F(x)+\frac1{2h}|x-x^k|^2\right\}
\]
for the Euclidean gradient flow $\dot x=-\nabla F(x)$.  The first term in
\eqref{eq:jko-scheme} rewards energy decrease, while the Wasserstein term
penalizes moving too much mass in one time step.

Testing the minimization problem with the admissible competitor
$\mu=\mu^k$ immediately gives the discrete energy inequality
\begin{equation}\label{eq:jko-energy-inequality}
 \mathcal F_\beta(\mu^{k+1})
 +\frac1{2h}W_2^2(\mu^{k+1},\mu^k)
 \leq\mathcal F_\beta(\mu^k).
\end{equation}
Summing over $k$ controls the total discrete action
$\sum_k W_2^2(\mu^{k+1},\mu^k)/h$.

For the free energy $\mathcal F_\beta$ in \eqref{eq:free-energy} with $\mathcal F$  in \eqref{eq:interaction-energy}, a transparent sufficient
convex regime is: $U$ is lower semicontinuous, $\lambda$-convex, and
quadratically confining; $W$ is even, lower semicontinuous, convex, and
bounded below; and $\beta\geq0$.  The confinement supplies compact sublevels,
while the potential energy, interaction energy, and entropy have the required
displacement-convexity properties.  For nonconvex interactions the discrete scheme can still have
subsequential generalized minimizing-movement limits, but uniqueness and
convergence of the whole family require additional structure and cannot be
inferred from lower semicontinuity and compactness alone.  The original
construction is due to Jordan, Kinderlehrer, and Otto
\cite{JordanKinderlehrerOtto1998}; the metric theory, including weaker
compactness and slope hypotheses, is developed in
\cite[Chapters~2, 4, and 11]{AmbrosioGigliSavare2008}.

\begin{remark}[The bridge to attention]
Attention also defines a measure-dependent velocity field, and hence a
transport equation in continuous depth.  The preceding geometry applies
only if that field has the stronger variational form
\[
 \Phi[\mu](x)=-\nabla_x
 \frac{\delta\mathcal E}{\delta\mu}[\mu](x)
\]
for some energy $\mathcal E$.  Symmetric pair interactions have this form;
generic row-normalized softmax attention need not.
Section~\ref{sec:continuous-depth} will use this
criterion constructively: simple attention choices reveal a potential for a
frozen measure, while also showing why self-consistency and row normalization
obstruct an automatic global Wasserstein energy.
\end{remark}

\paragraph{\textbf{Literature and storyline.}}
Optimal transport is an unusually clear example of an applied question
creating new mathematical geometry.  Kantorovich's work grew from a 1938
production-allocation problem at a plywood factory; the method led both to
linear programming and to his formulation of mass transportation
\cite{KantorovichMacTutor,Kantorovich1942,NobelKantorovich1975}.  For quadratic
cost, Brenier later
showed that, under the appropriate absolute-continuity hypothesis, the optimal
map is the gradient of a convex function \cite{Brenier1991}.  Benamou and
Brenier replaced the static coupling by a least-action path of densities and
velocities \cite{BenamouBrenier2000}.  The next conceptual step was Jordan,
Kinderlehrer, and Otto's observation that the Fokker--Planck equation can be
constructed by repeated Wasserstein proximal minimization
\cite{JordanKinderlehrerOtto1998}.  Thus a factory-allocation problem became a
metric calculus for nonlinear PDEs.  This progression also explains the
order of the mathematics here: the continuity equation supplies tangent
velocities, Benamou--Brenier supplies their metric cost, first variations
supply forces, and the JKO scheme turns formal dissipation into a construction.
For attention, this history gives a precise test rather than an analogy: one
must exhibit an energy whose Wasserstein gradient is the attention field.


%%%%%%%%%%%%%% ============
\section{Self-attention as an interacting particle system}
\label{sec:self-attention}
\paragraph{\textbf{Key question}}
Which parts of self-attention are genuinely particle-like, and which
architectural features fall outside the analogy?

\subsection{One head: an exact empirical-measure formula}
For tokens $x_1,\ldots,x_N\in\R^d$, let
\[
  q_i=Qx_i,\qquad k_j=Kx_j,\qquad v_j=Vx_j,
\]
where $Q,K:\R^d\to\R^{d_k}$ and $V:\R^d\to\R^{d_v}$, i.e, $Q,K\in \R^{d_k\times d}$ and $V\in \R^{d_v\times d}$.  To add the attention
output to the residual stream, assume for now that $d_v=d$; otherwise an output
map is inserted.  With temperature $\tau>0$ and $X=(x_1,\ldots,x_N)\in \R^{d\times N}$, define
\begin{equation}\label{eq:attention-weights}
  s_\theta(x,y)=\frac{\ip{Qx}{Ky}}{\tau},\qquad
  a_{ij}(X)=\frac{e^{s_\theta(x_i,x_j)}}
  {\sum_{r=1}^Ne^{s_\theta(x_i,x_r)}}.
\end{equation}
The commonly used temperature is $\tau=\sqrt{d_k}$.  A residual
attention layer with step size $h_\ell$ is
\begin{equation}\label{eq:attention-layer}
    x_i^{\ell+1}=x_i^\ell+
    h_\ell\sum_{j=1}^Na_{ij}(X^\ell)Vx_j^\ell.
\end{equation}

For any $\mu\in\cP(\R^d)$ for which the integrals are finite, define
\begin{equation}\label{eq:attention-field}
 \begin{split}
    Z_\theta[\mu](x)
      =\int e^{s_\theta(x,y)}\,\mu(\dd y),\quad 
    \Phi_\theta[\mu](x)
      =\frac{\displaystyle
        \int e^{s_\theta(x,y)}Vy\,\mu(\dd y)}
       {Z_\theta[\mu](x)}.
 \end{split}
\end{equation}
In particular, for the empirical measure $\mu^{N}=\frac1N\sum_i\delta_{x_i}$, since the factors $1/N$ in numerator and denominator cancel, we have 
\begin{equation}\label{eq:attention-exact}
  Z_\theta[\mu^N](x) = \sum_{j=1}^N e^{s_\theta(x,x_j)},\quad  \Phi_\theta[\mu^{N}](x_i)
  =\sum_{j=1}^N a_{ij}(X)Vx_j.
\end{equation}
Thus, the IPS representation of the idealized attention layer is an
\emph{identity}: let 
\[
  T_{\theta,\mu}^{h}(x)=x+h\Phi_\theta[\mu](x),
\]
then the corresponding empirical-measure update is exactly
\begin{equation}\label{eq:attention-pushforward}
    \mu^{N,\ell+1}
    =\bigl(T_{\theta^\ell,\mu^{N,\ell}}^{h_\ell}\bigr)_\#
      \mu^{N,\ell}.
\end{equation}
Because the map itself depends on the input measure, this is a nonlinear
pushforward operator.

In matrix form, the attention layer is 
\begin{equation}\label{eq:attention-matrix}
 X^{\ell+1}=X^\ell+h_\ell\operatorname{Att}(X^\ell), \quad  \operatorname{Att}(X)=VXA^\top,\qquad
  A=(a_{ij}(X))_{i,j=1}^N\in \R^{N\times N}.
\end{equation}

\begin{remark}[Row stochasticity is not a conservation law]
The matrix $A=(a_{ij})$ is row stochastic:
$a_{ij}\geq0$ and $\sum_j a_{ij}=1$.  It is generally neither symmetric nor
column stochastic.  For the consensus-style dynamics
\begin{equation}\label{eq:attn-consensus-form}
  \dot x_i=\sum_j a_{ij}(X)(x_j-x_i),
\end{equation}
the mean satisfies
\[
  \frac{\dd}{\dd t}\frac1N\sum_ix_i
  =\frac1N\sum_j\left(\sum_i a_{ij}-1\right)x_j.
\]
Hence the mean is conserved if $A$ is doubly stochastic, but not merely because
it is row stochastic.  Standard attention in
\eqref{eq:attention-layer} does not even have the difference form
$x_j-x_i$.  It is therefore not automatically a physical pair force and need
not conserve momentum, energy, or any token average.  Such structure must be
designed.
\end{remark}

\begin{remark}[Temperature is a mathematical conditioning parameter]
As $\tau\downarrow0$, attention approaches hard selection when the top score is
separated, but the global Lipschitz estimate becomes poor near score ties.  As
$\tau\uparrow\infty$, weights become nearly uniform and the interaction loses
query specificity.  Temperature therefore trades selectivity against
regularity; it is not merely a cosmetic scaling.
\end{remark}

The above idealized attention layer is permutation equivariant: if the input tokens are relabeled, the output tokens are relabeled in the same way.  This is a direct consequence of the empirical-measure representation \eqref{eq:attention-exact}.  The following proposition makes this statement precise. 
\begin{proposition}[Permutation equivariance]\label{prop:perm-equiv}
Without positional encodings or masks, the token map on column-token matrices
$\operatorname{Att}:\R^{d\times N}\to\R^{d\times N}$ commutes with every
relabeling:
\[
  \operatorname{Att}(XP)=\operatorname{Att}(X)P
\]
for each permutation matrix $P\in\R^{N\times N}$.
\end{proposition}
\begin{proof}
Let $S(X)_{ij}=s_\theta(x_i,x_j)$,  and
$A(X)=\operatorname{softmax}_{\rm row}(S(X)) $. That is,
$$S(X)= \frac{1}{\tau} X^\top Q^\top K X, \quad A(X)_{ij}= \frac{\exp(S(X)_{ij})}{\sum_l \exp(S(X)_{il})}.$$
Relabeling the columns of
$X$ by $P$ relabels both queries and keys, so
\[
 S(XP)=P^\top S(X)P,
 \qquad
 A(XP)=P^\top A(X)P.
\]
The second identity holds because conjugation by $P$ reorders the rows and,
within each row, the entries to which softmax is applied.  Since the columns
of the value matrix are relabeled as $VXP$, formula
\eqref{eq:attention-matrix} gives
\[
\begin{split}
 \operatorname{Att}(XP)
 &=VXP\,A(XP)^\top
  =VXP\bigl(P^\top A(X)P\bigr)^\top\\
 &=VXP P^\top A(X)^\top P
  =VXA(X)^\top P
  =\operatorname{Att}(X)P.
\end{split}
\]
Equivalently, \eqref{eq:attention-field} depends only on the current query
and the empirical measure, which is unchanged by a relabeling.
\end{proof}

Permutation equivariance does not mean that all tokens receive the same
output.  The query $x_i$ distinguishes the particles.  It means only that
renaming the tokens renames the outputs in the same way.



\subsection{Regularity of the attention field}

The next lemma shows that the attention field is locally Lipschitz in both the query and the empirical measure.  While the the softmax normalization is smooth, the Lipschitz constants can be large if the query and key vectors are large or if the temperature is small.  This is a mathematical manifestation of the fact that attention can be highly selective: if one key has a much higher score than all others, then the output is nearly constant in the query and insensitive to all other keys.  In this case, the attention field is nearly flat, so its derivative is small, but the Lipschitz constant is large because a small change in the query or in the empirical measure can change which key has the top score.  The lemma below quantifies this effect.

Recall that $W_1$ is the $1$-Wasserstein distance defined in \eqref{eq:w1-definition}.  
% \[
%  W_1(\mu,\nu)
%  =\inf_{\pi\in\Gamma(\mu,\nu)}
%    \int_{\R^d\times\R^d}\abs{y-y'}\,\pi(\dd y,\dd y'),
% \]
% where $\Gamma(\mu,\nu)$ is the set of probability measures on
% $\R^d\times\R^d$ with first marginal $\mu$ and second marginal $\nu$.

\begin{lemma}[Local Lipschitz continuity]\label{lem:attention-lipschitz}
Let $\mu,\nu$ be supported in the ball $B_R\subset\R^d$.  For fixed
$Q,K,V$ and $\tau>0$, there are finite constants $L_x,L_\mu$ such that, for
$x,x'\in B_R$,
\begin{equation}\label{eq:attn-lip}
 \abs{\Phi_\theta[\mu](x)-\Phi_\theta[\nu](x')}
 \leq L_x\abs{x-x'}+L_\mu W_1(\mu,\nu), 
\end{equation}
and the constants $L_x,L_\mu$ depend only on $R,\tau$ and the operator norms of $Q,K,V$, and may grow
exponentially in $\norm{Q}\norm{K}R^2/\tau$.
\end{lemma}
\begin{proof}
Set
\[
 \kappa=\frac{\norm{Q}\norm{K}}{\tau},
 \qquad S=\kappa R^2.
\]
For $x,y,y'\in B_R$, Cauchy--Schwarz gives
\[
 \abs{s_\theta(x,y)}\leq S,
 \qquad
 \abs{s_\theta(x,y)-s_\theta(x,y')}
 \leq \kappa R\abs{y-y'}.
\]
Since the derivative of the exponential is bounded by $e^S$ on
$[-S,S]$, the functions
\[
 f_x(y)=e^{s_\theta(x,y)},
 \qquad g_x(y)=e^{s_\theta(x,y)}Vy
\]
satisfy
\begin{align}
 \norm{f_x}_\infty&\leq e^S,
 &\Lip(f_x)&\leq e^S\kappa R,\notag\\
 \norm{g_x}_\infty&\leq e^S\norm{V}R,
 &\Lip(g_x)&\leq e^S\norm{V}(1+\kappa R^2).
 \label{eq:softmax-integrand-bounds}
\end{align}
Write
\[
 N_\mu(x)=\int g_x(y)\,\mu(\dd y),
 \qquad Z_\mu(x)=\int f_x(y)\,\mu(\dd y).
\]
In particular, $Z_\mu(x)\geq e^{-S}$ and
$\abs{N_\mu(x)}\leq e^S\norm{V}R$.  The coupling bound
\[
 \left|\int f\,\dd\mu-\int f\,\dd\nu\right|
 \leq\Lip(f)W_1(\mu,\nu)
\]
and \eqref{eq:softmax-integrand-bounds} imply
\begin{align*}
 \abs{N_\mu(x)-N_\nu(x)}
 &\leq e^S\norm{V}(1+\kappa R^2)W_1(\mu,\nu),\\
 \abs{Z_\mu(x)-Z_\nu(x)}
 &\leq e^S\kappa R W_1(\mu,\nu).
\end{align*}
Using
$a/b-c/d=(a-c)/b+c(d-b)/(bd)$ therefore gives
\begin{align}
 \abs{\Phi_\theta[\mu](x)-\Phi_\theta[\nu](x)}
 &= \abs{\frac{N_\mu(x)}{Z_\mu(x)}-\frac{N_\nu(x)}{Z_\nu(x)}}
 \leq
 \frac{\abs{N_\mu(x)-N_\nu(x)}}{Z_\mu(x)}
 +\frac{\abs{N_\nu(x)}
          \abs{Z_\nu(x)-Z_\mu(x)}}{Z_\mu(x)Z_\nu(x)} \notag \\
 &\leq \norm{V}\left[
 e^{2S}(1+\kappa R^2)+e^{4S}\kappa R^2\right]
 W_1(\mu,\nu) \notag \\
 & = L_\mu W_1(\mu,\nu), \quad L_\mu=\norm{V}\left[
 e^{2S}(1+\kappa R^2)+e^{4S}\kappa R^2\right].  \label{eq:attention-measure-lip}
\end{align}


It remains to vary the query.  For $x,x',y\in B_R$,
\[
 \abs{s_\theta(x,y)-s_\theta(x',y)}
 \leq\kappa R\abs{x-x'}.
\]
The same mean-value and quotient calculations, now with the measure fixed,
give 
\begin{equation}\label{eq:attention-state-lip}
 \abs{\Phi_\theta[\nu](x)-\Phi_\theta[\nu](x')}
 \leq L_x\abs{x-x'},
 \qquad
 L_x=\norm{V}\kappa R^2(e^{2S}+e^{4S}).
\end{equation}
Finally, insert and subtract $\Phi_\theta[\nu](x)$ and combine
\eqref{eq:attention-measure-lip}--\eqref{eq:attention-state-lip} to obtain
\eqref{eq:attn-lip}.
\end{proof}

%%%%%%%%%%% -------------------
\subsection{Marked particles, positions, and causal attention}
Sequence models are not exchangeable after position is discarded.  A clean
repair is to augment each particle by a mark:
\[
  z_i=(x_i,p_i,c_i)\in
  \R^d\times\mathcal P_{\rm pos}\times\mathcal P_{\rm type},\qquad
  \nu^N=\frac1N\sum_i\delta_{z_i}.
\]
Here $\mathcal P_{\rm pos}$ and $\mathcal P_{\rm type}$ are measurable
mark spaces, not probability measures.  For a sequence of maximum length
$L$, one may take the ordered discrete space
$\mathcal P_{\rm pos}=\{1,\ldots,L\}$; for continuously indexed data it may
instead be an ordered subset of $\R$.  The type space can be a finite set of
token roles or modalities, or a Euclidean metadata space.  Thus neither space
must be a subset of $\R$, although an order on $\mathcal P_{\rm pos}$ is
needed for the causal mask below.

Scores and values may depend on both states and marks.
Precisely, $s_\theta$ is a learned measurable score map
\[
 s_\theta:
 \bigl(\R^d\times\mathcal P_{\rm pos}\times\mathcal P_{\rm type}\bigr)^2
 \longrightarrow\R,
\]
and
$V_\theta:\R^d\times\mathcal P_{\rm pos}\times\mathcal P_{\rm type}
\to\R^d$ is a learned value map.  For example, choose embeddings
$e_{\rm pos}$ and $e_{\rm type}$, set
\[
 u(x,p,c)=\bigl(x,e_{\rm pos}(p),e_{\rm type}(c)\bigr)\in\R^D,
\]
and define
\[
 s_\theta(x,p,c;y,p',c')
 =\frac{\ip{Qu(x,p,c)}{Ku(y,p',c')}}{\tau},
 \qquad
 V_\theta(y,p',c')=OVu(y,p',c'),
\]
with matrices of compatible dimensions.  A mask
$m(p,p')\in\{0,1\}$ then gives
\begin{equation}\label{eq:masked-field}
  \Phi_\theta[\nu](x,p,c)
  =\frac{\int m(p,p')e^{s_\theta(x,p,c;y,p',c')}
                   V_\theta(y,p',c')
     \,\nu(\dd y,\dd p',\dd c')}
    {\int m(p,p')e^{s_\theta(x,p,c;y,p',c')}
     \,\nu(\dd y,\dd p',\dd c')},
\end{equation}
provided at least one visible key has positive mass.  The causal mask is
$m(p,p')=\one_{\{p'\leq p\}}$.  The marked cloud is still invariant under
reordering the list of \emph{state--mark pairs}; it is not invariant under
changing which state carries which position.  This distinction is essential:
unmarked IPS is best matched to sets and exchangeable in-context examples,
while marked IPS is needed for language and time series.

\begin{remark}[Equivariance and positions]
Let $\widetilde x_i=x_i+p_i$ for fixed absolute position vectors.  Permuting
the $x_i$ while keeping the $p_i$ fixed need not commute with attention,
whereas permuting the pairs $(x_i,p_i)$ does.  More precisely, retain the
column-token convention $X,p\in\R^{d\times N}$ and define
$\mathcal A(X,p)=\operatorname{Att}(X+p)$.  For a permutation matrix
$P\in\R^{N\times N}$, content-only permutation generally fails:
\[
 \mathcal A(XP,p)\ne\mathcal A(X,p)P,
\]
whereas simultaneous permutation of the state--position pairs satisfies
\[
 \mathcal A(XP,pP)
 =\operatorname{Att}\bigl((X+p)P\bigr)
 =\operatorname{Att}(X+p)P
 =\mathcal A(X,p)P.
\]
\end{remark}

%%%%%%%%%%%%%% -----------------------
\subsection{Multiple heads and the rest of a Transformer block}
For $H$ heads with parameters
$\theta_m=(Q_m,K_m,V_m,O_m,\tau_m)$, the field is
\begin{equation}\label{eq:multihead-field}
  \Phi_{\boldsymbol\theta}[\mu](x)
  =\sum_{m=1}^H O_m
   \frac{\int e^{\ip{Q_mx}{K_my}/\tau_m}V_my\,\mu(\dd y)}
        {\int e^{\ip{Q_mx}{K_my}/\tau_m}\,\mu(\dd y)}.
\end{equation}
Here $m$ indexes a head and $\tau_m$ is its temperature.  Multi-head attention
is therefore a multi-kernel interaction, not a random choice of one head.
Heads can probe distinct geometries, but nothing in the definition prevents
them from becoming redundant.

A token-wise MLP can be included as a local field $f_\ell(x)$:
\[
 x_i^{\ell+1}=x_i^\ell+h_\ell
 \bigl(f_\ell(x_i^\ell)+
       \Phi_{\boldsymbol\theta_\ell}[\mu^{N,\ell}](x_i^\ell)\bigr).
\]
Examples of such local fields include the layer normalization, dropout, and gating operations that are standard in Transformer blocks. Layer normalization is token-wise in its usual form, but it is a nonlinear projection/rescaling rather than a small vector-field increment.
For $x\in\R^d$, its usual form is
\[
 \operatorname{LN}_{a,b,\varepsilon}(x)
 =a\odot
 \frac{x-\bar x\mathbf 1}
 {\sqrt{d^{-1}\norm{x-\bar x\mathbf 1}^2+\varepsilon}}+b,
 \qquad
 \bar x=\frac1d\sum_{r=1}^d x_r,
\]
where $a,b\in\R^d$ are learned, $\varepsilon>0$, and $\odot$ denotes
componentwise multiplication.  A gate is a learned scalar or vector
$g(x)\in[0,1]^d$ that modulates an update, for example
$x^+=x+g(x)\odot F(x)$.  Dropout is a training-time random map
$D_\rho(u)=\xi\odot u/(1-\rho)$, where the coordinates of $\xi$ are
Bernoulli$(1-\rho)$; it therefore introduces randomness not present in a
deterministic vector field.  Residual placement specifies where the skip
connection is added.  Normalization order distinguishes, for example,
\[
 \text{pre-norm: }x^+=x+F(\operatorname{LN}(x)),
 \qquad
 \text{post-norm: }x^+=\operatorname{LN}(x+F(x)).
\]
These choices define different update maps and can therefore materially
change stability, invariants, and continuous-depth limits.
The single-field model should therefore be read as a controlled
idealization of one sublayer, not as an exact description of every operation in
a modern block

\paragraph{\textbf{Literature and storyline.}}
Attention entered neural sequence modeling as a learned alignment mechanism.
Bahdanau, Cho, and Bengio let a decoder form an input-dependent weighted
combination of encoder states instead of compressing an entire sentence into
one fixed vector \cite{BahdanauEtAl2015}.  Vaswani and collaborators then made
self-attention the principal interaction in the Transformer, removing the
recurrent and convolutional backbone \cite{VaswaniEtAl2017}.  Uszkoreit's
contemporaneous Google Research account emphasizes two design pressures:
short paths between distant positions and computation that parallelizes over
tokens \cite{UszkoreitGoogle2017}.  Those engineering choices lead directly to
the mathematical object in \eqref{eq:attention-field}: every token interacts
with the empirical measure of all tokens through a normalized, learned
kernel.  The IPS formulation is therefore exact for the attention sublayer,
but not for the whole Transformer block.  Positional encodings, masks,
residual connections, normalization, and token-wise nonlinearities must be
added explicitly.  This distinction is historically and mathematically
important: the useful question is not whether a Transformer ``is'' a particle
system, but which architectural operations admit a particle or mean-field
description and what that description allows us to prove.

\subsection{Connections forward}
The exact nonlinear map \eqref{eq:attention-pushforward} is the shared object of
the rest of Part~II.  Iterating it raises a numerical stability question;
taking many small steps suggests a transport PDE; studying its geometry leads
to consensus and collapse; and treating $\theta$ or the composite interaction
as unknown leads to an inverse problem.  In Part~IV, a context-dependent
attention field will be interpreted as an algorithm that infers a latent task.

\section{Continuous depth and kinetic limits}
\label{sec:continuous-depth}
\paragraph{\textbf{Key question}}
What is gained, and what must be assumed, when layer number is replaced by time?

\subsection{Residual layers as an Euler scheme}
Set $h=T/L$ and $t_\ell=\ell h$.  A residual stack has the form
\begin{equation}\label{eq:residual-euler}
 x_i^{\ell+1}
 =x_i^\ell+h\Phi_{\theta^\ell}[\mu^{N,\ell}](x_i^\ell).
\end{equation}
Suppose that the layer parameters sample a curve:
$\theta^\ell\approx\theta_{t_\ell}$.  The proposed continuous-depth particle
system is
\begin{equation}\label{eq:token-ode}
    \dot X_t^i=\Phi_{\theta_t}[\mu_t^N](X_t^i),
    \qquad \mu_t^N=\frac1N\sum_i\delta_{X_t^i}.
\end{equation}
This interpretation is exactly the explicit Euler method applied to an ODE on
$(\R^d)^N$.  It is useful because it imports consistency, stability, and step
size questions from numerical analysis.

\begin{assumption}[A controlled depth interpolation]\label{ass:depth}
On a domain $D\subset\R^d$ containing both the exact trajectories and the
explicit-Euler iterates considered below, assume uniformly in $t\in[0,T]$ that
\begin{align}
 \abs{\Phi_{\theta_t}[\mu](x)-\Phi_{\theta_t}[\nu](y)}
 &\leq L_x\abs{x-y}+L_\mu W_1(\mu,\nu),
 \label{eq:depth-lipschitz}\\
 \abs{\Phi_{\theta_t}[\mu](x)-\Phi_{\theta_s}[\mu](x)}
 &\leq L_t\abs{t-s},\qquad
 \abs{\Phi_{\theta_t}[\mu](x)}\leq B.
 \label{eq:depth-time}
\end{align}
The measures in these inequalities are supported in $D$.
\end{assumption}
Lemma~\ref{lem:attention-lipschitz} verifies the first condition for bounded
token states and fixed bounded matrices.  The time condition is substantive:
an arbitrary list of untied layer weights need not sample a regular curve.

\begin{proposition}[Euler error, uniform in particle number]\label{prop:euler}
Under Assumption~\ref{ass:depth}, let $X(t)$ solve \eqref{eq:token-ode} and let
$x^\ell$ solve \eqref{eq:residual-euler} from the same initial cloud, with
$\theta^\ell=\theta_{t_\ell}$.  Then
\begin{equation}\label{eq:euler-error}
 \max_{0\leq\ell\leq L}
 \left(\frac1N\sum_{i=1}^N
 \abs{x_i^\ell-X_{t_\ell}^i}^2\right)^{1/2}
 \leq C_T h,
\end{equation}
where $C_T$ depends on $B,L_x,L_\mu,L_t,T$ but not on $N$.
\end{proposition}
\begin{proof}
For $z=(z_1,\ldots,z_N)\in(\R^d)^N$, set
\[
 \norm{z}_N=\left(\frac1N\sum_{i=1}^N\abs{z_i}^2\right)^{1/2},
 \qquad
 F_t^N(z)_i=\Phi_{\theta_t}[\mu_z^N](z_i),
 \quad
 \mu_z^N=\frac1N\sum_{i=1}^N\delta_{z_i}.
\]
Matching particles with the same labels gives
\begin{equation}\label{eq:matched-empirical-w1}
 W_1(\mu_z^N,\mu_w^N)
 \leq\frac1N\sum_{i=1}^N\abs{z_i-w_i}
 \leq\norm{z-w}_N,
\end{equation}
where the last inequality is Cauchy--Schwarz.  Minkowski's inequality,
\eqref{eq:depth-lipschitz}, and \eqref{eq:matched-empirical-w1} now yield
\begin{equation}\label{eq:particle-field-lipschitz}
 \norm{F_t^N(z)-F_t^N(w)}_N
 \leq L_x\norm{z-w}_N+L_\mu W_1(\mu_z^N,\mu_w^N)
 \leq L\norm{z-w}_N,
 \qquad L:=L_x+L_\mu.
\end{equation}
The important point is that $L$ is independent of $N$.  Similarly,
\begin{equation}\label{eq:particle-field-time}
 \norm{F_t^N(z)-F_s^N(z)}_N\leq L_t\abs{t-s},
 \qquad \norm{F_t^N(z)}_N\leq B.
\end{equation}

Let $X_t=(X_t^1,\ldots,X_t^N)$.  The last bound in
\eqref{eq:particle-field-time} gives
\begin{equation}\label{eq:token-time-increment}
 \norm{X_s-X_{t_\ell}}_N
 \leq\int_{t_\ell}^s\norm{F_r^N(X_r)}_N\dd r
 \leq B(s-t_\ell),
 \qquad t_\ell\leq s\leq t_{\ell+1}.
\end{equation}
Define the one-step defect of the exact solution by
\[
 \rho^{\ell+1}
 :=X_{t_{\ell+1}}-X_{t_\ell}-hF_{t_\ell}^N(X_{t_\ell}).
\]
Using \eqref{eq:particle-field-lipschitz}--\eqref{eq:token-time-increment},
\begin{align*}
 \norm{\rho^{\ell+1}}_N
 &\leq\int_{t_\ell}^{t_{\ell+1}}
 \norm{F_s^N(X_s)-F_{t_\ell}^N(X_{t_\ell})}_N\dd s\\
 &\leq\int_{t_\ell}^{t_{\ell+1}}
 \left(L\norm{X_s-X_{t_\ell}}_N+L_t(s-t_\ell)\right)\dd s\\
 &\leq C_0h^2,
 \qquad C_0:=\frac12(LB+L_t).
\end{align*}

Let $e_\ell=\norm{x^\ell-X_{t_\ell}}_N$.  The Euler update and the
definition of $\rho^{\ell+1}$ give the exact error identity
\[
 x^{\ell+1}-X_{t_{\ell+1}}
 =x^\ell-X_{t_\ell}
 +h\bigl(F_{t_\ell}^N(x^\ell)-F_{t_\ell}^N(X_{t_\ell})\bigr)
 -\rho^{\ell+1}.
\]
Taking the normalized product norm, then applying the triangle inequality,
\eqref{eq:particle-field-lipschitz}, and the one-step defect bound, yields
\[
\begin{split}
 e_{\ell+1}
 &\leq e_\ell
 +h\norm{F_{t_\ell}^N(x^\ell)-F_{t_\ell}^N(X_{t_\ell})}_N
 +\norm{\rho^{\ell+1}}_N\\
 &\leq e_\ell+hLe_\ell+C_0h^2
 =(1+hL)e_\ell+C_0h^2,
\end{split}
\]
Since the discrete and continuous systems have the same initial data,
$e_0=0$.
Iteration of this scalar recursion gives, when $L>0$,
\[
 e_\ell
 \leq C_0h^2\sum_{r=0}^{\ell-1}(1+hL)^r
 =\frac{C_0}{L}h\bigl((1+hL)^\ell-1\bigr)
 \leq\frac{C_0}{L}(e^{LT}-1)h.
\]
When $L=0$, the same recursion gives $e_\ell\leq C_0Th$.  This proves
\eqref{eq:euler-error} with a constant independent of $N$.
\end{proof}

\begin{remark}[What an ODE limit does not say]
The proposition justifies continuous depth along a family of networks whose
step size tends to zero while a regular parameter curve is held fixed.  It does
not say that a fixed 12- or 96-layer network with order-one, independently
trained blocks is close to an autonomous ODE.  At best such a model is a
finite, nonautonomous composition.  Finite-depth estimates are often more
relevant than infinite-time equilibria.
\end{remark}



\subsection{The measure equation for continuous-depth attention}

The general characteristic and continuity-equation framework has
already been developed in Section~\ref{sec:gradient-flow}.  For attention, the
only new ingredient is the particular self-consistent velocity
$v_t=\Phi_{\theta_t}[\mu_t]$.  Indeed, for a smooth test function $\varphi$,
the finite token ODE \eqref{eq:token-ode} gives the exact identity
\begin{equation}\label{eq:token-weak}
    \frac{\dd}{\dd t}\int\varphi\,\dd\mu_t^N
    =\int\nabla\varphi(x)\cdot
       \Phi_{\theta_t}[\mu_t^N](x)\,\mu_t^N(\dd x).
       % \text{ i.e., }   \partial_t\mu_t^N+ \diver\bigl(\mu_t^N \Phi_{\theta_t}[\mu_t^N]\bigr)=0.
 \end{equation}
That is, $\mu_t^N$ is a weak solution of the continuity equation
\begin{equation}\label{eq:attention-pde}
    \partial_t\mu_t+
    \diver\bigl(\mu_t\Phi_{\theta_t}[\mu_t]\bigr)=0,
    \qquad \mu_{t=0}=\mu_0.
\end{equation}
Under Assumption~\ref{ass:depth}, its solution is transported by the
self-consistent characteristics
$\dot F_t(x)=\Phi_{\theta_t}[\mu_t](F_t(x))$ and
$\mu_t=(F_t)_\#\mu_0$. If $\mu_t^N\to\mu_t$ strongly enough to pass through the nonlinear field, the measure $\mu_t$ is also a solution to the equation.


%%%%%%%%%%%%%%%%
\subsection{What the gradient-flow viewpoint says about attention}

Continuous depth and gradient flow are different assertions.
Equation~\eqref{eq:attention-pde} is a continuity equation for every
sufficiently regular attention field. In the standard Wasserstein geometry it
is a gradient descent of a functional $\mathcal F$ only when
\begin{equation}\label{eq:attention-variational-criterion}
 \Phi_\theta[\mu](x)
 =-\nabla_x\frac{\delta\mathcal F}{\delta\mu}[\mu](x).
\end{equation}
Note that the parameter $\theta$ has to be fixed in this criterion. \footnote{It is possible to have a gradient flow with time-varying parameters $\theta_t$, but the functional $\mathcal F$ will also depend on time.  In this case, the energy balance becomes
\[
\frac{\dd}{\dd t}\mathcal F(\mu_t,t)
 =\frac{\partial \mathcal F}{\partial t}(\mu_t,t)
 -\int_{\R^d}\left|
 \nabla\frac{\delta\mathcal F}{\delta\mu}[\mu_t](x)
 \right|^2\mu_t(\dd x).
\]
It is nonpositive only when the explicit time dependence is dominated by the
dissipation term.
}

We now test this criterion first for unnormalized identity attention (with $Q=K=V=I_d$), then for
its row-normalized version.  This order separates the interaction energy from
the state-dependent factor $Z_\mu^{-1}$ introduced by softmax
normalization.

\paragraph{\textbf{Unnormalized identity-attention field.}}
Fix $\tau>0$, set
\[
 k_\tau(x,y)=e^{\ip{x}{y}/\tau},
 \qquad
 Z_\mu(x)=\int_{\R^d}k_\tau(x,y)\,\mu(\dd y),
\]
and define
\begin{equation}\label{eq:unnormalized-identity-attention-field}
 \widetilde\Phi[\mu](x)
 :=\int_{\R^d}k_\tau(x,y)y\,\mu(\dd y).
\end{equation}
For the variational calculation, assume that $\mu$ is compactly supported;
the same formulas hold under sufficient exponential-integrability and
differentiability assumptions.  Introduce the similarity functional
\begin{equation}\label{eq:identity-attention-energy}
 \mathcal J_\tau(\mu)
 :=\frac{\tau}{2}\iint_{\R^d\times\R^d}
 k_\tau(x,y)\,\mu(\dd x)\mu(\dd y)
 =\frac{\tau}{2}\int_{\R^d}Z_\mu(x)\,\mu(\dd x)
\end{equation}
and the interaction energy
\begin{equation}\label{eq:attention-free-energy}
 \mathcal F_{\mathrm{att}}(\mu):=-\mathcal J_\tau(\mu).
\end{equation}
The minus sign is essential: attention moves toward large similarity, so the
positive similarity $\mathcal J_\tau$ increases while the energy
$\mathcal F_{\mathrm{att}}$ decreases.

Let $\eta$ be a signed measure with zero total mass.  Symmetry
$k_\tau(x,y)=k_\tau(y,x)$ gives
\[
 \left.\frac{\dd}{\dd\varepsilon}
 \mathcal J_\tau(\mu+\varepsilon\eta)\right|_{\varepsilon=0}
 =\tau\int_{\R^d}Z_\mu(x)\,\eta(\dd x).
\]
Consequently,
\begin{equation}\label{eq:identity-attention-first-variation}
 \frac{\delta\mathcal F_{\mathrm{att}}}{\delta\mu}[\mu](x)
 =-\tau Z_\mu(x),
 \qquad
 -\nabla_x\frac{\delta\mathcal F_{\mathrm{att}}}{\delta\mu}[\mu](x)
 =\widetilde\Phi[\mu](x).
\end{equation}
Thus
\begin{equation}\label{eq:unnormalized-attention-gradient-flow}
 \partial_t\mu_t+\diver\bigl(\mu_t\widetilde\Phi[\mu_t]\bigr)=0
 \quad\Longleftrightarrow\quad
 \partial_t\mu_t
 =\diver\left(\mu_t\nabla
 \frac{\delta\mathcal F_{\mathrm{att}}}{\delta\mu}[\mu_t]\right)
\end{equation}
is the standard $W_2$-gradient descent of
$\mathcal F_{\mathrm{att}}$.  Along a smooth solution,
\[
 \frac{\dd}{\dd t}\mathcal F_{\mathrm{att}}(\mu_t)
 =-\int_{\R^d}\left|
 \nabla\frac{\delta\mathcal F_{\mathrm{att}}}{\delta\mu}[\mu_t](x)
 \right|^2\mu_t(\dd x)\leq0.
\]
Equivalently,
\begin{equation}\label{eq:identity-attention-energy-ascent}
 \frac{\dd}{\dd t}\mathcal J_\tau(\mu_t)
 =\int_{\R^d}\left|\widetilde\Phi[\mu_t](x)\right|^2
 \mu_t(\dd x)\geq0.
\end{equation}

\paragraph{\textbf{Normalized identity-attention.}}
Row normalization gives
\begin{equation}\label{eq:normalized-identity-attention-field}
 \Phi[\mu](x)
 :=\frac{\widetilde\Phi[\mu](x)}{Z_\mu(x)}.
\end{equation}
Since $Z_\mu(x)>0$, \eqref{eq:identity-attention-first-variation} yields
\[
 \Phi[\mu](x)
 =-\frac1{Z_\mu(x)}
 \nabla_x\frac{\delta\mathcal F_{\mathrm{att}}}{\delta\mu}[\mu](x).
\]
The energy is unchanged, but softmax normalization multiplies the
gradient by the positive factor $Z_\mu^{-1}$.  

Define the positive operator $\mathsf K_\mu: C^\infty(\R^d) \to \mathcal E'_0(\R^d)$ by the weak form
\begin{equation}\label{eq:normalized-attention-onsager-operator}
 \left\langle\mathsf K_\mu\xi,\varphi\right\rangle
 :=\int_{\R^d}\nabla\xi(x)\cdot\nabla\varphi(x)\,\frac{1}{Z_\mu}\mu(\dd x),
 \qquad \varphi\in C^\infty(\R^d),
\end{equation}
where $\mathcal E'_0(\R^d)$ denotes the compactly supported distributions
with zero total mass. Equivalently, 
 \[
 \mathsf K_\mu\xi
 =-\diver\left(\frac{\mu}{Z_\mu}\nabla\xi\right)
\]
in the sense of distributions. The distributional range is necessary in
general: the divergence of a vector-valued singular measure need not be a
measure.  Testing \eqref{eq:normalized-attention-onsager-operator} with
$\varphi=1$ shows that the output has zero total mass, and testing with
$\varphi=\xi$ gives the positivity identity
$
 \left\langle\mathsf K_\mu\xi,\xi\right\rangle
 =\int_{\R^d}\frac{1}{Z_\mu(x)}\abs{\nabla\xi(x)}^2\,\mu(\dd x)
 \geq0.
 $


Then the normalized equation is the generalized gradient descent
\begin{equation}\label{eq:normalized-attention-generalized-gradient-flow}
 \partial_t\mu_t
 =-\mathsf K_{\mu_t}
 \frac{\delta\mathcal F_{\mathrm{att}}}{\delta\mu}[\mu_t]
 =\diver\left(\frac{\mu_t}{Z_{\mu_t}}
 \nabla\frac{\delta\mathcal F_{\mathrm{att}}}{\delta\mu}[\mu_t]\right)
 =-\diver\bigl(\mu_t\Phi[\mu_t]\bigr).
\end{equation}
Its dissipation identity is
\[
 \frac{\dd}{\dd t}\mathcal F_{\mathrm{att}}(\mu_t)
 =-\int_{\R^d}\frac1{Z_{\mu_t}(x)}
 \left|\nabla
 \frac{\delta\mathcal F_{\mathrm{att}}}{\delta\mu}[\mu_t](x)
 \right|^2\mu_t(\dd x)\leq0,
\]
or, equivalently,
\begin{equation}\label{eq:normalized-attention-energy-ascent}
 \frac{\dd}{\dd t}\mathcal J_\tau(\mu_t)
 =\int_{\R^d}\frac{\left|\widetilde\Phi[\mu_t](x)\right|^2}
 {Z_{\mu_t}(x)}\,\mu_t(\dd x)\geq0.
\end{equation}
However, the normalized attention equation is not the standard $W_2$ gradient flow: the factor multiplying the gradient at $x$ depends nonlocally on the current law through $Z_\mu(x)$.

Entropy makes the distinction between the two geometries especially clear.
For $\mu(\dd x)=\rho(x)\dd x$, recall
$\operatorname{Ent}(\rho)=\int\rho\log\rho\,\dd x$ from
\eqref{eq:boltzmann-entropy}, and define
\begin{equation}\label{eq:diffusive-attention-free-energy}
 \mathcal F_{\mathrm{att},\beta}(\rho)
 :=\mathcal F_{\mathrm{att}}(\rho)
   +\beta\operatorname{Ent}(\rho),\qquad \beta>0.
\end{equation}
In the standard Wasserstein geometry this produces the noisy
\emph{unnormalized} equation
\begin{equation}\label{eq:unnormalized-noisy-attention-fp}
 \partial_t\rho
 =\diver\left(\rho\nabla
 \frac{\delta\mathcal F_{\mathrm{att},\beta}}{\delta\rho}\right)
 =-\diver\bigl(\rho\widetilde\Phi[\rho]\bigr)+\beta\Delta\rho.
\end{equation}
For the normalized equation, applying the same positive operator
$\mathsf K_\rho$ to the same free energy produces instead the compatible
weighted diffusion
\begin{equation}\label{eq:normalized-weighted-diffusion}
 \partial_t\rho  = - \mathsf K_\rho
 \frac{\delta\mathcal F_{\mathrm{att},\beta}}{\delta\rho}
 =-\diver\bigl(\rho\Phi[\rho]\bigr)
  +\beta\diver\left(\frac1{Z_\rho}\nabla\rho\right),
\end{equation}
and
\[
 \frac{\dd}{\dd t}\mathcal F_{\mathrm{att},\beta}(\rho_t)
 =-\int_{\R^d}\frac{\rho_t}{Z_{\rho_t}}
 \left|\nabla
 \frac{\delta\mathcal F_{\mathrm{att},\beta}}{\delta\rho}[\rho_t]
 \right|^2\dd x\leq0.
\]
By contrast, adding ordinary Brownian noise to normalized attention gives
\begin{equation}\label{eq:noisy-attention-fp}
 \partial_t\rho+\diver\bigl(\rho\Phi[\rho]\bigr)=\beta\Delta\rho.
\end{equation}
The attention drift in \eqref{eq:noisy-attention-fp} contains the factor
$Z_\rho^{-1}$, whereas its Brownian diffusion does not.  Consequently, that
equation does not in general satisfy the preceding dissipation identity for
$\mathcal F_{\mathrm{att},\beta}$.

\paragraph{\textbf{Finite particles and the mean-field limit.}}
The variational structure already exists before taking $N\to\infty$.  For
$\mu^N=N^{-1}\sum_{i=1}^N\delta_{X^i}$, set
\[
 \mathcal F_{\mathrm{att}}^N(X)
 :=\mathcal F_{\mathrm{att}}(\mu^N)
 =-\frac{\tau}{2N^2}\sum_{i,j=1}^N
 e^{\ip{X^i}{X^j}/\tau}.
\]
Differentiating both occurrences of $X^i$, including the diagonal term,
gives
\[
 \nabla_{X^i}\mathcal F_{\mathrm{att}}^N(X)
 =-\frac1N\widetilde\Phi[\mu^N](X^i).
\]
Consequently, the unnormalized particle dynamics are
\begin{equation}\label{eq:unnormalized-attention-ips}
 \dot X_t^i
 =\frac1N\sum_{j=1}^N
 e^{\ip{X_t^i}{X_t^j}/\tau}X_t^j
 =-N\nabla_{X^i}\mathcal F_{\mathrm{att}}^N(X_t),
\end{equation}
while normalized identity attention satisfies
\begin{equation}\label{eq:normalized-attention-ips}
 \dot X_t^i
 =\frac{\sum_{j=1}^N e^{\ip{X_t^i}{X_t^j}/\tau}X_t^j}
 {\sum_{j=1}^N e^{\ip{X_t^i}{X_t^j}/\tau}}
 =-\frac{N}{Z_{\mu_t^N}(X_t^i)}
 \nabla_{X^i}\mathcal F_{\mathrm{att}}^N(X_t).
\end{equation}
Thus the first system is an ordinary Euclidean gradient descent after the
energy-per-particle scaling by $N$, whereas the second has a positive diagonal, state-dependent rescaling.

For either system, the empirical measure satisfies its continuity equation
exactly in the weak sense of \eqref{eq:token-weak}.  On a compact state space,
or on a finite interval on which the particle states have an a priori bound,
the exponential kernel is locally Lipschitz and $Z_\mu$ is bounded away from
zero.  The stability and propagation-of-chaos arguments of
Section~\ref{sec:ips} then justify passage from the finite IPS to
\eqref{eq:unnormalized-attention-gradient-flow} or
\eqref{eq:normalized-attention-generalized-gradient-flow}.  On all of
$\R^d$, this conclusion requires moment or growth estimates because the
exponential score is not globally Lipschitz.

If independent Brownian noise is added to a drift
$\Psi\in\{\widetilde\Phi,\Phi\}$, the particle model is
\begin{equation}\label{eq:noisy-attention-sde}
 \dd X_t^i=\Psi[\mu_t^N](X_t^i)\dd t+\sqrt{2\beta}\,\dd B_t^i.
\end{equation}
Its mean-field equation is \eqref{eq:noisy-attention-fp} with $\Phi$
replaced by $\Psi$.  The entropy interpretation is a single standard
Wasserstein gradient flow when $\Psi=\widetilde\Phi$; for
$\Psi=\Phi$, ordinary Brownian diffusion gives
\eqref{eq:noisy-attention-fp}.

The weighted diffusion in
\eqref{eq:normalized-weighted-diffusion} can be realized by multiplicative
noise, but the It\^o equation also needs a drift correction.  Assuming that
$Z_{\rho_t}$ is smooth and bounded away from zero, let
$\rho_t=\Law(X_t)$ and consider
\begin{equation}\label{eq:normalized-weighted-sde}
\begin{aligned}
 \dd X_t={}&\left[\Phi[\rho_t](X_t)
 +\beta\nabla\left(Z_{\rho_t}^{-1}\right)(X_t)\right]\dd t +\sqrt{\frac{2\beta}{Z_{\rho_t}(X_t)}}\,\dd B_t.
 \end{aligned}
\end{equation}
The Fokker--Planck equation associated with
\eqref{eq:normalized-weighted-sde} is
\[
\begin{aligned}
 \partial_t\rho_t
 &=-\diver\left(\rho_t
    \left[\Phi[\rho_t]+\beta\nabla Z_{\rho_t}^{-1}\right]\right)
   +\beta\Delta\left(\rho_t Z_{\rho_t}^{-1}\right)\\
 &=-\diver\bigl(\rho_t\Phi[\rho_t]\bigr)
   +\beta\diver\left(Z_{\rho_t}^{-1}\nabla\rho_t\right),
 \end{aligned}
\]
which is exactly \eqref{eq:normalized-weighted-diffusion}. 
% Here, the second line use the identity $\Delta{(fg)} = \diver (g \nabla f) + \diver (f \nabla g)$ with $f=Z_{\rho_t}^{-1}$ and $g=\rho_t$.
With the
same multiplicative noise but without the correction
$\beta\nabla Z_{\rho_t}^{-1}$, the diffusion term would instead be
$\beta\Delta(\rho_t/Z_{\rho_t})$.

\begin{example}[A frozen-measure potential]
\label{ex:identity-attention-potential}
Freeze $\mu$ and define
\[
Z_\mu(x):= \int e^{\ip{x}{y}/\tau}\,\mu(\dd y),\qquad
 P_\mu(x):=\tau\log Z_\mu(x),\qquad
 \pi_x^\mu(\dd y):=\frac{e^{\ip{x}{y}/\tau}}{Z_\mu(x)}\,\mu(\dd y).
\]
Differentiation with respect to the query gives
\begin{equation}\label{eq:identity-attention-log-partition}
 \nabla P_\mu(x)=\Phi[\mu](x),\qquad
 \nabla^2P_\mu(x)=\frac1\tau
 \operatorname{Cov}_{\pi_x^\mu}(Y)\succeq0.
\end{equation}
To derive the Hessian, differentiate the partition function twice:
\begin{align*}
 \nabla Z_\mu(x)   & = \int e^{\ip{x}{y}/\tau}\frac{y}{\tau}\,\mu(\dd y) 
                     =\frac{Z_\mu(x)}{\tau}\E_{\pi_x^\mu}[Y],\\
 \nabla^2 Z_\mu(x) & = \int e^{\ip{x}{y}/\tau}\frac{yy^\top}{\tau^2}\,\mu(\dd y)
                     =\frac{Z_\mu(x)}{\tau^2}\E_{\pi_x^\mu}[YY^\top].
\end{align*}
Since $P_\mu=\tau\log Z_\mu$, it follows that
\[
 \begin{aligned}
 \nabla^2P_\mu(x) = \nabla\left(\tau\frac{\nabla Z_\mu(x)}{Z_\mu(x)}\right)
 &=\tau\left(
 \frac{\nabla^2Z_\mu(x)}{Z_\mu(x)}
 -\frac{\nabla Z_\mu(x)\nabla Z_\mu(x)^\top}{Z_\mu(x)^2}
 \right)\\
 &=\frac1\tau\left(
 \E_{\pi_x^\mu}[YY^\top]
 -\E_{\pi_x^\mu}[Y]\E_{\pi_x^\mu}[Y]^\top
 \right)
 =\frac1\tau\operatorname{Cov}_{\pi_x^\mu}(Y).
 \end{aligned}
\]

Hence residual attention with a frozen key--value law,
$\dot x_t=\Phi[\mu](x_t)$, is gradient ascent of the convex log-partition
potential:
\[
 \frac{\dd}{\dd t}P_\mu(x_t)
 =\left|\nabla P_\mu(x_t)\right|^2\geq0.
\]
This identity makes $P_\mu$ a Lyapunov function and rules out a
nonstationary periodic orbit, but it does not by itself imply convergence:
gradient ascent of a nonconfining convex function can escape to infinity.

The consensus modification has a different and more useful energy:
\begin{equation}\label{eq:identity-attention-consensus-potential}
 \dot x_t=\Phi[\mu](x_t)-x_t=-\nabla V_\mu(x_t),
 \qquad
 V_\mu(x):=\frac12\abs{x}^2-P_\mu(x).
\end{equation}
Now
\[
 \frac{\dd}{\dd t}V_\mu(x_t)=-|\nabla V_\mu(x_t)|^2\leq0,
\]
and equilibria are precisely the fixed points
$x_\star=\Phi[\mu](x_\star)$ satisfying
$0=\nabla V_\mu(x_\star)=x_\star-\nabla P_\mu(x_\star)$. If
$\supp(\mu)\subset B_R(0)$, then
\[
 \nabla^2V_\mu(x)
 =I-\frac1\tau\operatorname{Cov}_{\pi_x^\mu}(Y)
 \succeq\left(1-\frac{R^2}{\tau}\right)I.
\]
Thus, when $\tau>R^2$, $V_\mu$ is strongly convex, it has a unique
equilibrium, and
\[
 |x_t-x_\star|
 \leq e^{-(1-R^2/\tau)t}|x_0-x_\star|.
\]
For smaller temperature the potential can be nonconvex, so several attracting
critical points, saddles, and long plateaus may coexist.  This is the
variational source of initialization-dependent clustering and metastability.
In the self-consistent spherical model, compactness and analyticity allow
Łojasiewicz arguments to upgrade energy convergence to convergence of
finite-particle trajectories, while well-separated clusters can persist for
exponentially long times
\cite{GeshkovskiEtAl2024Metastability,Rigollet2026MeanField}.  A Lyapunov
function therefore identifies the direction of evolution, but convexity or a
separate critical-point analysis is still needed for a unique equilibrium and
a quantitative rate.
\end{example}

\paragraph{\textbf{Scope and limitations of the variational structure.}}
The preceding calculation uses two special facts: the identity score is
symmetric in $x,y$, and the value vector $y$ is exactly the query gradient
of the score.  For general matrices, the unnormalized field is
\[
 \widetilde\Phi_{Q,K,V}[\mu](x)
 =\int e^{x^\top Q^\top Ky/\tau}Vy\,\mu(\dd y).
\]
Writing $A=Q^\top K$, the gradient of
$\tau e^{x^\top Ay/\tau}$ with respect to $x$ is
$e^{x^\top Ay/\tau}Ay$.  A symmetric pair energy therefore requires special
algebraic compatibility, for example $A=A^\top$ and $V$ proportional to
$A$.  Generic $Q,K,V$ do not satisfy this condition, and row normalization
cannot create an energy when the underlying pair field is not variational.

Several further qualifications matter.  Causal masks destroy pair symmetry;
time-dependent or untied parameters produce a time-dependent energy with an
additional $\partial_t\mathcal F_t$ term; value biases, MLPs, normalization,
and residual forcing add vector fields that need not share the same potential.
Even when a continuous-depth energy exists, nonconvexity permits multiple
equilibria and metastability, and a finite residual layer is only a forward
Euler step, whose energy monotonicity may require a step-size restriction.
Thus identity attention---and specially tied symmetric extensions---has the
gradient structure derived above; a generic Transformer block inherits the
continuity-equation description, but not this Lyapunov theory.

% \begin{exercise}[Noisy attention]
% Starting from \eqref{eq:noisy-attention-sde}, derive
% \eqref{eq:noisy-attention-fp} in weak form, including the quadratic-variation
% term.  Identify the free energy when
% $\Phi[\mu]=-\nabla(U+W*\mu)$.
% \end{exercise}
% \ifexeSolutionOn
% \begin{proof}[Solution]
% Let $\mu_t^N=N^{-1}\sum_i\delta_{X_t^i}$ and
% $\varphi\in C_b^2(\R^d)$.  It\^o's formula gives
% \[
%  \dd\varphi(X_t^i)
%  =\nabla\varphi(X_t^i)\cdot
%    \Phi_{\theta_t}[\mu_t^N](X_t^i)\dd t
%  +\beta\Delta\varphi(X_t^i)\dd t
%  +\sqrt{2\beta}\nabla\varphi(X_t^i)\cdot\dd B_t^i.
% \]
% The coefficient of $\Delta\varphi$ is
% $\frac12(2\beta)=\beta$.  Averaging over $i$ yields the exact finite-particle
% identity
% \begin{align}
%  \dd\int\varphi\,\dd\mu_t^N
%  &=\int\left[
%  \nabla\varphi\cdot\Phi_{\theta_t}[\mu_t^N]
%  +\beta\Delta\varphi\right]\dd\mu_t^N\dd t
%  +\dd M_t^{N,\varphi},\label{eq:noisy-attention-weak-finite}\\
%  M_t^{N,\varphi}
%  &=\frac{\sqrt{2\beta}}{N}\sum_{i=1}^N
%  \int_0^t\nabla\varphi(X_s^i)\cdot\dd B_s^i.
% \end{align}
% Independence of the Brownian motions gives
% \begin{equation}\label{eq:noisy-attention-bracket}
%  \left\langle M^{N,\varphi}\right\rangle_t
%  =\frac{2\beta}{N^2}\sum_{i=1}^N
%  \int_0^t\abs{\nabla\varphi(X_s^i)}^2\dd s
%  \leq\frac{2\beta t}{N}\norm{\nabla\varphi}_\infty^2.
% \end{equation}
% By It\^o's isometry and Doob's inequality, the bracket estimate implies
% the stronger maximal bound
% \[
%  \E\sup_{0\leq t\leq T}\abs{M_t^{N,\varphi}}^2
%  \leq4\E\abs{M_T^{N,\varphi}}^2
%  =4\E\left\langle M^{N,\varphi}\right\rangle_T
%  \leq\frac{8\beta T}{N}\norm{\nabla\varphi}_\infty^2.
% \]
% Thus the martingale vanishes in
% $L^2(\Omega;C([0,T]))$ for each fixed test function.  If the
% empirical measures and nonlinear fields converge strongly enough to pass to
% the drift, the limiting weak equation is
% \[
%  \int\varphi\,\dd\mu_t-\int\varphi\,\dd\mu_0
%  =\int_0^t\int\left[
%  \nabla\varphi\cdot\Phi_{\theta_s}[\mu_s]
%  +\beta\Delta\varphi\right]\dd\mu_s\dd s,
% \]
% which is precisely the weak form of \eqref{eq:noisy-attention-fp}.

% Now suppose that $W(z)=W(-z)$ and
% $\Phi[\mu]=-\nabla(U+W*\mu)$.  For $\mu(\dd x)=\rho(x)\dd x$, define
% \begin{equation}\label{eq:noisy-attention-free-energy}
%  \mathcal F_\beta(\rho)
%  =\int U\rho\,\dd x
%  +\frac12\iint W(x-y)\rho(x)\rho(y)\dd x\dd y
%  +\beta\int\rho\log\rho\,\dd x.
% \end{equation}
% Its first variation is
% $\delta\mathcal F_\beta/\delta\rho
% =U+W*\rho+\beta(1+\log\rho)$.  Hence, formally,
% \[
%  \partial_t\rho
%  =\diver\left(\rho\nabla
%  \frac{\delta\mathcal F_\beta}{\delta\rho}\right)
%  =\diver\bigl(\rho\nabla(U+W*\rho)\bigr)+\beta\Delta\rho,
% \]
% so \eqref{eq:noisy-attention-free-energy} is the required free energy.
% \end{proof}
%\fi

%%%%%%%%%%%%%%----------
\subsection{Four limits and a warning about their order}
There are several distinct regimes:
\begin{enumerate}[label=(\roman*),leftmargin=*]
 \item fixed $N$ and many small layers: Euler convergence to
 \eqref{eq:token-ode};
 \item fixed finite depth and $N\to\infty$: convergence of iterated nonlinear
 measure maps;
 \item simultaneous $N,L\to\infty$: both particle and discretization errors
 must be controlled, preferably with constants uniform in depth;
 \item $d\to\infty$ as well: score concentration and temperature scaling can
 change the limiting interaction.
\end{enumerate}
These limits need not commute.  In particular, a finite-time propagation of
chaos estimate may grow like $e^{CT}$ and become useless on the long time scale
needed for clustering or metastable transitions.

A genuine gradient-flow structure can change this long-time picture but
does not remove the need to distinguish the limits.  For example, geodesic
convexity of the energy can yield contractivity of the measure evolution and
replace an exponentially growing stability constant by a uniform or decaying
one.  The calculation in Example~\ref{ex:identity-attention-potential} shows
that an energy identity alone is not enough for contractivity: one must
also prove geodesic or ordinary convexity, or another mechanism that controls
the critical points and the dissipation rate.

\paragraph{\textbf{Literature and storyline.}}
The continuous-depth interpretation was assembled from several ideas rather
than introduced by one paper.  Residual networks first made a deep layer look
like the increment $x^{\ell+1}-x^\ell=F_\ell(x^\ell)$
\cite{HeEtAl2016ResNet}.  E then argued explicitly that deep architectures
should be organized using dynamical-systems and numerical-discretization
principles \cite{E2017DynamicalSystems}.  Chen, Rubanova, Bettencourt, and
Duvenaud parameterized a hidden-state derivative and used an ODE solver in
place of a prescribed layer grid \cite{ChenEtAl2018NeuralODE}.  The paper
develops adaptive evaluation and memory-efficient training, while Duvenaud's
research page summarizes the central object as a continuous-depth neural
network \cite{DuvenaudResearchProfile}.  For interacting tokens, however, the state is
the entire empirical measure, so depth simultaneously evolves each particle
and the field seen by every other particle.  This is why the section develops
the argument in the order Euler consistency $\to$ particle ODE $\to$
continuity equation $\to$ stability of the joint limits.  ``Many layers'' is
not by itself a theorem about an ODE: one needs residual scaling, a controlled
interpolation of the parameters, and bounds uniform in the depth and particle
limits.


\subsection{Connections forward}
Continuous-depth attention defines a nonlinear solution operator
$\mu_0\mapsto\mu_T$, with the current token law entering the velocity that
transports it.  The Euler estimate separates approximation by depth from
approximation by particle number, while the four-limit warning records why
neither can be hidden inside the other.

The gradient-flow test gives a
three-way distinction: symmetric unnormalized identity attention is a standard
Wasserstein gradient flow; normalized identity attention is a generalized
gradient flow in which the gradient is multiplied by $Z_\mu^{-1}$; and
generic $Q,K,V$ attention has, in
general, only the transport description.  The linear consensus example in
Section~\ref{sec:ips} is the quadratic-potential case
$W(z)=\frac12\abs z^2$.  The self-consistent identity-softmax energy is
\eqref{eq:attention-free-energy}, while the log-partition potential
\eqref{eq:identity-attention-log-partition} describes one query when the
key--value law is frozen.  These distinctions will guide
the analysis of clustering, stability, and the inverse problem in the
subsequent sections.

\section{Clustering, consensus, and collapse}
\label{sec:clustering}
\paragraph{\textbf{Key question}}
When does nonlocal averaging create useful clusters, consensus, or collapse?

\subsection{A consensus model}
Consider
\begin{equation}\label{eq:consensus}
    \dot X_t^i=\sum_{j=1}^Na_{ij}(t)(X_t^j-X_t^i),
    \qquad a_{ij}(t)\geq0,\qquad\sum_ja_{ij}(t)=1.
\end{equation}
The weights may depend on the entire configuration, as attention weights do,
but the first result uses only their sign and row sums.  For a unit vector $u$,
set
\begin{equation}\label{eq:directional-diameter}
 M_u(t)=\max_i\ip{u}{X_t^i},\quad
 m_u(t)=\min_i\ip{u}{X_t^i},\quad
 D_u(t)=M_u(t)-m_u(t).
\end{equation}
The quantity $D_u(t)$ is the \emph{directional diameter} of the configuration.



Note first that the directional diameter is nonincreasing. To see this, write $z_i(t)=\ip{u}{X_t^i}$.  The projected velocity is simply the scalar derivative of this projection:
\begin{equation}\label{eq:projected-velocity}
 \dot z_i(t)=\ip{u}{\dot X_t^i}
 =\sum_{j=1}^Na_{ij}(t)(z_j(t)-z_i(t)).
\end{equation}
If $z_i=M_u$, every difference $z_j-z_i$ is nonpositive, so
$\dot z_i\leq0$; if $z_i=m_u$, every difference is nonnegative, so
$\dot z_i\geq0$.  Since each $z_i$ is continuously differentiable,
the finite maximum $M_u$ and minimum $m_u$ are locally Lipschitz.  Therefore
$D_u=M_u-m_u$ is absolutely continuous and differentiable almost everywhere,
but it need not be differentiable when the identity of an extremizer changes
(the elementary example $\max\{t,-t\}=|t|$ shows why).  At every time, tied
extrema are handled by the upper and lower right Dini derivatives.\footnote{For
a real function $f$, the upper and lower right Dini derivatives are
$D^+f(t)=\limsup_{h\downarrow0}[f(t+h)-f(t)]/h$ and
$D_+f(t)=\liminf_{h\downarrow0}[f(t+h)-f(t)]/h$.  For a finite maximum or
minimum of $C^1$ functions, the active-index formulas are recorded in
\eqref{eq:appendix-dini-maximum}.} In fact,
\[
 D^+M_u(t)=\max_{i:z_i(t)=M_u(t)}\dot z_i(t)\leq0,
 \qquad
 D_+m_u(t)=\min_{i:z_i(t)=m_u(t)}\dot z_i(t)\geq0,
\]
which follow by choosing an index that realizes the extremum at $t+h$, passing
to a constant-index subsequence as $h\downarrow0$, and applying the first-order
expansion of that $z_i$; the reverse inequality follows by testing every index
that is active at time $t$.
Consequently $D_u'(t)\leq0$ at almost every differentiability time, so
absolute continuity implies that $D_u$ is nonincreasing. Theorem~\ref{thm:diameter} below shows that under the uniform communication assumption, the directional diameter contracts exponentially, and the Euclidean diameter contracts as well.

Also, the convex hull of the configuration is forward invariant.
Let
$C_0=\operatorname{co}\{X_0^1,\ldots,X_0^N\}$ be the smallest convex set
containing the initial configuration.  Its support function is
$h_{C_0}(u)=\max_i\ip{u}{X_0^i}=M_u(0)$, and the supporting-half-space
representation of a compact convex set gives
\[
 C_0=\bigcap_{|u|=1}\{x\in\R^d:\ip{u}{x}\leq M_u(0)\}.
\]
For every $i$ and every unit $u$,
$\ip{u}{X_t^i}\leq M_u(t)\leq M_u(0)$; hence $X_t^i\in C_0$.  This is what it
means for the convex hull to be forward invariant: pure averaging cannot
create a token state outside the region spanned by the initial tokens.  The
fact is useful in three ways.  It supplies the a priori bound
$\max_i|X_t^i|\leq\max_i|X_0^i|$, keeps state-dependent attention scores in a
controlled range, and identifies a structural limitation of the consensus
model.  Value maps, MLPs, or residual forcing can move representations outside
this hull, so hull invariance should not be attributed to a general
Transformer block.



\begin{assumption}[Uniform communication]\label{ass:uniform-mixing}
There is $\alpha\in(0,1]$ such that
\[
  a_{ij}(t)\geq\frac{\alpha}{N}
  \qquad\text{for all }i,j,t.
\]
\end{assumption}

\begin{theorem}[Directional diameter contraction]\label{thm:diameter}
Under Assumption~\ref{ass:uniform-mixing},
\begin{equation}\label{eq:diameter-contraction}
    D_u(t)\leq e^{-\alpha t}D_u(0)
    \qquad\text{for every unit }u.
\end{equation}
Consequently,
$\max_{i,j}\abs{X_t^i-X_t^j}\leq
e^{-\alpha t}\max_{i,j}\abs{X_0^i-X_0^j}$.
\end{theorem}
\begin{proof}
If $\alpha=1$, the lower bound and the row sums force
$a_{ij}=1/N$ for all $i,j$, and the differential inequalities below follow
directly.  Suppose henceforth that $0<\alpha<1$ and write
\[
 a_{ij}=\frac{\alpha}{N}+(1-\alpha)\widetilde a_{ij},
 \qquad \sum_j\widetilde a_{ij}=1.
\]
Then, the projected equation \eqref{eq:projected-velocity} gives
\begin{equation}\label{eq:consensus-uniform-remainder}
 \dot z_i
 =\alpha(\overline z-z_i)
 +(1-\alpha)\sum_j\widetilde a_{ij}(z_j-z_i),
 \qquad \overline z=\frac1N\sum_j z_j.
\end{equation}
 Now, for each $i$ such that $z_i=M_u(t)$, every $z_j-M_u\leq0$, so
\[
 \dot z_i
 =\alpha(\overline z-M_u)
 +(1-\alpha)\sum_j\widetilde a_{ij}(z_j-M_u)
 \leq\alpha(\overline z-M_u),
\]
because every summand in the remainder is nonpositive.  For each
$i$ such that $z_i=m_u(t)$, the same calculation, with every $z_j-m_u\geq0$, gives
$\dot z_i\geq\alpha(\overline z-m_u)$.  Taking the maximum over $I_M(t)$ and
the minimum over $I_m(t)$ therefore gives
\[
 D^+M_u\leq\alpha(\overline z-M_u),
 \qquad
 D_+m_u\geq\alpha(\overline z-m_u).
\]
It follows that
\[
 D^+D_u\leq D^+M_u-D_+m_u
 \leq-\alpha(M_u-m_u)=-\alpha D_u.
\]
Equivalently, the same differential inequality holds for $D_u'$ almost
everywhere.  The integrating-factor form of Gronwall then gives
\eqref{eq:diameter-contraction}.


For the Euclidean conclusion, fix a
diameter-realizing pair $(i,j)$ at time $t$.  If its distance is zero, the
claim is immediate.  Otherwise choose
\[
 u=\frac{X_t^i-X_t^j}{\abs{X_t^i-X_t^j}}.
\]
By the definitions of the maximum and minimum,
\[
 D_u(t)=\max_k\ip{u}{X_t^k}-\min_r\ip{u}{X_t^r}
 \geq\ip{u}{X_t^i}-\ip{u}{X_t^j}
 =\abs{X_t^i-X_t^j}.
\]
Therefore
\[
 \abs{X_t^i-X_t^j}
 \leq D_u(t)
 \leq e^{-\alpha t}D_u(0)
 \leq e^{-\alpha t}\max_{k,r}\abs{X_0^k-X_0^r},
\]
which proves the stated bound.
\end{proof}



\begin{exercise}[Uniform consensus]
Solve \eqref{eq:consensus} when $a_{ij}=1/N$.  Identify the conserved quantity
and exact decay rate.  \emph{Hint:} first derive
$\dot X^i=\overline X-X^i$.
\end{exercise}
\begin{proof}[Solution]
When $a_{ij}=1/N$,
\[
 \dot X_t^i=\frac1N\sum_{j=1}^N(X_t^j-X_t^i)
 =\overline X_t-X_t^i,
 \qquad
 \overline X_t=\frac1N\sum_{j=1}^NX_t^j.
\]
Averaging this equation over $i$ gives $\dot{\overline X}_t=0$, so
$\overline X_t=\overline X_0$.  Consequently,
\begin{equation}\label{eq:uniform-consensus-solution}
 X_t^i=\overline X_0+e^{-t}(X_0^i-\overline X_0).
\end{equation}
Every pairwise difference therefore satisfies
$X_t^i-X_t^j=e^{-t}(X_0^i-X_0^j)$, and the exact directional and Euclidean
diameter decay rate is $e^{-t}$.
\end{proof}


\paragraph{\textbf{What is conserved?}}
Row stochasticity alone does not conserve the arithmetic mean.  Summing
\eqref{eq:consensus} gives
\begin{equation}\label{eq:mean-column-sum}
 \frac{\dd}{\dd t}\overline X_t
 =\frac1N\sum_j\left(\sum_i a_{ij}(t)-1\right)X_t^j.
\end{equation}
If $A(t)$ is also column stochastic, then $\overline X_t$ is constant.  If it
is symmetric, the variance
$\mathcal V_t=N^{-1}\sum_i\abs{X_t^i-\overline X_t}^2$ satisfies
\begin{equation}\label{eq:variance-dissipation}
 \frac{\dd}{\dd t}\mathcal V_t
 =-\frac1N\sum_{i,j}a_{ij}(t)\abs{X_t^i-X_t^j}^2\leq0.
\end{equation}
Indeed, symmetry and row stochasticity imply column stochasticity, so
$\overline X_t$ is constant.  Writing $Z_i=X_t^i-\overline X_t$ and suppressing
the time variable, symmetry gives
\begin{align*}
  N\frac{\dd}{\dd t}\mathcal V_t
  & = 2\sum_iZ_i\cdot\dot Z_i
    = 2\sum_iZ_i\cdot\dot X_i
    = 2\sum_{i,j}a_{ij}Z_i\cdot(Z_j-Z_i)  \\
  & =\sum_{i,j}a_{ij}\left[
 Z_i\cdot(Z_j-Z_i)+Z_j\cdot(Z_i-Z_j)\right]
 =-\sum_{i,j}a_{ij}\abs{Z_i-Z_j}^2.
\end{align*}
This is a Lyapunov identity.  Standard row-softmax attention does not supply
it, which is one motivation for symmetric or doubly-stochastic attention when
conservation and predictable contraction are desired.



\paragraph{\textbf{Why softmax can mix slowly}}
Suppose all scores in one head lie in $[-S,S]$.  Then
\[
 a_{ij}\geq\frac{e^{-S}}{Ne^S}=\frac{e^{-2S}}N.
\]
Theorem~\ref{thm:diameter} applies with $\alpha=e^{-2S}$.  If
$S\asymp\tau^{-1}$ is large, the guaranteed contraction time is exponentially
large in inverse temperature.  This is not merely a weakness of the proof:
sharp softmax can create nearly disconnected interaction graphs.

At high temperature, weights are nearly uniform and the system rapidly pools
everything into one cluster.  At low temperature, each query follows a few
leaders; within-group relaxation can be fast while communication between
groups is extremely slow.  The resulting multi-cluster state can be
\emph{metastable}: stable over the depth horizon of interest even if eventual
consensus occurs.




\paragraph{\textbf{Consensus results for identity attention.}}
 For Euclidean identity-softmax consensus,
\begin{equation}\label{eq:identity-softmax-consensus}
 \dot X_t^i=\Phi[\mu_t^N](X_t^i)-X_t^i
 =\sum_{j=1}^Na_{ij}(X_t)(X_t^j-X_t^i),
 \qquad
 a_{ij}(X)=\frac{e^{\ip{X^i}{X^j}/\tau}}
 {\sum_{k=1}^Ne^{\ip{X^i}{X^k}/\tau}}.
\end{equation}
Let $R_0=\max_i|X_0^i|$.  Forward invariance of the initial convex hull gives
$|X_t^i|\leq R_0$, and hence
\[
 a_{ij}(X_t)
 \geq\frac{e^{-R_0^2/\tau}}{Ne^{R_0^2/\tau}}
 =\frac{\alpha_\tau}{N},
 \qquad \alpha_\tau:=e^{-2R_0^2/\tau}.
\]
Theorem~\ref{thm:diameter} therefore applies to every initial configuration:
\[
 \max_{i,j}|X_t^i-X_t^j|
 \leq e^{-\alpha_\tau t}
 \max_{i,j}|X_0^i-X_0^j|,
\]
and the shrinking convex hulls converge to a single point
$X_\infty\in\operatorname{co}\{X_0^1,\ldots,X_0^N\}$.  The rate can be
exponentially small when $\tau$ is small, so this theorem is consistent with
long-lived apparent clusters.  The subtraction of $X_t^i$ is essential: the
pure residual flow $\dot X_t^i=\Phi[\mu_t^N](X_t^i)$ is not a convex-combination
consensus system and the preceding diameter argument does not apply.


For projected identity attention on the sphere in Section \ref{sec:spherical-attention}, the conclusion is instead almost-global: Theorem~\ref{thm:spherical-attention-clustering} gives
single-cluster convergence for almost every finite-particle initialization,
while exceptional antipodal or other unstable stationary configurations must
be excluded.  Mean-field exponential convergence to a Dirac mass is also
available under regularity and parameter assumptions
\cite{ChenEtAl2026Quantitative,Rigollet2026MeanField}; separated configurations
can nevertheless remain metastable for exponentially long times
\cite{GeshkovskiEtAl2024Metastability}.


\subsection{Clustering results for projected identity attention on the sphere}\label{sec:spherical-attention}
The preceding consensus theorem uses only stochastic-matrix bounds.  We
now keep the state dependence of softmax and analyze a standard idealized
model directly.  Let $x_i(t)\in\mathbb S^{d-1}$, let $\gamma\geq0$ be the
inverse temperature, and write
$$\mathbf P_x^\perp=I-xx^\top$$
for the orthogonal projection onto the tangent space
$T_x\mathbb S^{d-1}=\{v\in\R^d:x^\top v=0\}$.  Indeed,
$(\mathbf P_x^\perp)^\top=\mathbf P_x^\perp$,
$(\mathbf P_x^\perp)^2=\mathbf P_x^\perp$, and
$x^\top\mathbf P_x^\perp v=0$ for every $v\in\R^d$.  Thus
$\mathbf P_x^\perp$ fixes every tangent vector and removes the component in
the normal direction $x$, which characterizes the orthogonal projection.
The normalized spherical self-attention dynamics are
\begin{equation}\label{eq:spherical-sa}
 \dot x_i=\mathbf P_{x_i}^\perp\sum_{j=1}^Na_{ij}(X)x_j,
 \qquad
 a_{ij}(X)=\frac{e^{\gamma\ip{x_i}{x_j}}}{Z_i(X)},
 \qquad
 Z_i(X)=\sum_{k=1}^Ne^{\gamma\ip{x_i}{x_k}}.
\end{equation}
The unnormalized variant is
\begin{equation}\label{eq:spherical-usa}
 \dot x_i=\mathbf P_{x_i}^\perp\frac1N
 \sum_{j=1}^Ne^{\gamma\ip{x_i}{x_j}}x_j.
\end{equation}
In either case,
$\frac{\dd}{\dd t}\abs{x_i}^2=2\ip{x_i}{\dot x_i}=0$, so the sphere is
invariant.

Define the average similarity energy as
\begin{equation}\label{eq:spherical-attention-energy}
 \mathcal J_{\gamma,N}(X)
 =\frac1{2\gamma N^2}\sum_{i,j=1}^N
 \left(e^{\gamma\ip{x_i}{x_j}}-1\right),
\end{equation}
with the continuous extension
$\mathcal J_{0,N}(X)=(2N^2)^{-1}\sum_{i,j}\ip{x_i}{x_j}$ at
$\gamma=0$.  For $\gamma>0$, symmetry gives
\[
 \nabla_{x_i}\mathcal J_{\gamma,N}(X)
 =\frac1{N^2}\sum_{j=1}^Ne^{\gamma\ip{x_i}{x_j}}x_j,
 \qquad
 \nabla_{\mathbb S,x_i}\mathcal J_{\gamma,N}
 :=\mathbf P_{x_i}^\perp\nabla_{x_i}\mathcal J_{\gamma,N}.
\]
Consequently, \eqref{eq:spherical-usa} is the spherical gradient ascent
$\dot x_i=N\nabla_{\mathbb S,x_i}\mathcal J_{\gamma,N}$, while
\eqref{eq:spherical-sa} is a diagonally rescaled ascent
$ \dot x_i=\frac{N^2}{Z_i(X)}
 \nabla_{\mathbb S,x_i}\mathcal J_{\gamma,N}$.
Both identities yield an exact Lyapunov calculation:
\begin{align}
 \frac{\dd}{\dd t}\mathcal J_{\gamma,N}(X_t)
 &=N\sum_i\left|\nabla_{\mathbb S,x_i}
 \mathcal J_{\gamma,N}(X_t)\right|^2
 &&\text{for \eqref{eq:spherical-usa}},
 \label{eq:usa-energy-ascent}\\
 \frac{\dd}{\dd t}\mathcal J_{\gamma,N}(X_t)
 &=\sum_i\frac{N^2}{Z_i(X_t)}
 \left|\nabla_{\mathbb S,x_i}
 \mathcal J_{\gamma,N}(X_t)\right|^2
 &&\text{for \eqref{eq:spherical-sa}}.
 \label{eq:sa-energy-ascent}
\end{align}
Thus every nonstationary trajectory raises the average similarity energy.
This is the direct state-dependent replacement for the fixed-weight diameter
argument.

\begin{theorem}[Generic global clustering for spherical attention]
\label{thm:spherical-attention-clustering}
Let $N\geq2$, $d\geq3$, and $\gamma\geq0$.  For either
\eqref{eq:spherical-sa} or \eqref{eq:spherical-usa}, for almost every initial
configuration in $(\mathbb S^{d-1})^N$ there exists
$x_\infty\in\mathbb S^{d-1}$ such that
\[
 \lim_{t\to\infty}x_i(t)=x_\infty,
 \qquad i=1,\ldots,N.
\]
\end{theorem}
This result is presented, with its underlying synchronization references, in
\cite{Rigollet2026MeanField}; see also
\cite{GeshkovskiEtAl2023Clusters}.  Its proof has three structural steps.
First, compactness of the sphere gives global trajectories.  Second,
\eqref{eq:usa-energy-ascent}--\eqref{eq:sa-energy-ascent} and the
{\L}ojasiewicz inequality for the analytic energy imply convergence of each
trajectory to a stationary configuration, not merely convergence of its
energy.  Third, every non-consensus stationary configuration has an unstable
escape direction; the center-stable manifold theorem places the initial data
converging to such saddles in a measure-zero set.  The phrase ``almost every''
is essential because exceptional stationary configurations do exist.

%% Exe or skip 
% \begin{proof}
% We make the three steps above precise.  Let
% $\mathcal M=(\mathbb S^{d-1})^N$ with its product-sphere structure, and define
% \[
%  \varphi_\gamma(s)=
%  \begin{cases}
%   (e^{\gamma s}-1)/\gamma,&\gamma>0,\\
%   s,&\gamma=0.
%  \end{cases}
% \]
% Then
% $\mathcal J_{\gamma,N}(X)=(2N^2)^{-1}\sum_{i,j}
% \varphi_\gamma(\ip{x_i}{x_j})$ is real analytic on $\mathcal M$, and its
% gradient for the standard product metric is
% \[
%  \bigl(\operatorname{grad}\mathcal J_{\gamma,N}(X)\bigr)_i
%  =\frac1{N^2}\mathbf P_{x_i}^{\perp}
%    \sum_{j=1}^N e^{\gamma\ip{x_i}{x_j}}x_j.
% \]
% Thus \eqref{eq:spherical-usa} is gradient ascent for the metric
% \[
%  g_X^{\mathrm{USA}}(U,V)=\frac1N\sum_{i=1}^N\ip{u_i}{v_i},
% \]
% whereas \eqref{eq:spherical-sa} is gradient ascent for the analytic metric
% \[
%  g_X^{\mathrm{SA}}(U,V)
%  =\sum_{i=1}^N\frac{Z_i(X)}{N^2}\ip{u_i}{v_i}.
% \]
% Indeed, the corresponding metric gradients multiply the standard
% product-metric gradient by $N$ and by $N^2/Z_i(X)$, respectively.  Since
% $\mathcal M$ is compact and $Z_i>0$, both vector fields are smooth and
% complete.  The {\L}ojasiewicz gradient inequality for analytic functions and
% analytic Riemannian metrics therefore implies that every trajectory converges
% to one critical point of $\mathcal J_{\gamma,N}$.

% It remains to identify which critical points can attract a set of positive
% volume.  Write the configuration as the column matrix
% $\mathbf X=[x_1|\cdots|x_N]\in\R^{d\times N}$ and set
% \[
%  G=\mathbf X^\top\mathbf X,
%  \qquad A_{ij}=e^{\gamma G_{ij}},
%  \qquad D=\operatorname{diag}\bigl((GA)_{11},\ldots,(GA)_{NN}\bigr),
%  \qquad S=D-A.
% \]
% The first-order criticality condition is
% \begin{equation}\label{eq:spherical-critical-matrix}
%  \mathbf X S=0.
% \end{equation}
% We next show that a critical point whose Riemannian Hessian is negative
% semidefinite must be a consensus point.  Let
% $U=[u_1|\cdots|u_N]\in T_X\mathcal M$, so that $\ip{u_i}{x_i}=0$.  Along
% the product-sphere geodesic with initial velocity $U$, write
% $G_{ij}(s)=\ip{x_i(s)}{x_j(s)}$.  Then
% \[
%  G_{ij}'(0)=\ip{u_i}{x_j}+\ip{x_i}{u_j},
%  \qquad
%  G_{ij}''(0)=2\ip{u_i}{u_j}
%  -\bigl(\abs{u_i}^2+\abs{u_j}^2\bigr)G_{ij}.
% \]
% Consequently, direct differentiation gives
% \begin{align}
%  N^2\operatorname{Hess}\mathcal J_{\gamma,N}(X)[U,U]
%  ={}&-\left\langle S,U^\top U\right\rangle_{\mathrm F}\notag\\
%  &+\frac{\gamma}{2}\sum_{i,j=1}^N A_{ij}
%  \bigl(\ip{u_i}{x_j}+\ip{x_i}{u_j}\bigr)^2.
%  \label{eq:spherical-hessian-form}
% \end{align}
% The second line is zero when $\gamma=0$.  Negative semidefiniteness of the
% Hessian therefore implies
% \begin{equation}\label{eq:spherical-S-positive}
%  \left\langle S,U^\top U\right\rangle_{\mathrm F}\geq0
%  \qquad\text{for every }U\in T_X\mathcal M.
% \end{equation}

% To exploit \eqref{eq:spherical-S-positive}, let
% $\zeta\sim\mathcal N(0,I_d)$ and choose the common-direction tangent
% perturbation $u_i=(I-x_ix_i^\top)\zeta$.  Since
% \[
%  \E\ip{u_i}{u_j}
%  =\operatorname{tr}\bigl((I-x_ix_i^\top)(I-x_jx_j^\top)\bigr)
%  =d-2+G_{ij}^2,
% \]
% taking expectations in \eqref{eq:spherical-S-positive} yields
% \[
%  0\leq\left\langle
%  S,(d-2)\one\one^\top+G^{\odot2}\right\rangle_{\mathrm F}.
% \]
% The identity
% \[
%  (d-2)\one\one^\top+G^{\odot2}
%  =(d-3)\one\one^\top+(\one\one^\top-G)^{\odot2}+2G
% \]
% and \eqref{eq:spherical-critical-matrix} give
% $\langle S,G\rangle_{\mathrm F}=0$.  Moreover,
% \[
%  \left\langle S,\one\one^\top\right\rangle_{\mathrm F}
%  =-\sum_{i,j=1}^N A_{ij}(1-G_{ij})\leq0,
% \]
% and, because $(\one\one^\top-G)^{\odot2}$ has zero diagonal,
% \[
%  \left\langle S,(\one\one^\top-G)^{\odot2}\right\rangle_{\mathrm F}
%  =-\sum_{i,j=1}^N A_{ij}(1-G_{ij})^2.
% \]
% Since $d\geq3$, we conclude that
% \[
%  0\leq(d-3)\left\langle S,\one\one^\top\right\rangle_{\mathrm F}
%  -\sum_{i,j=1}^N A_{ij}(1-G_{ij})^2\leq0.
% \]
% Every $A_{ij}$ is strictly positive, hence $G_{ij}=1$ for all $i,j$, which
% means $x_1=\cdots=x_N$.  Thus every non-consensus critical point has a
% strictly positive Hessian direction: it is a strict saddle for gradient
% ascent.  This conclusion is independent of which of the two positive-definite
% metrics above is used.  At a critical point, the Hessian quadratic form is the
% second differential of the energy, and multiplication by a positive mobility
% preserves the existence of an unstable direction.

% Finally, the center-stable manifold theorem, in its standard avoidance form
% for gradient flows, implies that the set of initial conditions whose limiting
% critical point has a strictly positive Hessian direction has zero
% product-sphere volume.  Hence, for almost every initial condition, the limit
% established above has negative semidefinite Hessian and therefore, by the
% landscape calculation, is a consensus point
% $(x_\infty,\ldots,x_\infty)$.  This proves the claim.  The dynamical reduction
% and the direct landscape calculation are given in
% \cite[Theorem~1 and Appendix~C]{CriscitielloEtAl2026Synchronization}.
% \end{proof}


\begin{example}[Two tokens: the complete calculation]
\label{ex:two-token-spherical-attention}
Take $N=2$ in \eqref{eq:spherical-sa} and set
$c(t)=\ip{x_1(t)}{x_2(t)}$.  Since the projected self-interaction vanishes,
\[
 \dot x_1=a_\gamma(c)(x_2-cx_1),\qquad
 \dot x_2=a_\gamma(c)(x_1-cx_2),
 \qquad
 a_\gamma(c)=\frac1{1+e^{\gamma(1-c)}}.
\]
Taking the derivative of the inner product gives the closed scalar equation
\begin{equation}\label{eq:two-token-correlation}
 \dot c(t)= \ip{\dot x_1}{x_2} + \ip{x_1}{\dot x_2}
 =   2a_\gamma(c(t))(1-c(t)^2).
\end{equation}
If $c(0)>-1$, then $c$ is increasing and its only possible limit is $1$.
Moreover,
$a_\gamma(c)\geq(1+e^{2\gamma})^{-1}=:a_{\min}$ and
$1+c(t)\geq1+c(0)$, so
\[
 \frac{\dd}{\dd t}(1-c(t))
 \leq-2a_{\min}(1+c(0))(1-c(t)).
\]
Therefore
\begin{equation}\label{eq:two-token-exponential-clustering}
 1-c(t)\leq(1-c(0))
 e^{-2a_{\min}(1+c(0))t}.
\end{equation}
Since $\abs{x_1-x_2}^2=2(1-c)$, the two tokens synchronize exponentially.
The antipodal configuration $c(0)=-1$ is stationary and illustrates exactly
the exceptional saddle excluded by Theorem~\ref{thm:spherical-attention-clustering}.
\end{example}

TBV: Three refinements explain why this elementary conclusion is not the whole
story.  Multi-cluster states can persist for exponentially long low-temperature
times before merging \cite{GeshkovskiEtAl2024Metastability}; causal masking
destroys the mean-field gradient structure but can still lead to one-cluster
convergence \cite{KaragodinEtAl2024Causal}; and, at the PDE level, quantitative
exponential convergence to a Dirac mass is known only under stated regularity,
nondegeneracy, and parameter regimes
\cite{ChenEtAl2026Quantitative,Rigollet2026MeanField}.  These are results for
idealized tied-parameter flows.  They do not by themselves establish collapse
for a finite-depth trained Transformer with untied blocks, MLPs, residual
channels, or multiple heads.





\subsection{Useful clustering versus representation collapse}
Clustering can pool examples that share a latent task, select a leader, denoise
redundant patches, or create a task-relevant sufficient statistic.  Complete
consensus, however, removes token distinctions and can reduce the rank of the
representation.  A geometric statistic such as $\mathcal V_t$ therefore cannot
by itself declare success or failure.  One should also track:
\begin{enumerate}[label=(\roman*),leftmargin=*]
 \item a task-relevant readout or sufficient statistic;
 \item covariance spectrum or effective rank;
 \item cluster masses and margins;
 \item sensitivity of the terminal prediction to perturbations;
 \item whether information survives in residual coordinates not represented by
 the simplified consensus model.
\end{enumerate}
Finite-depth transient clusters may be more useful than the infinite-depth
limit.  Conversely, a model with several stationary clusters is not thereby a
reasoning system; the clusters must encode a verifiable computational
quantity.



\paragraph{\textbf{Literature and storyline.}}
Consensus models make the geometry of attention visible in its simplest
form.  DeGroot studied repeated averaging as a mechanism by which agents can
reach a common opinion \cite{DeGroot1974}; Cucker and Smale showed how
distance-dependent communication can generate flocking
\cite{CuckerSmale2007}.  The same mathematical motif---state-dependent
weighted averaging---reappears when tokens attend to one another.  Rigollet's
research account describes a recent program, developed with collaborators,
that treats self-attention as a velocity field of an interacting particle
system and asks when tokens cluster \cite{RigolletResearchProfile}; an MIT
profile records how his work moved across optimization, statistics, and
machine learning, the combination needed for this question
\cite{RigolletMITProfile2018}.  The new difficulty is that softmax weights are
directed, normalized by the current cloud, and often very sensitive to
temperature.  Classical consensus therefore supplies the baseline theorem,
not the conclusion: contraction of the convex hull is easy, strict consensus
requires persistent communication, and multiple clusters or metastable states
appear when that communication degenerates.  This is exactly the distinction
between useful representation formation and indiscriminate token collapse.

\subsection{Connections forward}
Consensus is an attractive gradient flow and hence links directly to the
energy-dissipation viewpoint of Section~\ref{sec:gradient-flow}.  In generative modeling, excessive
attraction resembles mode collapse, while diffusion or repulsion maintains
diversity.  In in-context learning, clusters can represent posterior mass on
different latent tasks, but averaging incompatible task predictions creates
bias.  For inverse problems, metastable trajectories explore only part of state
space and can make interaction laws outside the visited clusters
unidentifiable.  Architecturally, head diversity, repulsive heads, gating,
normalization, and doubly-stochastic constraints are mechanisms for positioning
the network between uncontrolled fragmentation and one-cluster collapse.



\section{Stability and propagation of perturbations}\label{sec:stability}
\paragraph{\textbf{Key question}}
How does an error in tokens or interaction kernels propagate through depth?

\subsection{Wasserstein stability under model discrepancy}\label{sec:w1-stability}
Let $\mu_t$ and $\nu_t$ solve continuity equations driven by fields $\Phi$ and
$\Psi$:
\[
 \partial_t\mu_t+\diver(\mu_t\Phi_t[\mu_t])=0,\qquad
 \partial_t\nu_t+\diver(\nu_t\Psi_t[\nu_t])=0.
\]
Differences can enter through the initial tokens, the interaction law, the
layer parameters, numerical discretization, or finite-particle sampling.  A
useful bound should keep these sources separate.

\begin{assumption}[Lipschitz field and model discrepancy]\label{ass:stability}
For measures supported in the domain visited up to time $T$, assume
\begin{align}
 \abs{\Phi_t[\mu](x)-\Phi_t[\nu](y)}
 &\leq L_x\abs{x-y}+L_\mu W_1(\mu,\nu),
 \label{eq:field-measure-lip}\\
 \sup_{t,x,\eta}
 \abs{\Phi_t[\eta](x)-\Psi_t[\eta](x)}
 &\leq\varepsilon.
 \label{eq:field-discrepancy}
\end{align}
\end{assumption}
The supremum in \eqref{eq:field-discrepancy} can be restricted to a
data-supported region; outside it the theorem gives no guarantee.


\begin{theorem}[Wasserstein stability]\label{thm:w1-stability}
Under Assumption~\ref{ass:stability}, let
$\mu_0,\nu_0\in\cP_1(\R^d)$, and let $L=L_x+L_\mu$.  Then
\begin{equation}\label{eq:stability}
    W_1(\mu_t,\nu_t)
    \leq e^{Lt}W_1(\mu_0,\nu_0)
       +\frac{e^{Lt}-1}{L}\varepsilon,
\end{equation}
where $(e^{Lt}-1)/L$ is interpreted as $t$ when $L=0$.
\end{theorem}
The exponential factor in \eqref{eq:stability} is sharp for
$\dot x=Lx$.  Even a small finite-time $W_1$ error can also change the selected
long-time attractor when the initial law is near a basin boundary.  

\begin{proof}
Choose an optimal initial coupling
$(X_0,Y_0)$.  Evolve it by
\[
 \dot X_t=\Phi_t[\mu_t](X_t),\qquad
 \dot Y_t=\Psi_t[\nu_t](Y_t).
\]
Its law at time $t$ is a coupling of $\mu_t$ and $\nu_t$.
Set $D_t=X_t-Y_t$.  Since $D_t$ is absolutely continuous, at times
when $D_t\ne0$,
$\frac{\dd}{\dd t}(|D_t|^2)^{1/2}= \langle D_t/|D_t|,\dot D_t\rangle
\leq|\dot D_t|$; the corresponding upper-derivative inequality holds when
$D_t=0$.  Integrating this pathwise inequality and then using Fubini gives the
first line below.  For the second line, use the decomposition
\[
\Phi_t[\mu_t](X_t)-\Psi_t[\nu_t](Y_t)
 =\bigl(\Phi_t[\mu_t](X_t)-\Phi_t[\nu_t](Y_t)\bigr)
  +\bigl(\Phi_t[\nu_t](Y_t)-\Psi_t[\nu_t](Y_t)\bigr).
\]
Assumption~\ref{ass:stability} bounds the two terms by
$L_x|D_t|+L_\mu W_1(\mu_t,\nu_t)$ and $\varepsilon$, respectively.
\begin{align*}
 \frac{\dd}{\dd t}\E\abs{X_t-Y_t}
 &\leq\E\abs{\Phi_t[\mu_t](X_t)
                 -\Psi_t[\nu_t](Y_t)}\\
 &\leq L_x\E\abs{X_t-Y_t}
       +L_\mu W_1(\mu_t,\nu_t)+\varepsilon\\
 &\leq L\E\abs{X_t-Y_t}+\varepsilon.
\end{align*}
Gronwall's inequality gives  $\E\abs{X_t-Y_t}\leq e^{Lt}\E\abs{X_0-Y_0}+\frac{e^{Lt}-1}{L}\varepsilon
 $. Then, we obtain \eqref{eq:stability} since
$W_1(\mu_t,\nu_t)\leq\E\abs{X_t-Y_t}$ and $\E\abs{X_0-Y_0}=W_1(\mu_0,\nu_0)$.
\end{proof}



\begin{remark} \textbf{Replacing uniform field error by a trajectory
$L^2$ error}
A squared-loss learning problem does not usually control the supremum in
\eqref{eq:field-discrepancy}. Suppose that the discrepancy is evaluated
at the process $Y_t\sim\nu_t$.  Set
\[
 e_t^\eta(x):=\Phi_t[\eta](x)-\Psi_t[\eta](x),
 \qquad
 \varepsilon_{2,\nu}(t)
 :=\left(\int_{\R^d}\abs{e_t^{\nu_t}(x)}^2\,\nu_t(\dd x)\right)^{1/2}.
\] 
In the preceding proof, Cauchy--Schwarz replaces the uniform
estimate by
\[
 \E\abs{e_t^{\nu_t}(Y_t)}\leq\varepsilon_{2,\nu}(t),
\]
and the time-dependent Gronwall inequality gives
\begin{equation}\label{eq:stability-occupation-L2}
 W_1(\mu_t,\nu_t)
 \leq e^{Lt}W_1(\mu_0,\nu_0)
 +\int_0^t e^{L(t-s)}\varepsilon_{2,\nu}(s)\,\dd s.
\end{equation}
The natural space--time norm is induced by the occupation measure
\[
 \Lambda_{\nu,T}(\dd s,\dd x):=\nu_s(\dd x)\,\dd s,
 \qquad
 \norm{e}_{L^2(\Lambda_{\nu,T})}^2
 :=\int_0^T\int_{\R^d}\abs{e_s^{\nu_s}(x)}^2
 \nu_s(\dd x)\,\dd s.
\]
A second application of Cauchy--Schwarz to
\eqref{eq:stability-occupation-L2} yields, for $t\leq T$,
\begin{equation}\label{eq:stability-space-time-L2}
 W_1(\mu_t,\nu_t)
 \leq e^{Lt}W_1(\mu_0,\nu_0)
 +\left(\frac{e^{2Lt}-1}{2L}\right)^{1/2}
  \norm{e}_{L^2(\Lambda_{\nu,t})},
\end{equation}
where the factor is interpreted as $\sqrt t$ when $L=0$.  If
$\mu_0,\nu_0\in\cP_2(\R^d)$, the same formula holds with $W_2$ in place of
$W_1$: starting from a $W_2$-optimal coupling and setting
$R_t=(\E\abs{X_t-Y_t}^2)^{1/2}$ gives
$R_t'\leq LR_t+\varepsilon_{2,\nu}(t)$.

Statistical learning usually supplies an $L^2$ error under a training design
$\rho_t^{\mathrm{tr}}$, not automatically under the rollout law.  If
$\nu_t\ll\rho_t^{\mathrm{tr}}$, then
\[
 \varepsilon_{2,\nu}(t)
 \leq  \left\|\frac{\dd\nu_t}{\dd\rho_t^{\mathrm{tr}}}\right\|_{L^\infty} ^{1/2}
 \norm{e_t^{\nu_t}}_{L^2(\rho_t^{\mathrm{tr}})}.
\]
This is the required state-space coverage, or distribution-shift, condition.
If training data follow $\mu_t$ rather than $\nu_t$, one can obtain the
analogous bound with the occupation measure of $\mu$ by splitting the fields
at $\Psi_t[\mu_t](X_t)$ and assuming the Lipschitz estimate for $\Psi$.

For uniform over a class $\mathcal A_t$ of all measures that can be visited, one may instead assume
\[
 \sup_{\eta\in\mathcal A_t}
 \norm{e_t^\eta}_{L^2(\eta)}\leq\varepsilon_2(t), 
\]
 then proceed as above with $\varepsilon_{2,\nu}(t)\leq\varepsilon_2(t)$.
\end{remark}

\paragraph{\textbf{The finite-depth version}}
For practical networks it is often better not to pass to continuous depth.
For step sizes $h_\ell\geq0$, let
\[
 \mu^{\ell+1}
  =(\Id+h_\ell\Phi_\ell[\mu^\ell])_\#\mu^\ell,\qquad
 \nu^{\ell+1}
  =(\Id+h_\ell\Psi_\ell[\nu^\ell])_\#\nu^\ell.
\]
Suppose \eqref{eq:field-measure-lip} holds at layer $\ell$ with
$L_\ell=L_{x,\ell}+L_{\mu,\ell}$ and the field error bound in \eqref{eq:field-discrepancy} holds with
$\varepsilon_\ell$.
Precisely, the latter means
\[
 \sup_{x,\eta}
 \abs{\Phi_\ell[\eta](x)-\Psi_\ell[\eta](x)}
 \leq\varepsilon_\ell
\]
on the region visited by the two layer sequences. Pushing a coupling through the two maps gives (see Exercise \ref{ex:layer-stability})
\begin{equation}\label{eq:layer-stability-recursion}
 d_{\ell+1}\leq(1+h_\ell L_\ell)d_\ell+h_\ell\varepsilon_\ell,
 \qquad d_\ell=W_1(\mu^\ell,\nu^\ell).
\end{equation}
Iteration yields the explicit error budget
\begin{equation}\label{eq:layer-stability-product}
\begin{split}
 d_L\leq{}&
 d_0\prod_{r=0}^{L-1}(1+h_rL_r) +\sum_{k=0}^{L-1}h_k\varepsilon_k
   \prod_{r=k+1}^{L-1}(1+h_rL_r).
\end{split}
\end{equation}
Thus an early-layer error is multiplied by every later amplification factor.
Equation~\eqref{eq:layer-stability-product} remains meaningful for untied
weights and order-one depth, where the ODE idealization may not.

\begin{exercise}[Derive the discrete error budget]\label{ex:layer-stability}
Prove \eqref{eq:layer-stability-recursion} by pushing a coupling through the two
layer maps, and then prove \eqref{eq:layer-stability-product} by induction.
\end{exercise}
\begin{proof}[Solution]
Let $\pi_\ell$ be an optimal coupling of $\mu^\ell$ and $\nu^\ell$, and let
$(X,Y)$ have law $\pi_\ell$.  Push this coupling through the two layer maps:
\[
 X^+=X+h_\ell\Phi_\ell[\mu^\ell](X),\qquad
 Y^+=Y+h_\ell\Psi_\ell[\nu^\ell](Y).
\]
The law of $(X^+,Y^+)$ is an admissible coupling of
$\mu^{\ell+1}$ and $\nu^{\ell+1}$.  Splitting the field difference at
$\Phi_\ell[\nu^\ell](Y)$ gives
\begin{align*}
 d_{\ell+1}
 &\leq\E\abs{X^+-Y^+}\\
 &\leq\E\abs{X-Y}
 +h_\ell\E\abs{
 \Phi_\ell[\mu^\ell](X)-\Psi_\ell[\nu^\ell](Y)}\\
 &\leq d_\ell+h_\ell\left(
 L_{x,\ell}d_\ell+L_{\mu,\ell}d_\ell
 +\varepsilon_\ell\right),
\end{align*}
which is \eqref{eq:layer-stability-recursion}.

Set $a_r=1+h_rL_r$.  We claim more generally that, for $1\leq m\leq L$,
\[
 d_m\leq d_0\prod_{r=0}^{m-1}a_r
 +\sum_{k=0}^{m-1}h_k\varepsilon_k
   \prod_{r=k+1}^{m-1}a_r,
\]
where an empty product equals one.  The case $m=1$ is the one-step recursion.
If the claim holds for $m$, multiplying it by $a_m$ and adding
$h_m\varepsilon_m$ proves it for $m+1$.  Taking $m=L$ gives
\eqref{eq:layer-stability-product}.
\end{proof}


%%%%%%%%%%%
\paragraph{\textbf{Contractivity from one-sided dissipation}}
Global Lipschitz estimates permit exponential growth even when the true
dynamics is contracting.  A one-sided condition captures dissipation.

\begin{proposition}[Contractive regime]\label{prop:contractive}
Suppose, for some $\kappa>L_\mu\geq0$,
\begin{equation}\label{eq:one-sided}
 \ip{x-y}{\Phi[\mu](x)-\Phi[\nu](y)}
 \leq-\kappa\abs{x-y}^2
      +L_\mu W_2(\mu,\nu)\abs{x-y}.
\end{equation}
Let $\mu_0,\nu_0\in\cP_2(\R^d)$.
If $\Phi$ is replaced by a field $\Psi$ with uniform error $\varepsilon$, then
\begin{equation}\label{eq:contractive-stability}
 W_2(\mu_t,\nu_t)
 \leq e^{-(\kappa-L_\mu)t}W_2(\mu_0,\nu_0)
 +\frac{1-e^{-(\kappa-L_\mu)t}}{\kappa-L_\mu}\varepsilon.
\end{equation}
The same estimate holds for diffusions driven by identical additive Brownian
noise.
\end{proposition}
\begin{proof}
Choose an optimal initial coupling and evolve it synchronously according to
\[
 \dot X_t=\Phi[\mu_t](X_t),\qquad
 \dot Y_t=\Psi[\nu_t](Y_t).
\]
Set $R(t)=(\E\abs{X_t-Y_t}^2)^{1/2}$ and $D_t=X_t-Y_t$.
At almost every time, insert and subtract $\Phi[\nu_t](Y_t)$ to obtain
\[
\begin{aligned}
 \frac12\frac{\dd}{\dd t}R(t)^2
 &=\E\left\langle D_t,
   \Phi[\mu_t](X_t)-\Psi[\nu_t](Y_t)\right\rangle\\
 &=\E\left\langle D_t,
   \Phi[\mu_t](X_t)-\Phi[\nu_t](Y_t)\right\rangle
   +\E\left\langle D_t,
   \Phi[\nu_t](Y_t)-\Psi[\nu_t](Y_t)\right\rangle\\
 &\leq-\kappa R(t)^2
   +L_\mu W_2(\mu_t,\nu_t)\E\abs{D_t}
   +\varepsilon\E\abs{D_t}\\
 &\leq-(\kappa-L_\mu)R(t)^2+\varepsilon R(t).
 \end{aligned}
\]
The first inequality uses \eqref{eq:one-sided} for the first
expectation and
$|\Phi[\nu_t](Y_t)-\Psi[\nu_t](Y_t)|\leq\varepsilon$ for the second.
The last inequality follows because the current joint law is a coupling, so
$W_2(\mu_t,\nu_t)\leq R(t)$, and Cauchy--Schwarz gives
$\E|D_t|\leq R(t)$.
Dividing by $R(t)$ when it is positive gives
\[
 R'(t)\leq-(\kappa-L_\mu)R(t)+\varepsilon.
\]
The same inequality holds in the upper-Dini sense at zeros of $R$; equivalently,
one may apply the calculation to $(R^2+\delta^2)^{1/2}$ and let
$\delta\downarrow0$.  Solving the scalar inequality yields
\[
 R(t)\leq e^{-(\kappa-L_\mu)t}R(0)
 +\frac{1-e^{-(\kappa-L_\mu)t}}{\kappa-L_\mu}\varepsilon.
\]
Finally $W_2(\mu_t,\nu_t)\leq R(t)$ and
$R(0)=W_2(\mu_0,\nu_0)$, proving \eqref{eq:contractive-stability}.
For additive diffusions driven by the same Brownian motion, the stochastic
increments cancel in $X_t-Y_t$, so the identical calculation applies.
\end{proof}
In the contractive regime, model error produces a bounded error floor rather
than exponential accumulation.  This is why monotonicity, spectral
normalization, or a provable Lyapunov structure can be more informative than a
small parameter norm.

\subsection{Finite particles and the high-dimensional obstruction}
For deterministic dynamics, the empirical measure of the $N$-particle system
is itself a weak solution of the nonlinear continuity equation.  Taking
$\nu_0=\mu_0^N$, $\Phi=\Psi$, and applying
Theorem~\ref{thm:w1-stability} gives
\begin{equation}\label{eq:empirical-from-initial}
 W_1(\mu_t,\mu_t^N)\leq e^{Lt}W_1(\mu_0,\mu_0^N).
\end{equation}
Thus concentration or quantization at the input propagates through finite
time.  For stochastic particles, the martingale in
\eqref{eq:empirical-weak} adds a dynamical fluctuation term.

There is no universal $N^{-1/2}$ rate for
$W_1(\mu_0,\mu_0^N)$ in high ambient dimension.  Without additional structure,
estimating a full $d$-dimensional distribution suffers a curse of
dimensionality.  Better rates require intrinsic low-dimensional support,
concentration of selected observables, a weaker metric, or special structure
of the attention field.  This is a central warning when importing classical
low-dimensional PDE estimates into embeddings with thousands of coordinates.


\begin{exercise}[Sampling dimension]
Set $\nu_0=\mu_0^N$ and $\Phi=\Psi$.  Explain how a bound on
$\E W_1(\mu_0,\mu_0^N)$ becomes a finite-particle approximation bound.  Compare
the cases where $\mu_0$ is supported on a curve and where it has a density on a
high-dimensional cube.
\end{exercise}
\begin{proof}[Solution]
Equation~\eqref{eq:empirical-from-initial} holds for every realization of the
initial sample.  Taking expectations therefore gives
\begin{equation}\label{eq:expected-particle-from-initial}
 \E W_1(\mu_t,\mu_t^N)
 \leq e^{Lt}\E W_1(\mu_0,\mu_0^N).
\end{equation}
Thus the dynamics propagates the statistical approximation error already
present at the input; it does not create a dimension-free sampling rate.

For a concrete intrinsic one-dimensional model, suppose
$\mu_0=T_\#\lambda$, where $T:[0,1]\to\R^d$ is Lipschitz and
$U_1,\ldots,U_N$ are i.i.d.\ with law $\lambda$.  With
$\lambda^N=N^{-1}\sum_i\delta_{U_i}$ and
$\mu_0^N=T_\#\lambda^N$, the contraction of Wasserstein distance under a
Lipschitz map gives
\[
 W_1(\mu_0,\mu_0^N)
 \leq\Lip(T)W_1(\lambda,\lambda^N).
\]
For a bounded one-dimensional law,
$\E W_1(\lambda,\lambda^N)\leq C N^{-1/2}$, so
\eqref{eq:expected-particle-from-initial} retains the parametric exponent.

By contrast, for a distribution with a density bounded above and below on a
$d$-dimensional cube, $d\geq3$, Theorem~1 of
\cite{FournierGuillin2015} gives
$\E W_1(\mu_0,\mu_0^N)\leq C N^{-1/d}$.  This exponent is unavoidable for
such a nondegenerate density.  Indeed, if $S$ is the support of any measure on
at most $N$ points, then
\[
 W_1(\mu_0,\nu)\geq\int\operatorname{dist}(x,S)\,\mu_0(\dd x).
\]
For a density bounded below, the volume bound
$\abs{\{x:\operatorname{dist}(x,S)\leq r\}}\leq Nv_dr^d$ and the choice
$r=(2Nv_d)^{-1/d}$ give a lower bound $cN^{-1/d}$.  Substitution in
\eqref{eq:expected-particle-from-initial} therefore gives the upper bound
$Ce^{Lt}N^{-1/d}$, with a dimension-dependent sampling exponent already
present at time zero.  The comparison is between intrinsic dimension one and
genuinely $d$-dimensional mass, not merely between two subsets embedded in the
same ambient space.
\end{proof}



\subsection{Miscellaneous remarks}


\paragraph{\textbf{From kernel error to rollout error}}
If an interaction has the form
\[
 b_\phi[\mu](x)=\int K_\phi(x,y)\,\mu(\dd y),
\]
then on a visited domain
\[
 \sup_{x,\mu}\abs{b_{\widehat\phi}[\mu](x)-b_\phi[\mu](x)}
 \leq\sup_{x,y}\abs{K_{\widehat\phi}(x,y)-K_\phi(x,y)}.
\]
The right side is a field-error $\varepsilon$ in
\eqref{eq:stability}.  More refined bounds replace the supremum by a
data-weighted norm.  This is the forward half of the inverse problem:
estimating the interaction in a sufficiently strong norm controls finite-time
rollout.  The converse is false in general; similar trajectories can hide
different kernels.

\paragraph{\textbf{Literature and storyline.}}
Dobrushin's stability estimate for Vlasov equations was a foundational shift:
instead of comparing two $N$-particle configurations coordinate by coordinate,
one compares their laws in a transportation metric
\cite{Dobrushin1979}.  That estimate made stability of the nonlinear limit the
engine behind mean-field convergence.  The later development of optimal
transport clarified why this metric is adapted to kinetic evolution.  The
official Fields Medal profile of Villani describes a research program linking
kinetic theory, entropy, functional inequalities, and optimal transport
\cite{VillaniFieldsProfile2010}; these are precisely the ingredients behind
long-time stability and relaxation.  The resulting proof strategy is the one
used here: construct a coupling, differentiate its transportation cost, split
state error from model error, and close the estimate by Gronwall.  It also
gives a useful statistical principle for the rest of the notes: establish the
forward stability modulus first, and only then insert the sampling or
estimation rate.  A small one-step loss has scientific meaning only through
the metric and time horizon in which the learned dynamics remain stable.


\subsection{Connections forward}
Generative samplers need precisely this decomposition: initial-law mismatch,
noise mismatch,
learned score or velocity error, and numerical integration error.  In operator
learning, \eqref{eq:layer-stability-product} converts one-step operator error
into rollout error.  In inverse problems, stability of the forward evolution
must be combined with a coercivity inequality in the reverse direction.  For
architecture design, the theorem motivates controlling the induced Jacobian
and measure sensitivity, while Proposition~\ref{prop:contractive} explains why
dissipative or normalized blocks can remain stable at greater depth.




\newpage





%%%%%%%%%%%%%%%%%%%%%%% 

\appendix
\section{ODE and SDE preliminaries}
\label{sec:ode-sde-preliminaries}
This appendix collects the deterministic and stochastic tools used repeatedly
in the notes.  Its purpose is also to distinguish three logically different
steps: a pathwise integral inequality, a maximal estimate for a martingale,
and a deterministic Gronwall argument after taking expectations.  Unless
stated otherwise, all processes are defined on a filtered probability space
$(\Omega,\mathcal F,(\mathcal F_t)_{t\geq0},\Prob)$ satisfying the usual
conditions.

We first record the elementary bridge between the pushforwards notation  and
random variables.

\begin{lemma}[Pushforwards, composition, and laws]
\label{lem:pushforward-random-variable}
Let $E,F,G$ be Polish spaces, let $\mu\in\cP(E)$, and let
$S:E\to F$ and $R:F\to G$ be Borel maps.  Then
\begin{align}
 \int_F h(y)\,(S_\#\mu)(\dd y)
 &=\int_E h(S(x))\,\mu(\dd x),
 \label{eq:pushforward-integration}\\
 R_\#(S_\#\mu)&=(R\circ S)_\#\mu.
 \label{eq:pushforward-composition}
\end{align}
The first identity holds for every bounded Borel $h$, and more generally
whenever either side is well defined.  If $X$ is an $E$-valued random
variable with $\Law(X)=\mu$, then
\begin{equation}\label{eq:law-of-map}
 \Law(S(X))=S_\#\mu.
\end{equation}
In particular, if $T:E\to F$ is Borel, then
$(\Id,T)_\#\mu$ is the joint law of $(X,T(X))$; its marginals are
$\mu$ and $T_\#\mu$.
\end{lemma}
\begin{proof}
By definition,
$(S_\#\mu)(B)=\mu(S^{-1}(B))$ for every Borel set $B\subset F$.
This proves \eqref{eq:pushforward-integration} first for indicator functions,
then for nonnegative and integrable Borel functions by the usual simple-function
approximation.  For a Borel set $C\subset G$,
\[
 [R_\#(S_\#\mu)](C)
 =(S_\#\mu)(R^{-1}(C))
 =\mu\bigl((R\circ S)^{-1}(C)\bigr),
\]
which proves \eqref{eq:pushforward-composition}.  Finally,
\[
 \Prob(S(X)\in B)=\Prob(X\in S^{-1}(B))
 =\mu(S^{-1}(B))=(S_\#\mu)(B),
\]
so \eqref{eq:law-of-map} follows.  Applying it to the map
$x\mapsto(x,T(x))$ gives the last assertion.
\end{proof}

The following lemma is used in Section~\ref{sec:ips} to justify the convergence of empirical measures
in Wasserstein distance, where $E=\R^d$, $p=2$, and $Y_i=\overline X_t^i$.

\begin{lemma}[Strong law for empirical measures in Wasserstein distance]
\label{lem:empirical-wasserstein-slln}
Let $(E,d)$ be a Polish space (i.e., a complete separable metric space), let $p\geq1$, and let $Y_1,Y_2,\ldots$ be i.i.d.\ random variables with law $\mu$.  Assume that,
for some $y_0\in E$,
$\int_E d(y,y_0)^p\,\mu(\dd y)<\infty$.  For
\[
 \mu^N=\frac1N\sum_{i=1}^N\delta_{Y_i},
\]
one has
\begin{equation}\label{eq:empirical-wasserstein-slln}
 W_p(\mu^N,\mu)\longrightarrow0
 \qquad\text{almost surely}.
\end{equation}
\end{lemma}
\begin{proof}
Fix $y_0\in E$.  Since $E$ is Polish, there is a countable
convergence-determining family
$\{f_m\}_{m\geq1}\subset C_b(E)$, i.e., if $\int f_m d\mu^N \to \int f_m d\mu$ for all $m$, then $\mu^N \Rightarrow \mu$.  For each $m$, the scalar strong law of large numbers gives
\[
 \int_E f_m\,\dd\mu^N
 =\frac1N\sum_{i=1}^N f_m(Y_i)
 \longrightarrow \E f_m(Y_1)
 =\int_E f_m\,\dd\mu
 \qquad\text{almost surely}.
\]
For each $m$, let $A_m$ be the event that $\int_E f_m\,\dd\mu^N \to \int_E f_m\,\dd\mu$, then $\Prob(A_m) = 1$ for each $m$. Intersecting these probability-one events over $m$, we get a probability-one event $A:=\bigcap_{m=1}^\infty A_m$, and on the event $A$, we have that $\int_E f_m\,\dd\mu^N \to \int_E f_m\,\dd\mu$ for all $m$. This implies that $\mu^N \Rightarrow \mu$ almost surely.

The moment assumption also permits the
strong law of large numbers for the integrable random variable $d(Y_i,y_0)^p$:
\[
 \int_E d(y,y_0)^p\,\mu^N(\dd y)
 =\frac1N\sum_{i=1}^N d(Y_i,y_0)^p
 \longrightarrow
 \int_E d(y,y_0)^p\,\mu(\dd y)
 \qquad\text{almost surely}.
\]
Weak convergence together with convergence of these $p$th moments is
equivalent to $W_p$ convergence by
\cite[Theorem~6.9]{Villani2009}, proving
\eqref{eq:empirical-wasserstein-slln}.  
\end{proof}


\subsection{Fixed points and deterministic Gronwall inequalities}

The fixed-point construction in Section~\ref{sec:ips} uses the following theorem.

\begin{theorem}[Banach contraction principle]
\label{thm:appendix-banach}
Let $(E,d)$ be a complete metric space and let $\Phi:E\to E$ satisfy
\[
 d(\Phi(x),\Phi(y))\leq qd(x,y),
 \qquad x,y\in E,\qquad 0\leq q<1.
\]
Then $\Phi$ has a unique fixed point $x_\ast$.  For every $x_0\in E$, the
iterates $x_{n+1}=\Phi(x_n)$ satisfy
\begin{equation}\label{eq:appendix-banach-rate}
 d(x_n,x_\ast)\leq\frac{q^n}{1-q}d(x_1,x_0).
\end{equation}
\end{theorem}
The proof sums the geometric bound
$d(x_{n+1},x_n)\leq q^nd(x_1,x_0)$ to show that $(x_n)$ is Cauchy, and then
uses continuity of $\Phi$ and completeness of $E$.  See
\cite[Theorem~2.1]{Teschl2012}.

The next form of Gronwall's inequality contains all continuous-time versions
used in these notes.

\begin{lemma}[Generalized integral Gronwall inequality]
\label{lem:appendix-gronwall}
Let $u,\alpha:[0,T]\to[0,\infty)$ be continuous and let
$\beta\in L^1(0,T)$ be nonnegative.  If
\begin{equation}\label{eq:appendix-gronwall-assumption}
 u(t)\leq\alpha(t)+\int_0^t\beta(s)u(s)\dd s,
 \qquad 0\leq t\leq T,
\end{equation}
then
\begin{equation}\label{eq:appendix-gronwall-general}
 u(t)\leq\alpha(t)+\int_0^t\alpha(s)\beta(s)
 \exp\left(\int_s^t\beta(r)\dd r\right)\dd s.
\end{equation}
If $\alpha$ is nondecreasing, this simplifies to
\begin{equation}\label{eq:appendix-gronwall-monotone}
 u(t)\leq\alpha(t)
 \exp\left(\int_0^t\beta(s)\dd s\right).
\end{equation}
\end{lemma}
\begin{proof}
Set
$v(t)=\int_0^t\beta(s)u(s)\dd s$ and
$B(t)=\int_0^t\beta(s)\dd s$.  Then $v$ is absolutely continuous and
\[
 v'(t)=\beta(t)u(t)
 \leq\beta(t)\alpha(t)+\beta(t)v(t)
 \qquad\text{for almost every }t.
\]
Multiplying by $e^{-B(t)}$ and integrating gives
\[
 v(t)\leq\int_0^t\alpha(s)\beta(s)e^{B(t)-B(s)}\dd s.
\]
Combining this with $u\leq\alpha+v$ proves
\eqref{eq:appendix-gronwall-general}.  If $\alpha$ is nondecreasing, replace
$\alpha(s)$ in the integral by $\alpha(t)$ and use
\[
 1+\int_0^t\beta(s)e^{B(t)-B(s)}\dd s=e^{B(t)}.
\]
\end{proof}

In particular, if
\begin{equation}\label{eq:appendix-course-gronwall}
 u(t)\leq a+c\int_0^tu(s)\dd s,
\end{equation}
then $u(t)\leq ae^{ct}$.  Thus an error $a=C/N$ remains of order $N^{-1}$
on every fixed time interval, although its constant can grow exponentially in
the horizon.  This is the version used in the propagation-of-chaos proof.  A
precise reference for the generalized form is
\cite[Lemma~2.7]{Teschl2012}.

For dissipative estimates, the integrating-factor form is often more useful.
If $u$ is absolutely continuous and
\[
 u'(t)\leq-\lambda u(t)+f(t)
 \quad\text{a.e.},\qquad \lambda>0,\quad f\geq0,
\]
then
\begin{equation}\label{eq:appendix-dissipative-gronwall}
 u(t)\leq e^{-\lambda t}u(0)
 +\int_0^te^{-\lambda(t-s)}f(s)\dd s.
\end{equation}
In particular, $f=0$ yields the exponential diameter and variance decay used
in Section~\ref{sec:clustering}.

\begin{lemma}[Discrete Gronwall inequality]
\label{lem:appendix-discrete-gronwall}
Let $h>0$, $L\geq0$, and suppose that nonnegative sequences $(e_n)$ and
$(r_n)$ satisfy
\[
 e_{n+1}\leq(1+hL)e_n+hr_n,
 \qquad n=0,\ldots,M-1.
\]
Then
\begin{align}
 e_n
 &\leq(1+hL)^ne_0
 +h\sum_{j=0}^{n-1}(1+hL)^{n-1-j}r_j\notag\\
 &\leq e^{Lt_n}\left(e_0+h\sum_{j=0}^{n-1}r_j\right),
 \qquad t_n=nh.\label{eq:appendix-discrete-gronwall}
\end{align}
Consequently, if $e_0=0$ and the local defect is $Ch^2$, so that
$r_n=Ch$, then $\max_{t_n\leq T}e_n\leq C_T h$.
\end{lemma}
\begin{proof}
Iterate the one-step recursion to obtain the first line of
\eqref{eq:appendix-discrete-gronwall}; the second follows from
$(1+hL)^k\leq e^{Lkh}$.  This is the estimate behind the continuous-depth
Euler bound in Section~\ref{sec:continuous-depth}; see also
\cite[Section~2.7]{Teschl2012}.
\end{proof}

\subsection{ODE well-posedness, stability, and characteristics}

\begin{theorem}[Global Picard--Lindel\"of theorem]
\label{thm:appendix-ode-wp}
Let $F:[0,T]\times\R^m\to\R^m$ be continuous and globally Lipschitz in its
second variable, uniformly in time:
\[
 \abs{F(t,x)-F(t,y)}\leq L\abs{x-y}.
\]
For every $x_0\in\R^m$, the initial-value problem
\begin{equation}\label{eq:appendix-ode}
 \dot x_t=F(t,x_t),\qquad x_{t=0}=x_0,
\end{equation}
has a unique solution on $[0,T]$.  If
$K=\sup_{0\leq t\leq T}\abs{F(t,0)}$, then
\begin{equation}\label{eq:appendix-ode-growth}
 \abs{x_t}\leq e^{Lt}\abs{x_0}
 +K\int_0^te^{L(t-s)}\dd s.
\end{equation}
Moreover, if $y$ solves $\dot y_t=G(t,y_t)$ and
\[
 \abs{F(t,z)-G(t,z)}\leq\varepsilon(t),
 \qquad z\in\R^m,\quad \varepsilon\in L^1(0,T),
\]
then
\begin{equation}\label{eq:appendix-ode-stability}
 \abs{x_t-y_t}
 \leq e^{Lt}\abs{x_0-y_0}
 +\int_0^te^{L(t-s)}\varepsilon(s)\dd s.
\end{equation}
\end{theorem}
\begin{proof}
Local existence and uniqueness follow by applying
Theorem~\ref{thm:appendix-banach} to the integral map
\[
 \Phi(x)(t)=x_0+\int_0^tF(s,x(s))\dd s
\]
on a sufficiently short time interval.  The linear-growth estimate
$\abs{F(t,x)}\leq K+L\abs{x}$ and
Lemma~\ref{lem:appendix-gronwall} prevent finite-time blowup, so the solution
extends to $[0,T]$.  Applying the same lemma to the integral equations for
$x$ and $y$ gives \eqref{eq:appendix-ode-growth} and
\eqref{eq:appendix-ode-stability}.  See
\cite[Theorem~2.2 and Corollary~2.6]{Teschl2012} for the local and global
statements.
\end{proof}

The empirical-measure identities in
Sections~\ref{sec:self-attention}--\ref{sec:clustering} are ordinary chain rules
applied simultaneously to all characteristics.

\begin{proposition}[Characteristics and the continuity equation]
\label{prop:appendix-characteristics}
Let $v:[0,T]\times\R^m\to\R^m$ be continuous, globally Lipschitz in space,
and of linear growth.  Let $F_t(x)$ solve
$\dot F_t(x)=v(t,F_t(x))$, $F_0(x)=x$, and set
$\mu_t=(F_t)_\#\mu_0$ for $\mu_0\in\cP_1(\R^m)$.  Then, for every
$\varphi\in C_c^1(\R^m)$,
\begin{equation}\label{eq:appendix-continuity-weak}
 \int\varphi\,\dd\mu_t-\int\varphi\,\dd\mu_0
 =\int_0^t\int\nabla\varphi(x)\cdot v(s,x)\,\mu_s(\dd x)\dd s.
\end{equation}
Thus $\mu_t$ is a weak solution of
$\partial_t\mu+\diver(\mu v)=0$.
\end{proposition}
\begin{proof}
Use the definition of pushforward and the ordinary chain rule:
\[
 \varphi(F_t(x))-\varphi(x)
 =\int_0^t\nabla\varphi(F_s(x))\cdot v(s,F_s(x))\dd s.
\]
Integrate with respect to $\mu_0$ and change variables by
$\mu_s=(F_s)_\#\mu_0$.
\end{proof}

The consensus proof also uses a standard fact about nonsmooth extrema.  If
$z_1,\ldots,z_N$ are differentiable, define
$M(t)=\max_i z_i(t)$ and $m(t)=\min_i z_i(t)$.  Their upper and lower right
Dini derivatives satisfy
\[
 D^+M(t):=\limsup_{h\downarrow0}\frac{M(t+h)-M(t)}h,
 \qquad
 D_+m(t):=\liminf_{h\downarrow0}\frac{m(t+h)-m(t)}h.
\]
\begin{equation}\label{eq:appendix-dini-maximum}
 D^+M(t)=\max_{i:z_i(t)=M(t)}\dot z_i(t),
 \qquad
 D_+m(t)=\min_{i:z_i(t)=m(t)}\dot z_i(t).
\end{equation}
This follows by separating the active indices from the finitely many indices
with a strict gap at time $t$.  It permits a differential inequality for a
diameter even when an extremizing particle changes or several particles tie.

\subsection{Stochastic integrals, brackets, and maximal inequalities}

Let $B$ be a standard Brownian motion in $\R^q$.  A predictable process
$H:[0,T]\times\Omega\to\R^q$ is square integrable if
$\E\int_0^T\abs{H_s}^2\dd s<\infty$.

\begin{theorem}[It\^o isometry and bracket]
\label{thm:appendix-ito-isometry}
For every square-integrable predictable $H$, the process
\[
 M_t=\int_0^tH_s\cdot\dd B_s
\]
is a continuous square-integrable martingale with $M_0=0$ and
\begin{equation}\label{eq:appendix-ito-isometry}
 \E M_t=0,\qquad
 \E\abs{M_t}^2=\E\int_0^t\abs{H_s}^2\dd s,\qquad
 \langle M\rangle_t=\int_0^t\abs{H_s}^2\dd s.
\end{equation}
If $M^i=\int H_s^i\cdot\dd B_s^i$ are driven by independent Brownian
motions, then $\langle M^i,M^j\rangle=0$ for $i\ne j$, and the bracket of a
finite sum is the sum of the individual brackets.
\end{theorem}
See \cite[Chapter~3, Section~2]{KaratzasShreve1991} for the construction of
the stochastic integral, its isometry, and quadratic variation.

\begin{theorem}[Doob's maximal inequality]
\label{thm:appendix-doob}
Let $M$ be a martingale with right-continuous paths and let $p>1$.  If
$M_T\in L^p$, then
\begin{equation}\label{eq:appendix-doob}
 \E\left[\sup_{0\leq t\leq T}\abs{M_t}^p\right]
 \leq\left(\frac{p}{p-1}\right)^p\E\abs{M_T}^p.
\end{equation}
For $p=2$, the constant is $4$.  Since $\abs{B_t}$ is a nonnegative
submartingale and $\E\abs{B_T}^2=qT$,
\begin{equation}\label{eq:appendix-brownian-max}
 \E\sup_{0\leq t\leq T}\abs{B_t}^2\leq4qT.
\end{equation}
\end{theorem}
The submartingale and continuous-time versions are given in
\cite[Chapter~II, Section~1]{RevuzYor1999}; see also
\cite[Chapter~1, Section~3]{KaratzasShreve1991}.

\begin{theorem}[Burkholder--Davis--Gundy inequality]
\label{thm:appendix-bdg}
Let $M$ be a real-valued continuous local martingale with $M_0=0$.  For every
$p>0$, there are constants $0<c_p\leq C_p<\infty$, depending only on $p$,
such that for every stopping time $\tau$,
\begin{equation}\label{eq:appendix-bdg}
 c_p\E\langle M\rangle_\tau^{p/2}
 \leq\E\left[\sup_{0\leq t\leq\tau}\abs{M_t}^p\right]
 \leq C_p\E\langle M\rangle_\tau^{p/2}.
\end{equation}
The quantities may be infinite.  In particular, for
$M_t=\int_0^tH_s\cdot\dd B_s$, the right side is controlled by
$\E(\int_0^\tau\abs{H_s}^2\dd s)^{p/2}$.
\end{theorem}
A precise reference is \cite[Chapter~IV, Theorem~4.1]{RevuzYor1999}.

Doob and BDG answer related but different questions.  Doob controls the
running maximum by an integrable terminal value.  BDG identifies the natural
scale of a continuous martingale directly through its bracket.  At $p=2$,
Doob's inequality and the It\^o isometry give the useful upper bound
\begin{equation}\label{eq:appendix-doob-isometry}
 \E\sup_{0\leq t\leq T}\abs{M_t}^2
 \leq4\E\langle M\rangle_T.
\end{equation}

For martingale terms produced by It\^o's formula, BDG is commonly combined
with Young's inequality.  For example, if
$M_t=\int_0^tY_sH_s\cdot\dd B_s$, then
\begin{align}
 \E\sup_{0\leq r\leq t}\abs{M_r}
 &\leq C_1\E\left[
 \sup_{0\leq r\leq t}\abs{Y_r}
 \left(\int_0^t\abs{H_s}^2\dd s\right)^{1/2}\right]\notag\\
 &\leq\varepsilon\E\sup_{0\leq r\leq t}\abs{Y_r}^2
 +\frac{C_1^2}{4\varepsilon}
 \E\int_0^t\abs{H_s}^2\dd s,
 \qquad\varepsilon>0.\label{eq:appendix-bdg-young}
\end{align}
The first term can then be absorbed into the left side of a moment estimate.

\subsection{It\^o's formula and globally Lipschitz SDEs}

\begin{theorem}[Multidimensional It\^o formula]
\label{thm:appendix-ito-formula}
Let $X_t\in\R^m$ satisfy
\begin{equation}\label{eq:appendix-sde-general}
 \dd X_t=b_t\dd t+\Sigma_t\dd B_t,
\end{equation}
where $b_t\in\R^m$, $\Sigma_t\in\R^{m\times q}$, and the usual local
integrability assumptions hold.  For $f\in C^{1,2}([0,T]\times\R^m)$,
\begin{align}
 \dd f(t,X_t)
 &=\left[\partial_tf(t,X_t)+b_t\cdot\nabla f(t,X_t)
 +\frac12\operatorname{tr}\left(
 \Sigma_t\Sigma_t^\top D^2f(t,X_t)\right)\right]\dd t\notag\\
 &\quad+\nabla f(t,X_t)^\top\Sigma_t\dd B_t.
 \label{eq:appendix-ito-formula}
\end{align}
\end{theorem}
See \cite[Chapter~3, Section~3]{KaratzasShreve1991}.  Taking expectations in
\eqref{eq:appendix-ito-formula} removes the stochastic integral when it is a
true martingale.  This is the step that converts an SDE into its weak
Fokker--Planck equation.  For constant
$\Sigma=\sigma I_m$, the second-order term is
$(\sigma^2/2)\Delta f$.

\begin{theorem}[Strong well-posedness and a second-moment estimate]
\label{thm:appendix-sde-wp}
Let
$b:[0,T]\times\R^m\to\R^m$ and
$\Sigma:[0,T]\times\R^m\to\R^{m\times q}$ be measurable and suppose that,
for some $L,K<\infty$,
\begin{align}
 \abs{b(t,x)-b(t,y)}
 +\norm{\Sigma(t,x)-\Sigma(t,y)}_{\mathrm F}
 &\leq L\abs{x-y},\label{eq:appendix-sde-lipschitz}\\
 \abs{b(t,0)}^2+\norm{\Sigma(t,0)}_{\mathrm F}^2
 &\leq K.\label{eq:appendix-sde-origin}
\end{align}
If $X_0$ is $\mathcal F_0$-measurable with $\E\abs{X_0}^2<\infty$ and is
independent of $B$, then
\begin{equation}\label{eq:appendix-lipschitz-sde}
 \dd X_t=b(t,X_t)\dd t+\Sigma(t,X_t)\dd B_t,
 \qquad X_{t=0}=X_0,
\end{equation}
has a unique global strong solution.  Moreover,
\begin{equation}\label{eq:appendix-sde-moment}
 \E\sup_{0\leq t\leq T}\abs{X_t}^2
 \leq C_{T,L,K,m,q}\left(1+\E\abs{X_0}^2\right).
\end{equation}
\end{theorem}
\begin{proof}
Existence and pathwise uniqueness follow by Picard iteration in the space of
adapted continuous processes; see
\cite[Chapter~5, Section~2]{KaratzasShreve1991}.  We record the moment argument
because it is used repeatedly.  Conditions
\eqref{eq:appendix-sde-lipschitz}--\eqref{eq:appendix-sde-origin} imply
\[
 \abs{b(t,x)}^2+\norm{\Sigma(t,x)}_{\mathrm F}^2
 \leq C(1+\abs{x}^2).
\]
Writing \eqref{eq:appendix-lipschitz-sde} in integral form, using
$\abs{x+y+z}^2\leq3(\abs{x}^2+\abs{y}^2+\abs{z}^2)$, and applying
Cauchy--Schwarz in time gives, for
$Q(t)=\E\sup_{0\leq r\leq t}\abs{X_r}^2$,
\begin{align*}
 Q(t)
 &\leq3\E\abs{X_0}^2
 +3t\int_0^t\E\abs{b(s,X_s)}^2\dd s\\
 &\quad+3\E\sup_{0\leq r\leq t}
 \left\lvert\int_0^r\Sigma(s,X_s)\dd B_s\right\rvert^2.
\end{align*}
Apply Doob's $L^2$ inequality componentwise, or its Hilbert-space version,
and then the It\^o isometry:
\[
 \E\sup_{0\leq r\leq t}
 \left\lvert\int_0^r\Sigma(s,X_s)\dd B_s\right\rvert^2
 \leq4\E\int_0^t\norm{\Sigma(s,X_s)}_{\mathrm F}^2\dd s.
\]
The linear-growth bound therefore yields
\[
 Q(t)\leq C_T\left(1+\E\abs{X_0}^2\right)
 +C_T\int_0^tQ(s)\dd s.
\]
Lemma~\ref{lem:appendix-gronwall} proves
\eqref{eq:appendix-sde-moment}.
\end{proof}

If two solutions use the same Brownian motion and have the same
state-independent diffusion matrix, their noise terms cancel upon
subtraction.  Their difference then satisfies an ordinary random integral
equation, as in the synchronous coupling of Section~\ref{sec:ips}.  For state-dependent
diffusion, a stochastic integral remains; Doob or BDG must be applied before
Gronwall.

\subsection{The estimate pipeline used in the notes}
The recurring argument can be summarized as follows.
\begin{enumerate}[leftmargin=*]
 \item Write the ODE or SDE in integral form and isolate the difference or
 moment to be estimated.
 \item Take a time supremum.  Use Cauchy--Schwarz for drift integrals and
 Doob or BDG for stochastic integrals.
 \item Take expectations and use linear growth, Lipschitz continuity,
 independence, or a coupling estimate to close the right side.
 \item Apply the deterministic integral Gronwall inequality to the resulting
 scalar function of time.
\end{enumerate}
Thus the arguments in these notes do not invoke a separate stochastic
Gronwall theorem: the stochastic term is controlled first, and the final
Gronwall step is deterministic.  For reference, the ODE results above are
collected in \cite[Chapter~2]{Teschl2012}; stochastic integration and It\^o's
formula are in \cite[Chapter~3]{KaratzasShreve1991}; strong SDE
well-posedness is in \cite[Chapter~5, Section~2]{KaratzasShreve1991}; and the
Doob and BDG inequalities are in
\cite[Chapter~II, Section~1 and Chapter~IV, Theorem~4.1]{RevuzYor1999}.

\paragraph{\textbf{Literature and storyline.}}
The estimates in this appendix also have a useful human history.  Gronwall
worked across engineering and analysis; the biographical account of his work
explains how the inequality that now bears his name emerged from differential
equations rather than from an abstract lemma factory
\cite{GronwallMacTutor}.  It\^o's Kyoto Prize profile dates his construction of
stochastic integration to 1942 and emphasizes how it supplied a calculus for
random paths that ordinary differentiation cannot handle
\cite{ItoKyotoPrize1998}.  Doob then built martingale theory into a systematic
part of probability; his National Academy memoir records the development of
that viewpoint and its influence on modern stochastic processes
\cite{DoobNASMemoir2016}.  Their results form the repeated proof pipeline of
these notes: deterministic growth is closed by Gronwall, stochastic maxima by
Doob or BDG, and nonlinear dynamics by applying those inequalities before a
final deterministic comparison.  Remembering the problems these tools were
created to solve makes them easier to use correctly than treating them as a
list of interchangeable black boxes.

\begin{thebibliography}{99}

\bibitem{AmbrosioGigliSavare2008}
L.~Ambrosio, N.~Gigli, and S.~Savar\'e,
\emph{Gradient Flows in Metric Spaces and in the Space of Probability
Measures},
2nd ed., Lectures in Mathematics ETH Z\"urich, Birkh\"auser, 2008.
\href{https://doi.org/10.1007/978-3-7643-8722-8}
{doi:10.1007/978-3-7643-8722-8}.

\bibitem{BenamouBrenier2000}
J.-D.~Benamou and Y.~Brenier,
\emph{A computational fluid mechanics solution to the Monge--Kantorovich
mass transfer problem},
Numer. Math. \textbf{84} (2000), no.~3, 375--393.
\href{https://doi.org/10.1007/s002110050002}
{doi:10.1007/s002110050002}.

\bibitem{Cavalier2008}
L.~Cavalier,
\emph{Nonparametric statistical inverse problems},
Inverse Problems \textbf{24} (2008), no.~3, 034004.
\href{https://doi.org/10.1088/0266-5611/24/3/034004}
{doi:10.1088/0266-5611/24/3/034004}.

\bibitem{BreschJabinWang2023}
D.~Bresch, P.-E.~Jabin, and Z.~Wang,
\emph{Mean field limit and quantitative estimates with singular attractive
kernels},
Duke Math. J. \textbf{172} (2023), no.~13, 2591--2641.
\href{https://doi.org/10.1215/00127094-2022-0088}
{doi:10.1215/00127094-2022-0088}.

\bibitem{CarmonaDelarue2018}
R.~Carmona and F.~Delarue,
\emph{Probabilistic Theory of Mean Field Games with Applications I:
Mean Field FBSDEs, Control, and Games},
Probability Theory and Stochastic Modelling, vol.~83, Springer, 2018.
\href{https://doi.org/10.1007/978-3-319-58920-6}
{doi:10.1007/978-3-319-58920-6}.

\bibitem{ChodronRosenzweigSerfaty2025}
A.~Chodron de Courcel, M.~Rosenzweig, and S.~Serfaty,
\emph{Sharp uniform-in-time mean-field convergence for singular periodic
Riesz flows},
Ann. Inst. H. Poincar\'e C Anal. Non Lin\'eaire
\textbf{42} (2025), no.~2, 391--472.
\href{https://doi.org/10.4171/AIHPC/105}{doi:10.4171/AIHPC/105}.

\bibitem{FournierGuillin2015}
N.~Fournier and A.~Guillin,
\emph{On the rate of convergence in Wasserstein distance of the empirical
measure},
Probab. Theory Related Fields \textbf{162} (2015), no.~3--4, 707--738.
\href{https://doi.org/10.1007/s00440-014-0583-7}
{doi:10.1007/s00440-014-0583-7}.

\bibitem{HaoRocknerZhang2024}
Z.~Hao, M.~R\"ockner, and X.~Zhang,
\emph{Strong convergence of propagation of chaos for McKean--Vlasov SDEs
with singular interactions},
SIAM J. Math. Anal. \textbf{56} (2024), no.~2, 2661--2713.
\href{https://doi.org/10.1137/23M1556666}{doi:10.1137/23M1556666}.

\bibitem{JabinWang2018}
P.-E.~Jabin and Z.~Wang,
\emph{Quantitative estimates of propagation of chaos for stochastic systems
with $\dot W^{-1,\infty}$ kernels},
Invent. Math. \textbf{214} (2018), no.~1, 523--591.
\href{https://doi.org/10.1007/s00222-018-0808-y}
{doi:10.1007/s00222-018-0808-y}.

\bibitem{JordanKinderlehrerOtto1998}
R.~Jordan, D.~Kinderlehrer, and F.~Otto,
\emph{The variational formulation of the Fokker--Planck equation},
SIAM J. Math. Anal. \textbf{29} (1998), no.~1, 1--17.
\href{https://doi.org/10.1137/S0036141096303359}
{doi:10.1137/S0036141096303359}.

\bibitem{KaratzasShreve1991}
I.~Karatzas and S.~E.~Shreve,
\emph{Brownian Motion and Stochastic Calculus},
Graduate Texts in Mathematics, vol.~113, 2nd ed., Springer, 1991.
\href{https://doi.org/10.1007/978-1-4612-0949-2}
{doi:10.1007/978-1-4612-0949-2}.

\bibitem{BrunnerEtAl2020Identifiability}
G.~Brunner, Y.~Liu, D.~Pascual, O.~Richter, M.~Ciaramita, and R.~Wattenhofer,
\emph{On identifiability in Transformers},
ICLR, 2020;
\href{https://arxiv.org/abs/1908.04211}{arXiv:1908.04211}.

\bibitem{ChenEtAl2026Quantitative}
S.~Chen, Z.~Lin, Y.~Polyanskiy, and P.~Rigollet,
\emph{Quantitative clustering in mean-field Transformer models},
\href{https://arxiv.org/abs/2504.14697}{arXiv:2504.14697}, revised 2026.

\bibitem{CriscitielloEtAl2026Synchronization}
C.~Criscitiello, Q.~Rebjock, A.~D.~McRae, and N.~Boumal,
\emph{Synchronization on circles and spheres with nonlinear interactions},
SIAM J. Appl. Dyn. Syst. \textbf{25} (2026), no.~2, 1207--1237;
\href{https://doi.org/10.1137/24M1699759}{doi:10.1137/24M1699759}.

\bibitem{EdelmanEtAl2022InductiveBias}
B.~L.~Edelman, S.~Goel, S.~Kakade, and C.~Zhang,
\emph{Inductive biases and variable creation in self-attention mechanisms},
Proc. 39th ICML, PMLR \textbf{162} (2022), 5793--5831.
\href{https://proceedings.mlr.press/v162/edelman22a.html}{PMLR version}.

\bibitem{GaoLangLu2025SelfTest}
Y.~Gao, Q.~Lang, and F.~Lu,
\emph{Self-test loss functions for learning weak-form operators and gradient
flows},
\href{https://arxiv.org/abs/2412.03506}{arXiv:2412.03506}, revised 2025.

\bibitem{GeshkovskiEtAl2023Clusters}
B.~Geshkovski, C.~Letrouit, Y.~Polyanskiy, and P.~Rigollet,
\emph{The emergence of clusters in self-attention dynamics},
Adv. Neural Inf. Process. Syst. \textbf{36} (2023), 57026--57037;
\href{https://arxiv.org/abs/2305.05465}{arXiv:2305.05465}.

\bibitem{GeshkovskiEtAl2024Metastability}
B.~Geshkovski, H.~Koubbi, Y.~Polyanskiy, and P.~Rigollet,
\emph{Dynamic metastability in the self-attention model},
\href{https://arxiv.org/abs/2410.06833}{arXiv:2410.06833}, 2024.

\bibitem{IldizEtAl2024Markov}
M.~E.~Ildiz, Y.~Huang, Y.~Li, A.~S.~Rawat, and S.~Oymak,
\emph{From self-attention to Markov models: Unveiling the dynamics of
generative Transformers},
\href{https://arxiv.org/abs/2402.13512}{arXiv:2402.13512}, 2024.

\bibitem{KaragodinEtAl2024Causal}
N.~Karagodin, Y.~Polyanskiy, and P.~Rigollet,
\emph{Clustering in causal attention masking},
NeurIPS, 2024;
\href{https://arxiv.org/abs/2411.04990}{arXiv:2411.04990}.

\bibitem{Rigollet2026MeanField}
P.~Rigollet,
\emph{The mean-field dynamics of Transformers},
to appear in the Proceedings of the International Congress of Mathematicians,
2026;
\href{https://arxiv.org/abs/2512.01868}{arXiv:2512.01868}.

\bibitem{LangLu2022MeanField}
Q.~Lang and F.~Lu,
\emph{Learning interaction kernels in mean-field equations of first-order
systems of interacting particles},
SIAM J. Sci. Comput. \textbf{44} (2022), no.~1, A260--A285.
\href{https://doi.org/10.1137/20M1377072}{doi:10.1137/20M1377072}.

\bibitem{LangLu2023Identifiability}
Q.~Lang and F.~Lu,
\emph{Identifiability of interaction kernels in mean-field equations of
interacting particles},
Found. Data Sci. \textbf{5} (2023), no.~4, 480--502.
\href{https://doi.org/10.3934/fods.2023007}{doi:10.3934/fods.2023007}.

\bibitem{LiEtAl2021Identifiability}
Z.~Li, F.~Lu, M.~Maggioni, S.~Tang, and C.~Zhang,
\emph{On the identifiability of interaction functions in systems of
interacting particles},
Stochastic Process. Appl. \textbf{132} (2021), 135--163.
\href{https://doi.org/10.1016/j.spa.2020.10.005}
{doi:10.1016/j.spa.2020.10.005}.

\bibitem{LMT2022Stochastic}
F.~Lu, M.~Maggioni, and S.~Tang,
\emph{Learning interaction kernels in stochastic systems of interacting
particles from multiple trajectories},
Found. Comput. Math. \textbf{22} (2022), 1013--1067.
\href{https://doi.org/10.1007/s10208-021-09521-z}
{doi:10.1007/s10208-021-09521-z}.

\bibitem{LuEtAl2019}
F.~Lu, M.~Zhong, S.~Tang, and M.~Maggioni,
\emph{Nonparametric inference of interaction laws in systems of agents from
trajectory data},
Proc. Natl. Acad. Sci. USA \textbf{116} (2019), no.~29, 14424--14433.
\href{https://doi.org/10.1073/pnas.1822012116}
{doi:10.1073/pnas.1822012116}.

\bibitem{WangSeroussiLu2025Minimax}
X.~Wang, I.~Seroussi, and F.~Lu,
\emph{Optimal minimax rate of learning nonlocal interaction kernels},
\href{https://arxiv.org/abs/2311.16852}{arXiv:2311.16852}, revised 2025.

\bibitem{WeiLu2026Unlabeled}
V.~Wei and F.~Lu,
\emph{Learning interacting particle systems from unlabeled data},
\href{https://arxiv.org/abs/2604.02581}{arXiv:2604.02581}, 2026.

\bibitem{ZuckerEtAl2026AttentionMinimax}
S.~Zucker, X.~Wang, F.~Lu, and I.~Seroussi,
\emph{Minimax rates for learning pairwise interactions in attention-style
models},
ICLR, 2026;
\href{https://arxiv.org/abs/2510.11789}{arXiv:2510.11789}.

\bibitem{RevuzYor1999}
D.~Revuz and M.~Yor,
\emph{Continuous Martingales and Brownian Motion},
Grundlehren der mathematischen Wissenschaften, vol.~293, 3rd ed.,
Springer, 1999.
\href{https://doi.org/10.1007/978-3-662-06400-9}
{doi:10.1007/978-3-662-06400-9}.

\bibitem{Sznitman1991}
A.-S.~Sznitman,
``Topics in propagation of chaos,'' in
\emph{\'Ecole d'\'Et\'e de Probabilit\'es de Saint-Flour XIX---1989},
Lecture Notes in Mathematics, vol.~1464, Springer, 1991, pp.~165--251.
\href{https://doi.org/10.1007/BFb0085169}{doi:10.1007/BFb0085169}.

\bibitem{Teschl2012}
G.~Teschl,
\emph{Ordinary Differential Equations and Dynamical Systems},
Graduate Studies in Mathematics, vol.~140,
American Mathematical Society, Providence, RI, 2012.
\href{https://www.mat.univie.ac.at/~gerald/ftp/book-ode/ode.pdf}
{Author's online edition}.

\bibitem{Villani2009}
C.~Villani,
\emph{Optimal Transport: Old and New},
Grundlehren der mathematischen Wissenschaften, vol.~338,
Springer, 2009.
\href{https://doi.org/10.1007/978-3-540-71050-9}
{doi:10.1007/978-3-540-71050-9}.

\bibitem{Wang2026}
S.~Wang,
\emph{Sharp local propagation of chaos for mean field particles with
$W^{-1,\infty}$ kernels},
J. Funct. Anal. \textbf{290} (2026), no.~3, Article~111240.
\href{https://doi.org/10.1016/j.jfa.2025.111240}
{doi:10.1016/j.jfa.2025.111240}.

\bibitem{Kac1956}
M.~Kac,
\emph{Foundations of kinetic theory},
in \emph{Proceedings of the Third Berkeley Symposium on Mathematical
Statistics and Probability}, vol.~III, University of California Press,
1956, pp.~171--197.
\href{https://digicoll.lib.berkeley.edu/record/112857}
{UC Berkeley Library}.

\bibitem{Kac1985Enigmas}
M.~Kac,
\emph{Enigmas of Chance: An Autobiography},
Harper \& Row, New York, 1985.

\bibitem{McKean1966}
H.~P.~McKean, Jr.,
\emph{A class of Markov processes associated with nonlinear parabolic
equations},
Proc. Natl. Acad. Sci. USA \textbf{56} (1966), no.~6, 1907--1911.
\href{https://doi.org/10.1073/pnas.56.6.1907}
{doi:10.1073/pnas.56.6.1907}.

\bibitem{ErdosSchleinYau2007NLS}
L.~Erd\H{o}s, B.~Schlein, and H.-T.~Yau,
\emph{Derivation of the cubic non-linear Schr\"odinger equation from quantum
dynamics of many-body systems},
Invent. Math. \textbf{167} (2007), no.~3, 515--614.
\href{https://doi.org/10.1007/s00222-006-0022-1}
{doi:10.1007/s00222-006-0022-1}.

\bibitem{DengHani2023WKE}
Y.~Deng and Z.~Hani,
\emph{Full derivation of the wave kinetic equation},
Invent. Math. \textbf{233} (2023), no.~2, 543--724.
\href{https://doi.org/10.1007/s00222-023-01189-2}
{doi:10.1007/s00222-023-01189-2}.

\bibitem{DengHani2026Chaos}
Y.~Deng and Z.~Hani,
\emph{Propagation of chaos and higher order statistics in wave kinetic theory},
J. Eur. Math. Soc. \textbf{28} (2026), no.~2, 673--733.
\href{https://doi.org/10.4171/JEMS/1488}{doi:10.4171/JEMS/1488}.

\bibitem{Kantorovich1942}
L.~V.~Kantorovich,
\emph{On the translocation of masses},
Management Sci. \textbf{5} (1958), no.~1, 1--4;
English translation of Dokl. Akad. Nauk SSSR \textbf{37} (1942), 199--201.
\href{https://doi.org/10.1287/mnsc.5.1.1}
{doi:10.1287/mnsc.5.1.1}.

\bibitem{NobelKantorovich1975}
The Royal Swedish Academy of Sciences,
\emph{Press release: The Sveriges Riksbank Prize in Economic Sciences in
Memory of Alfred Nobel 1975}, 1975.
\href{https://www.nobelprize.org/prizes/economic-sciences/1975/press-release/1000/}
{Nobel Prize report}.

\bibitem{KantorovichMacTutor}
MacTutor History of Mathematics Archive,
\emph{Leonid Vital'evich Kantorovich}, University of St Andrews.
\href{https://mathshistory.st-andrews.ac.uk/Biographies/Kantorovich/}
{Biographical profile} (accessed August~30, 2026).

\bibitem{Brenier1991}
Y.~Brenier,
\emph{Polar factorization and monotone rearrangement of vector-valued
functions},
Comm. Pure Appl. Math. \textbf{44} (1991), no.~4, 375--417.
\href{https://doi.org/10.1002/cpa.3160440402}
{doi:10.1002/cpa.3160440402}.

\bibitem{BahdanauEtAl2015}
D.~Bahdanau, K.~Cho, and Y.~Bengio,
\emph{Neural machine translation by jointly learning to align and translate},
International Conference on Learning Representations, 2015.
\href{https://arxiv.org/abs/1409.0473}{arXiv:1409.0473}.

\bibitem{VaswaniEtAl2017}
A.~Vaswani, N.~Shazeer, N.~Parmar, J.~Uszkoreit, L.~Jones, A.~N.~Gomez,
\L.~Kaiser, and I.~Polosukhin,
\emph{Attention is all you need},
Advances in Neural Information Processing Systems \textbf{30} (2017),
5998--6008.
\href{https://research.google/pubs/attention-is-all-you-need/}
{Google Research publication page}.

\bibitem{UszkoreitGoogle2017}
J.~Uszkoreit,
\emph{Transformer: A novel neural network architecture for language
understanding},
Google Research Blog, August~31, 2017.
\href{https://research.google/blog/transformer-a-novel-neural-network-architecture-for-language-understanding/}
{Google Research report}.

\bibitem{HeEtAl2016ResNet}
K.~He, X.~Zhang, S.~Ren, and J.~Sun,
\emph{Deep residual learning for image recognition},
Proceedings of the IEEE Conference on Computer Vision and Pattern Recognition,
2016, pp.~770--778.
\href{https://doi.org/10.1109/CVPR.2016.90}
{doi:10.1109/CVPR.2016.90}.

\bibitem{E2017DynamicalSystems}
W.~E,
\emph{A proposal on machine learning via dynamical systems},
Commun. Math. Stat. \textbf{5} (2017), no.~1, 1--11.
\href{https://doi.org/10.1007/s40304-017-0103-z}
{doi:10.1007/s40304-017-0103-z}.

\bibitem{ChenEtAl2018NeuralODE}
R.~T.~Q.~Chen, Y.~Rubanova, J.~Bettencourt, and D.~Duvenaud,
\emph{Neural ordinary differential equations},
Advances in Neural Information Processing Systems \textbf{31} (2018).
\href{https://arxiv.org/abs/1806.07366}{arXiv:1806.07366}.

\bibitem{DuvenaudResearchProfile}
D.~Duvenaud,
\emph{Research overview}, University of Toronto.
\href{https://www.cs.toronto.edu/~duvenaud/}{Author's research page}
(accessed August~30, 2026).

\bibitem{DeGroot1974}
M.~H.~DeGroot,
\emph{Reaching a consensus},
J. Amer. Statist. Assoc. \textbf{69} (1974), no.~345, 118--121.
\href{https://doi.org/10.1080/01621459.1974.10480137}
{doi:10.1080/01621459.1974.10480137}.

\bibitem{CuckerSmale2007}
F.~Cucker and S.~Smale,
\emph{Emergent behavior in flocks},
IEEE Trans. Automat. Control \textbf{52} (2007), no.~5, 852--862.
\href{https://doi.org/10.1109/TAC.2007.895842}
{doi:10.1109/TAC.2007.895842}.

\bibitem{RigolletResearchProfile}
P.~Rigollet,
\emph{Research overview}, Massachusetts Institute of Technology.
\href{https://math.mit.edu/~rigollet/index.html}{Author's research page}
(accessed August~30, 2026).

\bibitem{RigolletMITProfile2018}
L.~Hardesty,
\emph{From undisciplined to interdisciplinary}, MIT News, January~21, 2018.
\href{https://news.mit.edu/2018/faculty-profile-philippe-rigollet-0122}
{MIT profile}.

\bibitem{Dobrushin1979}
R.~L.~Dobrushin,
\emph{Vlasov equations},
Funct. Anal. Appl. \textbf{13} (1979), no.~2, 115--123.
\href{https://doi.org/10.1007/BF01077243}
{doi:10.1007/BF01077243}.

\bibitem{VillaniFieldsProfile2010}
International Mathematical Union,
\emph{Fields Medal: C\'edric Villani}, 2010.
\href{https://www.mathunion.org/fileadmin/IMU/ICM2010/offline/www.icm2010.in/prize-winners-2010/fields-medal-cedric-villani.html}
{Official prize profile}.

\bibitem{JHUFeiLuCAREER2023}
Johns Hopkins University Department of Mathematics,
\emph{Fei Lu awarded an NSF CAREER grant}, February~3, 2023.
\href{https://mathematics.jhu.edu/2023/02/03/fei-lu-awarded-an-nsf-career-grant/}
{Departmental report}.

\bibitem{BrownEtAl2020GPT3}
T.~B.~Brown et al.,
\emph{Language models are few-shot learners},
Advances in Neural Information Processing Systems \textbf{33} (2020),
1877--1901.
\href{https://proceedings.neurips.cc/paper_files/paper/2020/hash/1457c0d6bfcb4967418bfb8ac142f64a-Abstract.html}
{NeurIPS proceedings}.

\bibitem{OpenAIGPT3Report2020}
OpenAI,
\emph{Language models are few-shot learners}, May~28, 2020.
\href{https://openai.com/index/language-models-are-few-shot-learners/}
{Research report}.

\bibitem{LuEtAl2021DeepONet}
L.~Lu, P.~Jin, G.~Pang, Z.~Zhang, and G.~E.~Karniadakis,
\emph{Learning nonlinear operators via DeepONet based on the universal
approximation theorem of operators},
Nature Mach. Intell. \textbf{3} (2021), 218--229.
\href{https://doi.org/10.1038/s42256-021-00302-5}
{doi:10.1038/s42256-021-00302-5}.

\bibitem{LiEtAl2021FNO}
Z.~Li et al.,
\emph{Fourier neural operator for parametric partial differential equations},
International Conference on Learning Representations, 2021.
\href{https://arxiv.org/abs/2010.08895}{arXiv:2010.08895}.

\bibitem{SongEtAl2021ScoreSDE}
Y.~Song, J.~Sohl-Dickstein, D.~P.~Kingma, A.~Kumar, S.~Ermon, and B.~Poole,
\emph{Score-based generative modeling through stochastic differential
equations},
International Conference on Learning Representations, 2021.
\href{https://arxiv.org/abs/2011.13456}{arXiv:2011.13456}.

\bibitem{SongResearchProfile}
Y.~Song,
\emph{Research overview}.
\href{https://yang-song.net/}{Author's research page}
(accessed August~30, 2026).

\bibitem{GronwallMacTutor}
L.~Maligranda,
\emph{Thomas Hakon Gronwall}, MacTutor History of Mathematics Archive,
University of St Andrews.
\href{https://mathshistory.st-andrews.ac.uk/Biographies/Gronwall/}
{Biographical profile} (accessed August~30, 2026).

\bibitem{ItoKyotoPrize1998}
Inamori Foundation,
\emph{Kiyosi It\^o}, Kyoto Prize laureate profile, 1998.
\href{https://www.kyotoprize.org/en/laureates/kiyosi_ito/}
{Official laureate profile}.

\bibitem{DoobNASMemoir2016}
D.~W.~Stroock,
\emph{Joseph L. Doob, 1910--2004},
Biographical Memoirs, National Academy of Sciences, 2016.
\href{https://www.nasonline.org/wp-content/uploads/2024/06/doob-joseph.pdf}
{National Academy memoir}.

\end{thebibliography}

\end{document}
